%% The evidence for the deconstruction, as a document of its own: latexmk -pdf evidence.tex
%% Cross-references to the paper's sections are resolved from main.aux, so build main.tex first.
\documentclass[11pt]{article}
\usepackage[T1]{fontenc}
\usepackage[utf8]{inputenc}
\usepackage{lmodern}
\usepackage[margin=1in]{geometry}
\usepackage{microtype}
\usepackage{amsmath,amssymb}
\usepackage{booktabs}
\usepackage{tabularx}
\usepackage{longtable}
\usepackage{array}
\usepackage{enumitem}
\usepackage{xcolor}
\usepackage{xr-hyper}
\usepackage[colorlinks=true,linkcolor=black,citecolor=blue!50!black,urlcolor=blue!50!black]{hyperref}
\externaldocument{main}
\newcommand{\prim}[1]{\textmd{\textsc{#1}}}
\setlist{nosep,leftmargin=*}

\title{Evidence for the Deconstruction\\[4pt]\large Supplement to \emph{Agent = Harness(Model): Five Mechanistic Primitives, and What Is Not One}}
\author{Chase Dovey\\Average Joes Lab\\\texttt{https://cdovey.dev}}
\date{}

\begin{document}
\maketitle
\noindent This document is Appendix~C of the paper: every assignment in its deconstruction (\S\ref{sec:deconstruction}), with its evidence. The first six harnesses are deconstructed in full, the thirteen of the wider set at summary level, two to five components per primitive. For harnesses with published source, a citation is a path and a line range at the pinned commit; for a harness read from documentation it is a page and section heading on the date read. The same files are published in the repository under \texttt{paper/evidence/}, with the brief the deconstruction followed. Each harness's entry opens with its source, its pinned version and the path conventions used; \emph{not found} means the mechanism was looked for and not found, not that it is absent.

\appendix
\setcounter{section}{2}
\section{Evidence for the Deconstruction}
\begingroup\footnotesize\sloppy
\section{Evidence for the Deconstruction}
\label{app:evidence}
Every assignment in \S\ref{sec:deconstruction}, with its evidence. For harnesses with published source, a citation is a path and a line range at the pinned commit; for Claude Code it is a documentation page and section heading on the date read. The same files are published in the repository under \texttt{paper/evidence/}. \textbf{Path conventions} are given at the start of each harness; \textbf{not found} means the mechanism was looked for and not found, not that it is absent.
\small
\subsection{Claude Code}

\textbf{Source:} https://code.claude.com/docs/en/overview (official docs; source not published) $\cdot$ \textbf{Version:} documentation as of 2026-10-15; the docs state no single CLI version (pages reference behaviours by version, the highest seen being v2.1.299; agent-loop page bundles "Claude Code v2.1.290 or later" in SDK v0.3.290). Engineering posts: claude.dev/blog "Prompt caching is everything" (2026-04-30), "A harness for every task" (2026-06-02), "Getting started with loops" (2026-06-30) $\cdot$ \textbf{Language:} not documented (docs describe a native binary and an npm package; troubleshooting refers to a "JavaScript heap")
\textbf{Path conventions:} \texttt{D} = https://code.claude.com/docs/en ; \texttt{SDK} = \texttt{D/\allowbreak{}agent-\allowbreak{}sdk} ; \texttt{B} = https://claude.dev/blog ; every citation is \texttt{page § Section heading} and was read on 2026-10-15. Where a page did not say how something works I write "not documented".
\textbf{What it is:} Claude Code is Anthropic's agentic coding harness: a terminal/IDE/desktop/web/CI program that wraps Claude models with file, search, shell and web tools, permission gating, memory loading and the loop that chains actions (\texttt{D/\allowbreak{}glossary § Agentic harness}). It reaches the model over the Anthropic Messages API (\texttt{/\allowbreak{}v1/\allowbreak{}messages}, SSE streaming) or via Amazon Bedrock, Google Cloud's Agent Platform, Microsoft Foundry or an LLM gateway (\texttt{D/\allowbreak{}llm-\allowbreak{}gateway-\allowbreak{}protocol § API formats}). The same loop is exposed as a library (Agent SDK, Python/TypeScript) and as \texttt{claude -\allowbreak{}p} for scripts (\texttt{SDK/\allowbreak{}overview § Compare the Agent SDK to other Claude tools}).

\subsubsection*{Control flow}
\begin{itemize}
\item \textbf{Core loop:} "Claude continues calling tools and processing results until it produces a response with no tool calls"; each cycle (model output with tool calls $\rightarrow$ tools run $\rightarrow$ results fed back) is one turn, with no control returned to the caller between turns. \texttt{SDK/\allowbreak{}agent-\allowbreak{}loop § The loop at a glance}, \texttt{§ Turns and messages}
\item \textbf{Stop condition / result:} the loop ends on a text-only response; the \texttt{ResultMessage.\allowbreak{}subtype} is \texttt{success}, \texttt{error\_\allowbreak{}max\_\allowbreak{}turns}, \texttt{error\_\allowbreak{}max\_\allowbreak{}budget\_\allowbreak{}usd}, \texttt{error\_\allowbreak{}during\_\allowbreak{}execution} or \texttt{error\_\allowbreak{}max\_\allowbreak{}structured\_\allowbreak{}output\_\allowbreak{}retries}; \texttt{stop\_\allowbreak{}reason} carries \texttt{end\_\allowbreak{}turn}, \texttt{max\_\allowbreak{}tokens} or \texttt{refusal}. \texttt{SDK/\allowbreak{}agent-\allowbreak{}loop § Handle the result}
\item \textbf{Step budget:} \texttt{max\_\allowbreak{}turns}/\texttt{-\allowbreak{}-\allowbreak{}max-\allowbreak{}turns} "counts tool-use turns only"; no limit by default (print mode only). A message still queued when the limit ends a turn starts a new turn with a fresh count. \texttt{SDK/\allowbreak{}agent-\allowbreak{}loop § Turns and budget}; \texttt{D/\allowbreak{}cli-\allowbreak{}reference § CLI flags} (\texttt{-\allowbreak{}-\allowbreak{}max-\allowbreak{}turns})
\item \textbf{Token/cost budget:} \texttt{max\_\allowbreak{}budget\_\allowbreak{}usd} is compared "after model responses arrive", so spend can pass the cap by one response; subagent spend counts; at the cap, spawning fails with \texttt{Budget limit reached} and background subagents are stopped. \texttt{SDK/\allowbreak{}agent-\allowbreak{}loop § Budget headroom}
\item \textbf{Permission gate order (harness consumes the answer):} hooks $\rightarrow$ deny rules $\rightarrow$ ask rules $\rightarrow$ permission mode $\rightarrow$ allow rules $\rightarrow$ \texttt{canUseTool}/prompt; a hook \texttt{allow} never skips deny/ask; \texttt{dontAsk} denies at the last step. \texttt{SDK/\allowbreak{}permissions § How permissions are evaluated}; rule order "deny, then ask, then allow \ldots{} first match" \texttt{D/\allowbreak{}permissions § Manage permissions}
\item \textbf{Permission modes:} \texttt{default} (Manual), \texttt{acceptEdits}, \texttt{plan}, \texttt{auto}, \texttt{dontAsk}, \texttt{bypassPermissions}; auto is the built-in starting mode for interactive terminal/VS Code on v2.1.283+. \texttt{D/\allowbreak{}permission-\allowbreak{}modes § Available modes}, \texttt{§ Which mode a session starts in}
\item \textbf{Auto-mode classifier gate:} a second model decides most tool calls; "runs on Claude Sonnet 5 by default rather than on your /model selection"; reads/working-dir edits skip it; it sees user messages, non-read-only tool calls and CLAUDE.md but "Tool results are stripped"; falls back to prompting after 3 blocks in a row or 20 total. \texttt{D/\allowbreak{}permission-\allowbreak{}modes § How auto mode evaluates actions}, \texttt{§ When auto mode falls back}
\item \textbf{Plan mode:} edits stay blocked until a plan is approved; the approval prompt (\texttt{Yes, and use auto mode} / \texttt{Yes, manually approve edits} / \texttt{No, keep planning}) switches the session's permission mode. \texttt{D/\allowbreak{}permission-\allowbreak{}modes § Analyze before you edit with plan mode}, \texttt{§ Review and approve a plan}. Entering/leaving is itself a tool call (\texttt{EnterPlanMode}/\texttt{ExitPlanMode}) so the tool set never changes. \texttt{B/\allowbreak{}lessons-\allowbreak{}from-\allowbreak{}building-\allowbreak{}claude-\allowbreak{}code-\allowbreak{}prompt-\allowbreak{}caching-\allowbreak{}is-\allowbreak{}everything § Using Plan Mode to design around the cache}
\item \textbf{Stop hook as loop extender:} a \texttt{Stop} hook returning \texttt{decision:\allowbreak{} "block"} keeps Claude working; capped at 8 consecutive continuations (reset on tool use), \texttt{CLAUDE\_\allowbreak{}CODE\_\allowbreak{}STOP\_\allowbreak{}HOOK\_\allowbreak{}BLOCK\_\allowbreak{}CAP}. \texttt{D/\allowbreak{}hooks § Stop}; \texttt{D/\allowbreak{}env-\allowbreak{}vars § Variables} (\texttt{CLAUDE\_\allowbreak{}CODE\_\allowbreak{}STOP\_\allowbreak{}HOOK\_\allowbreak{}BLOCK\_\allowbreak{}CAP})
\item \textbf{Goal loop:} \texttt{/\allowbreak{}goal} "is a wrapper around a session-scoped prompt-based Stop hook"; after each turn the condition plus conversation go to the small fast model, which returns Not yet met / Met / Impossible; the loop halts if no tool use for several turns. \texttt{D/\allowbreak{}goal § How evaluation works}
\item \textbf{PostToolBatch gate:} fires once after all tool calls in a batch resolve, "before the next model request"; \texttt{decision:\allowbreak{} "block"} or \texttt{continue:\allowbreak{} false} stops the loop before the next model call. \texttt{D/\allowbreak{}hooks § PostToolBatch}
\item \textbf{Sub-agent delegation (Agent tool):} the model spawns a subagent in its own context; depth default 3 layers (\texttt{CLAUDE\_\allowbreak{}CODE\_\allowbreak{}MAX\_\allowbreak{}SUBAGENT\_\allowbreak{}SPAWN\_\allowbreak{}DEPTH}, \texttt{1} disables nesting); at most 20 concurrent (\texttt{CLAUDE\_\allowbreak{}CODE\_\allowbreak{}MAX\_\allowbreak{}CONCURRENT\_\allowbreak{}SUBAGENTS}); only the final message returns to the parent. \texttt{D/\allowbreak{}sub-\allowbreak{}agents § Let subagents spawn their own subagents}, \texttt{§ Concurrent subagent limit}; \texttt{SDK/\allowbreak{}subagents § What subagents inherit}
\item \textbf{Foreground vs background sequencing:} order of decision: teammates $\rightarrow$ \texttt{CLAUDE\_\allowbreak{}CODE\_\allowbreak{}DISABLE\_\allowbreak{}BACKGROUND\_\allowbreak{}TASKS=1} $\rightarrow$ fork mode on (background, Claude can't request foreground) $\rightarrow$ fork mode off (background by default, foreground when Claude needs the result); background results arrive later as a completion notification. \texttt{D/\allowbreak{}sub-\allowbreak{}agents § Run subagents in foreground or background}
\item \textbf{Model keyed to work:} \texttt{opusplan} "Uses opus in plan mode, then switches to sonnet for execution"; subagent model order: per-invocation \texttt{model} $\rightarrow$ frontmatter \texttt{model} $\rightarrow$ \texttt{CLAUDE\_\allowbreak{}CODE\_\allowbreak{}SUBAGENT\_\allowbreak{}MODEL} $\rightarrow$ main model; background functionality (titles, summaries, goal evaluator) uses the Haiku-class model. \texttt{D/\allowbreak{}model-\allowbreak{}config § Model aliases}, \texttt{§ Environment variables}; \texttt{D/\allowbreak{}sub-\allowbreak{}agents § Choose a model}; \texttt{D/\allowbreak{}sessions § Name your sessions}
\item \textbf{Dynamic workflows (model-written program sequencing model calls):} "A dynamic workflow is a JavaScript script that orchestrates many subagents"; \texttt{agent(\allowbreak{})}, \texttt{parallel(\allowbreak{})}, \texttt{pipeline(\allowbreak{})}; up to 16 concurrent agents, 4,096 items per call, 1,000 agents per run; no mid-run user input; \texttt{Date.\allowbreak{}now(\allowbreak{})}/\texttt{Math.\allowbreak{}random(\allowbreak{})} throw so a relaunch replays deterministically. \texttt{D/\allowbreak{}workflows § What the saved script looks like}, \texttt{§ Behavior and limits}, \texttt{§ Edit a saved script}
\item \textbf{Workflow resume:} completed agents return saved results; the first agent whose prompt differs, and every agent after it, runs again. \texttt{D/\allowbreak{}workflows § Which agents run again}
\item \textbf{Agent teams (outer harness of peer sessions):} lead + teammates with a shared task list and a per-agent mailbox JSON file at \texttt{\textasciitilde{}/\allowbreak{}.\allowbreak{}claude/\allowbreak{}teams/\allowbreak{}\{team-\allowbreak{}name\}/\allowbreak{}inboxes/\allowbreak{}\{agent-\allowbreak{}name\}.\allowbreak{}json}; experimental, \texttt{CLAUDE\_\allowbreak{}CODE\_\allowbreak{}EXPERIMENTAL\_\allowbreak{}AGENT\_\allowbreak{}TEAMS=1}; no nested teams. \texttt{D/\allowbreak{}agent-\allowbreak{}teams § Architecture}, \texttt{§ Limitations}
\item \textbf{Prompt- and agent-based hooks (harness-sequenced model calls):} \texttt{type:\allowbreak{} "prompt"} sends hook input to a model (default: the background-functionality model, 30 s timeout) to allow/block; \texttt{type:\allowbreak{} "agent"} spawns a subagent with tools (60 s). \texttt{D/\allowbreak{}hooks § Prompt-\allowbreak{}based hooks}, \texttt{§ Agent-\allowbreak{}based hooks}
\item \textbf{Scheduled turns:} \texttt{/\allowbreak{}loop} creates cron tasks; "The scheduler checks every second for due tasks and enqueues them at low priority"; fires between turns, never mid-response; 7-day expiry; self-paced loops end via \texttt{ScheduleWakeup} with \texttt{stop:\allowbreak{} true}. \texttt{D/\allowbreak{}scheduled-\allowbreak{}tasks § How scheduled tasks run}, \texttt{§ Stop a loop}
\item \textbf{Non-interactive exit wait:} after the turn ends, a \texttt{claude -\allowbreak{}p} run stays open for background work; "stops whatever is still running after 10 minutes of continuous idle waiting" (\texttt{CLAUDE\_\allowbreak{}CODE\_\allowbreak{}PRINT\_\allowbreak{}BG\_\allowbreak{}WAIT\_\allowbreak{}CEILING\_\allowbreak{}MS}). \texttt{D/\allowbreak{}headless § Background tasks at exit}
\end{itemize}

\subsubsection*{Input}
\begin{itemize}
\item \textbf{Prompt capture:} interactive terminal prompt, \texttt{claude "query"}, \texttt{claude -\allowbreak{}p "query"}, and piped stdin ("Piped stdin is capped at 10MB"). \texttt{D/\allowbreak{}cli-\allowbreak{}reference § CLI commands}; \texttt{D/\allowbreak{}headless § Pipe data through Claude}
\item \textbf{Mid-turn message queue:} a message typed while Claude works is queued; "Claude Code passes it to Claude as soon as those tool calls finish, within the same turn"; commands are held until the turn ends. \texttt{D/\allowbreak{}interactive-\allowbreak{}mode § When Claude Code sends what you queued}
\item \textbf{Interrupt:} \texttt{Esc} cancels the running tool call and Claude waits; queued messages are sent next. \texttt{D/\allowbreak{}how-\allowbreak{}claude-\allowbreak{}code-\allowbreak{}works § Interrupt and steer}; \texttt{D/\allowbreak{}interactive-\allowbreak{}mode § General controls}
\item \textbf{Model-consumed answers:} \texttt{AskUserQuestion} returns the user's choice to the model; questions stay open until answered unless \texttt{askUserQuestionTimeout} (\texttt{60s}/\texttt{5m}/\texttt{10m}) is set. \texttt{D/\allowbreak{}tools-\allowbreak{}reference § AskUserQuestion tool behavior}
\item \textbf{Tool results as input:} each tool execution yields a \texttt{UserMessage} carrying the tool result back to the model. \texttt{SDK/\allowbreak{}agent-\allowbreak{}loop § Message types}
\item \textbf{Event/stream inputs:} \texttt{-\allowbreak{}-\allowbreak{}input-\allowbreak{}format stream-\allowbreak{}json} accepts user messages over stdin; \texttt{Monitor} feeds "each output line back to Claude" and can treat WebSocket messages as events; channels push external events into the session. \texttt{D/\allowbreak{}cli-\allowbreak{}reference § CLI flags}; \texttt{D/\allowbreak{}tools-\allowbreak{}reference § Tools reference} (\texttt{Monitor} row); \texttt{D/\allowbreak{}glossary § Channel}
\item \textbf{Non-human turns captured through the same entry:} \texttt{UserPromptSubmit} "also fires for scheduled tasks, background subagent reports, and messages from other sessions". \texttt{D/\allowbreak{}hooks § UserPromptSubmit}
\item \textbf{Rich input:} \texttt{Ctrl+V} pastes an image as an \texttt{[Image \#N]} chip; \texttt{@path} mentions expand file content; \texttt{@server:\allowbreak{}protocol:\allowbreak{}/\allowbreak{}/\allowbreak{}resource} fetches MCP resources as attachments. \texttt{D/\allowbreak{}interactive-\allowbreak{}mode § General controls}; \texttt{SDK/\allowbreak{}modifying-\allowbreak{}system-\allowbreak{}prompts § Turn off the context your agent replaces} (\texttt{CLAUDE\_\allowbreak{}CODE\_\allowbreak{}DISABLE\_\allowbreak{}FILE\_\allowbreak{}MENTIONS}); \texttt{D/\allowbreak{}mcp § Use MCP resources}
\item \textbf{MCP elicitation:} a server's mid-task form/URL request is shown as a dialog and the answer goes back to the server (not the model); an \texttt{Elicitation} hook can answer it. \texttt{D/\allowbreak{}mcp § Respond to MCP elicitation requests}
\end{itemize}

\subsubsection*{Context}
\begin{itemize}
\item \textbf{Request layering (what one call contains):} system prompt (core instructions + tool definitions) $\rightarrow$ project context (CLAUDE.md, auto memory, unscoped rules) $\rightarrow$ conversation (messages, tool results); ordered most- to least-stable for prefix caching. \texttt{D/\allowbreak{}prompt-\allowbreak{}caching § How the cache is organized}
\item \textbf{Two system prompts:} full vs "a shorter one with abbreviated tool descriptions"; newer models get the short one unless \texttt{CLAUDE\_\allowbreak{}CODE\_\allowbreak{}SIMPLE\_\allowbreak{}SYSTEM\_\allowbreak{}PROMPT=0}. \texttt{D/\allowbreak{}env-\allowbreak{}vars § Variables} (\texttt{CLAUDE\_\allowbreak{}CODE\_\allowbreak{}SIMPLE\_\allowbreak{}SYSTEM\_\allowbreak{}PROMPT})
\item \textbf{System prompt attribution block:} "a short attribution block to the system prompt containing the client version and a fingerprint"; stripped by \texttt{api.\allowbreak{}anthropic.\allowbreak{}com}; \texttt{CLAUDE\_\allowbreak{}CODE\_\allowbreak{}ATTRIBUTION\_\allowbreak{}HEADER=0} omits it. \texttt{D/\allowbreak{}llm-\allowbreak{}gateway-\allowbreak{}protocol § System prompt attribution block}
\item \textbf{CLAUDE.md delivery:} loaded "as a user message after the system prompt"; all discovered files are concatenated root$\rightarrow$working directory, \texttt{CLAUDE.\allowbreak{}local.\allowbreak{}md} after \texttt{CLAUDE.\allowbreak{}md}; subdirectory files load "once Claude reads, writes, or edits another file in that subdirectory". \texttt{D/\allowbreak{}glossary § CLAUDE.\allowbreak{}md}; \texttt{D/\allowbreak{}memory § How CLAUDE.\allowbreak{}md files load}
\item \textbf{Imports and rules:} \texttt{@path} imports expand at launch; \texttt{.\allowbreak{}claude/\allowbreak{}rules/\allowbreak{}*.\allowbreak{}md} load every session, or on demand when \texttt{paths:\allowbreak{}} frontmatter matches a file Claude touches; a CLAUDE.md over 4 MiB is skipped. \texttt{D/\allowbreak{}memory § Import additional files}, \texttt{§ Path-\allowbreak{}specific rules}, \texttt{§ How it works}
\item \textbf{Auto memory (model-kept record):} \texttt{\textasciitilde{}/\allowbreak{}.\allowbreak{}claude/\allowbreak{}projects/\allowbreak{}<project>/\allowbreak{}memory/\allowbreak{}MEMORY.\allowbreak{}md}; "The first 200 lines of MEMORY.md, or the first 25KB, whichever comes first, are loaded"; topic files read on demand; a write past the limit returns an error telling Claude to rewrite the index. \texttt{D/\allowbreak{}memory § How it works}, \texttt{§ Storage location}
\item \textbf{System reminders:} harness-added context delivered "wrapped in \texttt{<system-\allowbreak{}reminder>} tags inside a user message, or \ldots{} a separate message with the \texttt{system} role": CLAUDE.md, output style, hook \texttt{additionalContext}, skill list, file-changed notes, attribution lines. \texttt{D/\allowbreak{}glossary § System reminder}
\item \textbf{Output style injection:} "Claude Code sends the active style's instructions with every request" as a conversation message, so the cached prefix survives a switch. \texttt{D/\allowbreak{}output-\allowbreak{}styles § How output styles work}; \texttt{D/\allowbreak{}prompt-\allowbreak{}caching § Changing output style}
\item \textbf{Skills in context:} descriptions (truncated at 1,536 chars) are always in context; the body loads on invocation as a user message and stays for the session; \texttt{disable-\allowbreak{}model-\allowbreak{}invocation:\allowbreak{} true} keeps the description out. \texttt{D/\allowbreak{}skills § Frontmatter fields}, \texttt{§ Controlling visibility and context cost}; \texttt{D/\allowbreak{}prompt-\allowbreak{}caching § Invoking skills and commands}
\item \textbf{Skill dynamic context:} \texttt{} !\texttt{command} \texttt{} runs a shell command before the model sees the skill and splices the output in. \texttt{D/\allowbreak{}skills § Dynamic context injection}
\item \textbf{MCP tool search (deferred definitions):} only names and server instructions load; schemas arrive via the \texttt{ToolSearch} tool; \texttt{ENABLE\_\allowbreak{}TOOL\_\allowbreak{}SEARCH} \texttt{auto} loads upfront while definitions are under 10\% of context; description caps 4,096 (upfront) / 16,384 (deferred) chars. \texttt{D/\allowbreak{}mcp § Scale with MCP tool search}, \texttt{§ Configure tool search}
\item \textbf{Environment and git snapshot:} working directory, platform, shell, OS and a git branch/status/recent-commits block are added at session start; \texttt{includeGitInstructions:\allowbreak{} false} drops the git snapshot. \texttt{D/\allowbreak{}context-\allowbreak{}window} (Environment info event); \texttt{SDK/\allowbreak{}modifying-\allowbreak{}system-\allowbreak{}prompts § Turn off the context your agent replaces}; \texttt{D/\allowbreak{}prompt-\allowbreak{}caching § Cache scope}
\item \textbf{Fitting tool results:} Bash output inline "up to roughly 30,000 characters"; past that a file path plus 2,000-char preview; MCP results capped at 25,000 tokens and text over 50,000 chars spilled to a file; \texttt{Read} returns at most 25,000 tokens per call and refuses whole files over 256 KB. \texttt{D/\allowbreak{}tools-\allowbreak{}reference § Output limits}, \texttt{§ Large files}; \texttt{D/\allowbreak{}mcp § Raise the limit for a specific tool}; \texttt{D/\allowbreak{}env-\allowbreak{}vars § Variables} (\texttt{MAX\_\allowbreak{}MCP\_\allowbreak{}OUTPUT\_\allowbreak{}TOKENS})
\item \textbf{Auto-compaction trigger:} by default "compacts when the conversation reaches the model's context limit"; native-1M models compact "at about 967K tokens"; window settable 100K--1M via \texttt{/\allowbreak{}autocompact}, \texttt{autoCompactWindow}, \texttt{-\allowbreak{}-\allowbreak{}autocompact}, \texttt{CLAUDE\_\allowbreak{}CODE\_\allowbreak{}AUTO\_\allowbreak{}COMPACT\_\allowbreak{}WINDOW}. \texttt{D/\allowbreak{}model-\allowbreak{}config § Default auto-\allowbreak{}compact thresholds}, \texttt{§ Set the auto-\allowbreak{}compact window}
\item \textbf{Compaction order:} "It clears older tool outputs first, then summarizes the conversation if needed." \texttt{D/\allowbreak{}how-\allowbreak{}claude-\allowbreak{}code-\allowbreak{}works § When context fills up}
\item \textbf{Compaction as a forked request:} the summarization call reuses the same system prompt, tools and history "plus a summarization instruction appended as a final user message"; inherits the session's thinking config. \texttt{D/\allowbreak{}prompt-\allowbreak{}caching § Compacting the conversation}; \texttt{D/\allowbreak{}context-\allowbreak{}window § What survives compaction}; \texttt{B/\allowbreak{}lessons-\allowbreak{}from-\allowbreak{}building-\allowbreak{}claude-\allowbreak{}code-\allowbreak{}prompt-\allowbreak{}caching-\allowbreak{}is-\allowbreak{}everything § COMPACTING WITHOUT BREAKING THE CACHE}
\item \textbf{What survives compaction:} system prompt/output style still apply; project-root CLAUDE.md, auto memory, plan file re-injected from disk; up to five recently modified files re-read (over 5,000 tokens $\rightarrow$ path reference); invoked skill bodies re-injected "capped at 5,000 tokens per skill and 25,000 tokens total"; hook context is summarized away. \texttt{D/\allowbreak{}context-\allowbreak{}window § What survives compaction}
\item \textbf{Hook-supplied context:} \texttt{additionalContext} is wrapped in a system reminder; \texttt{additionalContext}, \texttt{systemMessage}, \texttt{initialUserMessage} capped at 10,000 chars (longer saved to file with a 2,000-char preview); \texttt{PostToolUse.\allowbreak{}updatedToolOutput} replaces what Claude sees. \texttt{D/\allowbreak{}hooks § JSON output}, \texttt{§ Add context for Claude}, \texttt{§ PostToolUse}
\item \textbf{Sub-agent context:} a non-fork subagent gets its own system prompt plus environment details, the task message, CLAUDE.md hierarchy (not for Explore/Plan), git snapshot, preloaded skills; not the parent history, output style or auto memory. A fork gets the full conversation, system prompt and tools. \texttt{D/\allowbreak{}sub-\allowbreak{}agents § What loads at startup}, \texttt{§ How forks differ from other subagents}
\item \textbf{Sub-agent output scanning:} before the parent reads a report, control-tag imitations get a backslash and a \texttt{[harness:\allowbreak{} .\allowbreak{}.\allowbreak{}.\allowbreak{}]} marker line is prepended; nothing is removed. \texttt{D/\allowbreak{}sub-\allowbreak{}agents § Subagent output scanning}; \texttt{SDK/\allowbreak{}subagents § What subagents inherit}
\end{itemize}

\subsubsection*{Model interface}
\begin{itemize}
\item \textbf{Wire format:} Anthropic Messages format, "Inference requests post to \texttt{/\allowbreak{}v1/\allowbreak{}messages?beta=true}"; optional \texttt{/\allowbreak{}v1/\allowbreak{}messages/\allowbreak{}count\_\allowbreak{}tokens} (falls back to a character estimate); headers \texttt{anthropic-\allowbreak{}version:\allowbreak{} 2023-\allowbreak{}06-\allowbreak{}01}, \texttt{anthropic-\allowbreak{}beta}, \texttt{x-\allowbreak{}claude-\allowbreak{}code-\allowbreak{}session-\allowbreak{}id}, \texttt{x-\allowbreak{}claude-\allowbreak{}code-\allowbreak{}agent-\allowbreak{}id}. \texttt{D/\allowbreak{}llm-\allowbreak{}gateway-\allowbreak{}protocol § API formats}, \texttt{§ Optional endpoints and startup traffic}, \texttt{§ Request headers}
\item \textbf{Providers/backends:} \texttt{CLAUDE\_\allowbreak{}CODE\_\allowbreak{}USE\_\allowbreak{}BEDROCK}, \texttt{CLAUDE\_\allowbreak{}CODE\_\allowbreak{}USE\_\allowbreak{}VERTEX}, \texttt{CLAUDE\_\allowbreak{}CODE\_\allowbreak{}USE\_\allowbreak{}FOUNDRY}, Claude Platform on AWS, or any gateway via \texttt{ANTHROPIC\_\allowbreak{}BASE\_\allowbreak{}URL}; Bedrock uses \texttt{invoke-\allowbreak{}with-\allowbreak{}response-\allowbreak{}stream} binary eventstream. \texttt{D/\allowbreak{}env-\allowbreak{}vars § Variables}; \texttt{D/\allowbreak{}llm-\allowbreak{}gateway-\allowbreak{}protocol § API formats}
\item \textbf{Streaming receipt:} "Claude Code reads each streaming inference response event by event as it arrives"; expects the full sequence through \texttt{message\_\allowbreak{}delta}/\texttt{message\_\allowbreak{}stop}, else treats it as a dropped connection. \texttt{D/\allowbreak{}llm-\allowbreak{}gateway-\allowbreak{}protocol § Streaming}
\item \textbf{Retry policy:} "retries transient failures up to 10 times with exponential backoff" (\texttt{CLAUDE\_\allowbreak{}CODE\_\allowbreak{}MAX\_\allowbreak{}RETRIES}, capped at 15); retries server errors/529/timeouts before output starts, dropped connections (up to 2 after thinking), stalled streams (once), no-headers (once per request), temporary 429s; never re-runs after a completed text or tool block "because that could execute the same tool calls twice". \texttt{D/\allowbreak{}errors § Automatic retries}
\item \textbf{529 backoff:} \texttt{CLAUDE\_\allowbreak{}CODE\_\allowbreak{}OVERLOADED\_\allowbreak{}RETRY\_\allowbreak{}BASE\_\allowbreak{}DELAY\_\allowbreak{}MS} 500 $\rightarrow$ \texttt{CLAUDE\_\allowbreak{}CODE\_\allowbreak{}OVERLOADED\_\allowbreak{}RETRY\_\allowbreak{}MAX\_\allowbreak{}DELAY\_\allowbreak{}MS} 32000; \texttt{retry-\allowbreak{}after} honoured, and >60 s stops retries outside watchdog sessions; \texttt{x-\allowbreak{}should-\allowbreak{}retry} read. \texttt{D/\allowbreak{}errors § Automatic retries}; \texttt{D/\allowbreak{}llm-\allowbreak{}gateway-\allowbreak{}protocol § Response headers}
\item \textbf{Indefinite retry for unattended runs:} \texttt{CLAUDE\_\allowbreak{}CODE\_\allowbreak{}RETRY\_\allowbreak{}WATCHDOG=1} "Retries 429 and 529 capacity errors indefinitely", backing off up to 5 minutes. \texttt{D/\allowbreak{}env-\allowbreak{}vars § Variables}
\item \textbf{Timeouts and watchdogs:} \texttt{API\_\allowbreak{}TIMEOUT\_\allowbreak{}MS} default 600000; first-byte deadline 180 s (direct API) / 300 s (elsewhere) + 1 s per 32 KB of body; event-level watchdog 300 s; byte-level watchdog 180 s/300 s; body idle timeout 5 min. \texttt{D/\allowbreak{}env-\allowbreak{}vars § Variables} (\texttt{API\_\allowbreak{}TIMEOUT\_\allowbreak{}MS}); \texttt{D/\allowbreak{}network-\allowbreak{}config § Streaming idle watchdogs}
\item \textbf{Backup model when the request can't be served:} \texttt{-\allowbreak{}-\allowbreak{}fallback-\allowbreak{}model sonnet,haiku} / \texttt{fallbackModel} runs "when the primary model is overloaded, unavailable, or returns another non-retryable server error"; chain of $\leq$3; applies to the current turn only; subagents use it too. \texttt{D/\allowbreak{}model-\allowbreak{}config § Fallback model chains}
\item \textbf{Safety-flag model fallback:} when a server safety classifier flags a request, Claude Code "can re-run it on a fallback model" (e.g. Opus 4.8 for cybersecurity); \texttt{switchModelsOnFlag}. \texttt{D/\allowbreak{}model-\allowbreak{}config § Automatic model fallback}
\item \textbf{Capability-rejection recovery:} on a 400 rejecting \texttt{thinking}, a mid-conversation system message, \texttt{cache\_\allowbreak{}control}, the advisor tool or \texttt{output\_\allowbreak{}config.\allowbreak{}effort}, Claude Code retries without that capability for the rest of the conversation; a rejected thinking signature strips earlier thinking blocks. \texttt{D/\allowbreak{}llm-\allowbreak{}gateway-\allowbreak{}protocol § Automatic retry and error forwarding}
\item \textbf{Context-limit rejection:} a request rejected because input + \texttt{max\_\allowbreak{}tokens} exceeds the limit is retried "with a reduced max\_tokens", then compacted when no reduction fits. \texttt{D/\allowbreak{}errors § Automatic retries}
\item \textbf{Prompt-cache markers and TTL:} \texttt{cache\_\allowbreak{}control} on system blocks and messages; main conversation gets 1 h TTL on subscription usage, 5 min otherwise; subagents/compaction/titles get 5 min; \texttt{promptCacheTtl}, \texttt{FORCE\_\allowbreak{}PROMPT\_\allowbreak{}CACHING\_\allowbreak{}5M}, \texttt{DISABLE\_\allowbreak{}PROMPT\_\allowbreak{}CACHING}. \texttt{D/\allowbreak{}prompt-\allowbreak{}caching § Which TTL each request gets}, \texttt{§ Choose the TTL yourself}, \texttt{§ Disable prompt caching}
\item \textbf{Reasoning parameters on the wire:} \texttt{thinking:\allowbreak{} \{"type":\allowbreak{} "adaptive"\}} for 4.6+ models; effort and structured output in \texttt{output\_\allowbreak{}config} with paired beta headers; \texttt{MAX\_\allowbreak{}THINKING\_\allowbreak{}TOKENS} fixed budget ($\geq$1,024, $\leq$ max output - 1); \texttt{CLAUDE\_\allowbreak{}CODE\_\allowbreak{}MAX\_\allowbreak{}OUTPUT\_\allowbreak{}TOKENS} (default 32000 / cap 128000 for unknown IDs). \texttt{D/\allowbreak{}llm-\allowbreak{}gateway-\allowbreak{}protocol § Feature pass-\allowbreak{}through}; \texttt{D/\allowbreak{}env-\allowbreak{}vars § Variables}
\item \textbf{Retry visibility:} a \texttt{system/\allowbreak{}api\_\allowbreak{}retry} event (attempt, max\_retries, retry\_delay\_ms, error\_status, error category) is emitted in stream-json before each retry. \texttt{D/\allowbreak{}headless § Handle API retries}
\end{itemize}

\subsubsection*{Output}
\begin{itemize}
\item \textbf{Built-in tool set:} Agent, AskUserQuestion, Bash, Edit, EndConversation, Glob, Grep, LSP, Monitor, NotebookEdit, PowerShell, Read, SendMessage, Skill, TaskCreate/TaskUpdate/TodoWrite, ToolSearch, WebFetch, WebSearch, Workflow, Write (+ MCP tools \texttt{mcp\_\allowbreak{}\_\allowbreak{}<server>\_\allowbreak{}\_\allowbreak{}<tool>}). \texttt{D/\allowbreak{}tools-\allowbreak{}reference § Tools reference} (table); \texttt{SDK/\allowbreak{}agent-\allowbreak{}loop § Built-\allowbreak{}in tools}
\item \textbf{Concurrent execution of one response's calls:} "Read-only tools (like Read, Glob, Grep, and MCP tools marked as read-only) can run concurrently. Tools that modify state (like Edit, Write, and Bash) run sequentially"; custom tools opt in via \texttt{readOnlyHint}. \texttt{SDK/\allowbreak{}agent-\allowbreak{}loop § Parallel tool execution}
\item \textbf{Bash executor:} default timeout 2 min (\texttt{BASH\_\allowbreak{}DEFAULT\_\allowbreak{}TIMEOUT\_\allowbreak{}MS}), ceiling 10 min (\texttt{BASH\_\allowbreak{}MAX\_\allowbreak{}TIMEOUT\_\allowbreak{}MS}); a foreground command that hits its timeout "moves it to the background instead of stopping it"; background commands get 30 min (max 2 h); env vars don't persist between commands; shell aliases/options captured in \texttt{\textasciitilde{}/\allowbreak{}.\allowbreak{}claude/\allowbreak{}shell-\allowbreak{}snapshots/\allowbreak{}}. \texttt{D/\allowbreak{}tools-\allowbreak{}reference § Timeout and output limits}, \texttt{§ Foreground commands that move to the background}, \texttt{§ Time limit for background commands}, \texttt{§ What persists between commands}; \texttt{D/\allowbreak{}claude-\allowbreak{}directory § Cleaned up automatically}
\item \textbf{Sandbox technology:} macOS Seatbelt; Linux/WSL2 bubblewrap + socat; no sandbox on native Windows; covers Bash, PowerShell and Monitor commands only, not file tools, MCP servers or hooks. \texttt{D/\allowbreak{}sandboxing § OS-\allowbreak{}level enforcement}, \texttt{§ What runs outside the sandbox}
\item \textbf{Sandbox network path:} "Claude Code runs the sandbox proxy on your machine, outside the sandbox, and directs commands to it with HTTP\_PROXY, HTTPS\_PROXY, ALL\_PROXY"; the proxy checks hostnames against \texttt{allowedDomains}/\texttt{deniedDomains}, refuses names resolving only to local addresses, and can mask credentials with a sentinel (requires \texttt{tlsTerminate}). \texttt{D/\allowbreak{}sandboxing § Network isolation}, \texttt{§ Hostnames that resolve to local addresses}, \texttt{§ Mask credentials}
\item \textbf{Sandbox filesystem path:} writable: working dir + per-user \texttt{\$TMPDIR}; "Protected paths" (settings, hooks, MCP config) are denied even inside writable dirs; \texttt{excludedCommands} run unsandboxed; strict mode \texttt{allowUnsandboxedCommands:\allowbreak{} false}. \texttt{D/\allowbreak{}sandboxing § Filesystem isolation}, \texttt{§ Protected paths}, \texttt{§ Turn off the retry with strict sandbox mode}
\item \textbf{WebFetch executor:} converts HTML to Markdown, "runs the prompt against the content in a separate model call, and Claude receives the result of that call rather than the raw page"; 15-min response cache; hostname first checked against an Anthropic blocklist at \texttt{api.\allowbreak{}anthropic.\allowbreak{}com} (\texttt{skipWebFetchPreflight}). \texttt{D/\allowbreak{}tools-\allowbreak{}reference § WebFetch tool behavior}; \texttt{D/\allowbreak{}data-\allowbreak{}usage § WebFetch domain safety check}
\item \textbf{WebSearch:} server-side tool; "An interactive terminal session has a limit of 200 WebSearch calls" shared with subagents. \texttt{D/\allowbreak{}tools-\allowbreak{}reference § WebSearch tool behavior}, \texttt{§ Session search limit}
\item \textbf{Read/Edit/Write executors:} Read returns numbered lines, resizes images (>500 KB re-encoded as JPEG), reads PDFs in $\leq$20-page ranges via \texttt{pdftoppm}; Edit/Write enforce read-before-edit on older models. \texttt{D/\allowbreak{}tools-\allowbreak{}reference § Read tool behavior}, \texttt{§ Edit tool behavior}, \texttt{§ Write tool behavior}
\item \textbf{MCP executor:} stdio, HTTP (falls back to SSE), WebSocket transports; \texttt{MCP\_\allowbreak{}TIMEOUT} 30 s startup, \texttt{MCP\_\allowbreak{}TOOL\_\allowbreak{}TIMEOUT} \textasciitilde{}28 h, 60 s per HTTP request, idle timeout 5 min remote/30 min stdio, main-conversation calls auto-background after 2 min (\texttt{CLAUDE\_\allowbreak{}CODE\_\allowbreak{}MCP\_\allowbreak{}AUTO\_\allowbreak{}BACKGROUND\_\allowbreak{}MS}); invalid schemas and names >128 chars are excluded per tool. \texttt{D/\allowbreak{}mcp § Transports}, \texttt{§ Timeouts}, \texttt{§ Tools with invalid input schemas}, \texttt{§ Tools with names longer than 128 characters}
\item \textbf{Monitor executor:} Claude writes a script run in the background, or opens a WebSocket; deadline 5 min by default, max 30 min (10 min under \texttt{-\allowbreak{}p}). \texttt{D/\allowbreak{}tools-\allowbreak{}reference § Monitor tool}
\item \textbf{Rendering the stream:} an \texttt{AssistantMessage} is emitted per content block; \texttt{stream\_\allowbreak{}event} deltas (\texttt{text\_\allowbreak{}delta}, \texttt{input\_\allowbreak{}json\_\allowbreak{}delta}) with \texttt{-\allowbreak{}-\allowbreak{}include-\allowbreak{}partial-\allowbreak{}messages}; the spinner shows a live output-token count. \texttt{SDK/\allowbreak{}streaming-\allowbreak{}output § Message flow}; \texttt{D/\allowbreak{}costs § Read the token count beside the spinner}
\item \textbf{Structured output:} \texttt{-\allowbreak{}-\allowbreak{}output-\allowbreak{}format json|stream-\allowbreak{}json}, \texttt{-\allowbreak{}-\allowbreak{}json-\allowbreak{}schema} validated output in \texttt{structured\_\allowbreak{}output}; up to 5 attempts (\texttt{MAX\_\allowbreak{}STRUCTURED\_\allowbreak{}OUTPUT\_\allowbreak{}RETRIES}). \texttt{D/\allowbreak{}headless § Get structured output}; \texttt{D/\allowbreak{}env-\allowbreak{}vars § Variables}
\item \textbf{Hook executors:} handler types \texttt{command}, \texttt{http}, \texttt{mcp\_\allowbreak{}tool}, \texttt{prompt}, \texttt{agent}; "All matching hooks run in parallel. Identical handlers from multiple settings files run once"; defaults 600 s command/http/mcp\_tool; exit 2 blocks (PreToolUse) with stderr fed to Claude; \texttt{async:\allowbreak{} true} detaches; \texttt{PreToolUse.\allowbreak{}updatedInput} rewrites tool arguments before execution. \texttt{D/\allowbreak{}hooks § Configuration}, \texttt{§ Exit code output}, \texttt{§ Run hooks in the background}, \texttt{§ PreToolUse}
\item \textbf{User-facing artifacts:} file checkpoints snapshot files "before each prompt you send that starts a turn", 100 most recent kept, \texttt{/\allowbreak{}rewind} restores code and/or conversation; Bash-made and subagent edits are not tracked. \texttt{D/\allowbreak{}checkpointing § Automatic tracking}, \texttt{§ Bash command changes not tracked}, \texttt{§ Subagent edits not restored}
\item \textbf{EndConversation:} never prompts, PreToolUse hooks don't run for it, and deny rules/\texttt{-\allowbreak{}-\allowbreak{}tools} can't remove it while any other tool remains. \texttt{D/\allowbreak{}tools-\allowbreak{}reference § EndConversation tool behavior}
\item \textbf{Permission prompt surface:} Manual mode asks for Bash (except a built-in read-only set: \texttt{ls}, \texttt{cat}, \texttt{git status}\ldots{}), file modification, WebFetch (except preapproved doc domains), WebSearch; "Yes, don't ask again" persists per repository. \texttt{D/\allowbreak{}permissions § Permission system}, \texttt{§ Read-\allowbreak{}only commands}
\end{itemize}

\subsubsection*{Production concerns}
\begin{longtable}{@{}p{0.13\linewidth}p{0.17\linewidth}p{0.36\linewidth}p{0.26\linewidth}@{}}
\toprule \textbf{Concern} & \textbf{Where it lives (primitive)} & \textbf{Mechanism} & \textbf{Evidence} \\ \midrule \endhead
Persistence & Input/Context (transcript, memory); Output (checkpoints); Control flow (resume) & Every message/tool use is appended to \texttt{\textasciitilde{}/\allowbreak{}.\allowbreak{}claude/\allowbreak{}projects/\allowbreak{}<project>/\allowbreak{}<session-\allowbreak{}id>.\allowbreak{}jsonl}; swept after \texttt{cleanupPeriodDays} (30 d default); auto memory excluded from the sweep; \texttt{file-\allowbreak{}history/\allowbreak{}<session>/\allowbreak{}} snapshots, \texttt{plans/\allowbreak{}}, \texttt{tool-\allowbreak{}results/\allowbreak{}}; \texttt{-\allowbreak{}-\allowbreak{}resume}/\texttt{-\allowbreak{}-\allowbreak{}continue} restore history, model, agent, permission mode, goal and cron tasks; \texttt{-\allowbreak{}-\allowbreak{}fork-\allowbreak{}session}/\texttt{/\allowbreak{}branch} copy to a new ID; SDK \texttt{session\_\allowbreak{}store} mirrors transcripts; \texttt{CLAUDE\_\allowbreak{}CODE\_\allowbreak{}TRANSCRIPT\_\allowbreak{}LOCAL\_\allowbreak{}GC} trims history before the last compaction once the file passes 5 MB & \texttt{D/\allowbreak{}sessions § Where transcripts are stored}, \texttt{§ What a resumed session restores}; \texttt{D/\allowbreak{}claude-\allowbreak{}directory § Cleaned up automatically}; \texttt{D/\allowbreak{}memory § Storage location}; \texttt{SDK/\allowbreak{}agent-\allowbreak{}loop § Sessions and continuity}; \texttt{D/\allowbreak{}env-\allowbreak{}vars § Variables} \\
Safety \& security & Control flow (gates); Output (sandbox, executors); Context (scanning) & deny$\rightarrow$ask$\rightarrow$allow rules plus modes; auto-mode classifier with tool results stripped and a "server-side probe" over incoming tool results; critical paths (\texttt{rm -\allowbreak{}rf /\allowbreak{}}, \texttt{\textasciitilde{}}, cwd) never approvable by allow rule or hook; protected paths; OS sandbox + domain allowlist + credential masking; subagent output scanning; teammate messages can't approve prompts; WebFetch summarised by a separate model call; \texttt{-\allowbreak{}p} runs project hooks/MCP without a trust dialog (\texttt{-\allowbreak{}-\allowbreak{}bare} to skip); credentials in Keychain or a 0600 file; prompt-injection list & \texttt{D/\allowbreak{}permission-\allowbreak{}modes § What the classifier blocks by default}, \texttt{§ Critical paths}, \texttt{§ Protected paths}; \texttt{D/\allowbreak{}sandboxing § What the sandbox restricts}; \texttt{D/\allowbreak{}sub-\allowbreak{}agents § Subagent output scanning}; \texttt{D/\allowbreak{}agent-\allowbreak{}teams § Messages between agents}; \texttt{D/\allowbreak{}security § Protect against prompt injection}, \texttt{§ Additional safeguards}; \texttt{D/\allowbreak{}headless § Start faster with bare mode} \\
Observability & Model interface (api events); Output (tool events); operator records & OpenTelemetry metrics (\texttt{claude\_\allowbreak{}code.\allowbreak{}cost.\allowbreak{}usage}, \texttt{token.\allowbreak{}usage}, \texttt{session.\allowbreak{}count}, \texttt{lines\_\allowbreak{}of\_\allowbreak{}code.\allowbreak{}count}\ldots{}) and events (\texttt{user\_\allowbreak{}prompt}, \texttt{tool\_\allowbreak{}result}, \texttt{api\_\allowbreak{}request}, \texttt{api\_\allowbreak{}error}, \texttt{tool\_\allowbreak{}decision}) via \texttt{CLAUDE\_\allowbreak{}CODE\_\allowbreak{}ENABLE\_\allowbreak{}TELEMETRY=1}; content opt-ins incl. \texttt{OTEL\_\allowbreak{}LOG\_\allowbreak{}RAW\_\allowbreak{}API\_\allowbreak{}BODIES=file:\allowbreak{}<dir>}; \texttt{/\allowbreak{}usage} session cost + \texttt{Prompt cache (\allowbreak{}main)} line; status-line \texttt{current\_\allowbreak{}usage}/\texttt{prompt\_\allowbreak{}cache}; debug log \texttt{\textasciitilde{}/\allowbreak{}.\allowbreak{}claude/\allowbreak{}debug/\allowbreak{}<session-\allowbreak{}id>.\allowbreak{}txt} with hook durations; stream-json \texttt{system/\allowbreak{}init}, \texttt{api\_\allowbreak{}retry}, \texttt{compact\_\allowbreak{}boundary}, \texttt{permission\_\allowbreak{}denials}; gateway hint headers & \texttt{D/\allowbreak{}monitoring-\allowbreak{}usage § Metrics}, \texttt{§ Events}, \texttt{§ Content opt-\allowbreak{}in}; \texttt{D/\allowbreak{}costs § Prompt cache statistics}; \texttt{D/\allowbreak{}hooks § Debug hooks}; \texttt{D/\allowbreak{}headless § Read session metadata}; \texttt{D/\allowbreak{}llm-\allowbreak{}gateway-\allowbreak{}protocol § Gateway hint headers} \\
Resilience & Model interface (retries/watchdogs/fallback); Context (compaction); Control flow (resume/stall) & 10 retries w/ exponential backoff, no re-run after a completed block (partial output kept + "may be incomplete" notice, \texttt{-\allowbreak{}p} auto-continues $\leq$3$\times$); four streaming watchdogs; fallback model chain; context-limit retry with reduced \texttt{max\_\allowbreak{}tokens} then compaction; auto-compact "stops auto-compacting after a few attempts" on thrashing; interrupted tool call marked cut-off on resume (\texttt{CLAUDE\_\allowbreak{}CODE\_\allowbreak{}RESUME\_\allowbreak{}INTERRUPTED\_\allowbreak{}TURN}); SIGTERM runs \texttt{SessionEnd} hooks and exits 143; workflow agents restart on stall up to 5$\times$; \texttt{/\allowbreak{}goal} auto-retries 3$\times$ then pauses; usage-limit auto-wait (\texttt{autoContinueAtUsageLimit}); MCP auto-reconnect & \texttt{D/\allowbreak{}errors § Automatic retries}, \texttt{§ The response above may be incomplete}; \texttt{D/\allowbreak{}network-\allowbreak{}config § Streaming idle watchdogs}; \texttt{D/\allowbreak{}model-\allowbreak{}config § Fallback model chains}; \texttt{D/\allowbreak{}how-\allowbreak{}claude-\allowbreak{}code-\allowbreak{}works § When context fills up}; \texttt{D/\allowbreak{}troubleshooting § Auto-\allowbreak{}compaction stops with a thrashing error}; \texttt{D/\allowbreak{}sessions § What a resumed session restores}; \texttt{D/\allowbreak{}headless § Stop a run with SIGTERM}; \texttt{D/\allowbreak{}workflows § When an agent stalls and restarts}; \texttt{D/\allowbreak{}goal § Other errors retry or pause the goal} \\
Performance & Context (layering, deferral, trimming); Model interface (cache TTL); Output (parallel tools); Control flow (background) & Prefix-ordered requests and cache-safe updates via messages (system reminders); MCP schemas deferred; compaction as a forked cached request; \texttt{/\allowbreak{}rewind} truncates to a cached prefix; idle compaction before 1 h TTL expiry; read-only tools run concurrently; background subagents and auto-backgrounded commands; workflow fan-out holds siblings $\leq$5 s so they share the first agent's cache; WebFetch 15-min cache; \texttt{-\allowbreak{}-\allowbreak{}bare} skips discovery; cache hit rate is alerted on internally & \texttt{D/\allowbreak{}prompt-\allowbreak{}caching § How the cache is organized}, \texttt{§ Compacting the conversation}, \texttt{§ Rewinding the conversation}, \texttt{§ Subagents and the cache}; \texttt{D/\allowbreak{}workflows § Prompt caching in a fan-\allowbreak{}out}; \texttt{SDK/\allowbreak{}agent-\allowbreak{}loop § Parallel tool execution}; \texttt{D/\allowbreak{}headless § Start faster with bare mode}; \texttt{B/\allowbreak{}lessons-\allowbreak{}from-\allowbreak{}building-\allowbreak{}claude-\allowbreak{}code-\allowbreak{}prompt-\allowbreak{}caching-\allowbreak{}is-\allowbreak{}everything § LESSONS LEARNED} \\
Evaluation & not documented for the harness itself & No harness test suite or agent evals are documented. Adjacent: \texttt{/\allowbreak{}insights} analyses up to 200 local sessions into a report; skills are evaluated by running prompts with the skill on/off (skill-creator plugin); the auto-mode classifier's first request "validates the Sonnet 5 default"; "Evals" listed as a workflow use case & \texttt{D/\allowbreak{}costs § Analyze your usage patterns}; \texttt{D/\allowbreak{}skills § Evaluating skills}; \texttt{D/\allowbreak{}permission-\allowbreak{}modes § How auto mode evaluates actions}; \texttt{B/\allowbreak{}a-\allowbreak{}harness-\allowbreak{}for-\allowbreak{}every-\allowbreak{}task-\allowbreak{}dynamic-\allowbreak{}workflows-\allowbreak{}in-\allowbreak{}claude-\allowbreak{}code § Use cases} \\
\bottomrule \end{longtable}

\subsubsection*{Design decisions (one line each: the choice this harness made within a primitive)}
\begin{itemize}
\item Input: one entry point for everything that starts a turn (typed prompts, piped stdin, queued mid-turn messages delivered after the current tool batch, scheduled/cron prompts, cross-session and channel messages, background-subagent reports), all passing through \texttt{UserPromptSubmit}. \texttt{D/\allowbreak{}hooks § UserPromptSubmit}; \texttt{D/\allowbreak{}interactive-\allowbreak{}mode § When Claude Code sends what you queued}
\item Context: a stable cached prefix (system prompt + tools, then project context) with every later update appended as a message or \texttt{<system-\allowbreak{}reminder>}; MCP schemas deferred behind \texttt{ToolSearch}; compaction on reaching the window via a forked request that re-injects CLAUDE.md, memory, five recent files and invoked skills. \texttt{D/\allowbreak{}prompt-\allowbreak{}caching § How the cache is organized}; \texttt{D/\allowbreak{}context-\allowbreak{}window § What survives compaction}
\item Model interface: direct Anthropic Messages API with streaming, 10 exponential-backoff retries that never re-issue after a completed tool block, four idle watchdogs, a user-configured fallback model chain, and capability-rejection downgrade for gateways. \texttt{D/\allowbreak{}errors § Automatic retries}; \texttt{D/\allowbreak{}llm-\allowbreak{}gateway-\allowbreak{}protocol § Automatic retry and error forwarding}
\item Output: a fixed built-in tool set plus MCP; read-only calls run concurrently, mutating calls sequentially; shell commands run under an OS sandbox (Seatbelt/bubblewrap) with a local domain-allowlist proxy; large results spill to files; file edits are checkpointed for \texttt{/\allowbreak{}rewind}. \texttt{SDK/\allowbreak{}agent-\allowbreak{}loop § Parallel tool execution}; \texttt{D/\allowbreak{}sandboxing § OS-\allowbreak{}level enforcement}; \texttt{D/\allowbreak{}checkpointing § Automatic tracking}
\item Control flow: loop until a text-only response, with the decision to stop extendable by Stop hooks and \texttt{/\allowbreak{}goal}'s evaluator model; a six-step permission pipeline whose default gate is a second model (auto mode); delegation through an \texttt{Agent} tool (depth 3, 20 concurrent, background by default) and, at scale, a model-written JavaScript workflow the runtime executes. \texttt{SDK/\allowbreak{}agent-\allowbreak{}loop § Turns and messages}; \texttt{SDK/\allowbreak{}permissions § How permissions are evaluated}; \texttt{D/\allowbreak{}workflows § When to use a workflow}
\end{itemize}

\subsubsection*{Hard cases}
\begin{itemize}
\item \textbf{WebFetch's inner model call:} a tool (output) that itself calls the model and decides what the main model sees (a summary, not the page). Resolved by rule 21: the tool call is output; the summarisation is a context decision (and a prompt-injection defence, \texttt{D/\allowbreak{}security § Additional safeguards}); the extra request is model interface. \texttt{D/\allowbreak{}tools-\allowbreak{}reference § WebFetch tool behavior}
\item \textbf{Auto-mode classifier:} an approval gate the harness consumes (control flow) implemented as a separate model request with its own stripped-down context (\texttt{D/\allowbreak{}permission-\allowbreak{}modes § How auto mode evaluates actions}). Resolved: the decision is control flow; the request is model interface; the stripping of tool results is context assembly for that side call. Its "server-side classifier review" variant runs inside the API as part of the session's model requests, i.e. partly outside the harness. \texttt{D/\allowbreak{}permission-\allowbreak{}modes § Server-\allowbreak{}side classifier review}
\item \textbf{Prompt/agent hooks and the \texttt{/\allowbreak{}goal} evaluator:} harness-sequenced model calls that decide whether to go again. Resolved by rule 27 (a program that sequences model calls is control flow). \texttt{D/\allowbreak{}hooks § Prompt-\allowbreak{}based hooks}; \texttt{D/\allowbreak{}goal § How evaluation works}
\item \textbf{Dynamic workflows:} the model writes a JS program at run time whose \texttt{agent(\allowbreak{})} calls sequence model calls $\rightarrow$ control flow (rule 27); the \texttt{Workflow} tool call that launches it is output; the script's intermediate results live "in script variables" rather than any context. \texttt{D/\allowbreak{}workflows § How a workflow runs}
\item \textbf{Compaction:} the trigger (window threshold) and the result (summary replaces history) are context; the summarisation is a model call (model interface); the thrashing stop after "a few attempts" is a loop decision (control flow/resilience). \texttt{D/\allowbreak{}model-\allowbreak{}config § Default auto-\allowbreak{}compact thresholds}; \texttt{D/\allowbreak{}troubleshooting § Auto-\allowbreak{}compaction stops with a thrashing error}
\item \textbf{Context-limit retry with reduced \texttt{max\_\allowbreak{}tokens}:} a retry of a failed request (model interface) that changes the request's parameters (context). Placed under model interface because the trigger is a rejected request and the content is unchanged; noted as a boundary case. \texttt{D/\allowbreak{}errors § Automatic retries}
\item \textbf{Hook rewrites:} \texttt{PreToolUse.\allowbreak{}updatedInput} changes the action before execution (output executor side); \texttt{PostToolUse.\allowbreak{}updatedToolOutput} changes what the model is shown (context). Decomposed per rule 1. \texttt{D/\allowbreak{}hooks § PreToolUse}, \texttt{§ PostToolUse}
\item \textbf{Two kinds of model switch:} \texttt{opusplan} (plan vs execution) and subagent \texttt{model} fields are keyed to the work $\rightarrow$ control flow; \texttt{-\allowbreak{}-\allowbreak{}fallback-\allowbreak{}model} (overloaded/unavailable) and safety-flag fallback $\rightarrow$ model interface (rule 22). \texttt{D/\allowbreak{}model-\allowbreak{}config § Model aliases}, \texttt{§ Fallback model chains}, \texttt{§ Automatic model fallback}
\item \textbf{Subagent output scanning:} modifies a tool result before the model reads it $\rightarrow$ context (what the model is shown) and a safety concern; not output, because nothing is executed. \texttt{D/\allowbreak{}sub-\allowbreak{}agents § Subagent output scanning}
\item \textbf{Checkpoints and session titles:} file snapshots are a user artifact (output) that persists; the title is produced by a cheap-model call keyed to the work (control flow) and stored for the operator/user (persistence). \texttt{D/\allowbreak{}checkpointing § Automatic tracking}; \texttt{D/\allowbreak{}sessions § Name your sessions}
\item \textbf{Queued messages mid-turn:} capturing the message is input; deciding to hand it to the model "as soon as those tool calls finish, within the same turn" is a loop decision (control flow). \texttt{D/\allowbreak{}interactive-\allowbreak{}mode § When Claude Code sends what you queued}
\item \textbf{Agent teams' mailbox:} messages arrive as turns for the recipient (input) but the delivery, idle notifications and task-claim locking are an outer harness's control flow; the classifier reviewing each \texttt{SendMessage} is a safety gate. Depth: lead $\rightarrow$ teammates (no nested teams); results return via idle notifications containing the final answer. \texttt{D/\allowbreak{}agent-\allowbreak{}teams § Context and communication}, \texttt{§ Limitations}
\item \textbf{Server-side classifier review and the server-side prompt-injection probe:} DOES NOT FIT the five primitives cleanly, because the safety decision is made inside the API during the model request rather than by the harness; the harness only consumes the verdict (denies when none arrives after ten responses). \texttt{D/\allowbreak{}permission-\allowbreak{}modes § Server-\allowbreak{}side classifier review}; \texttt{D/\allowbreak{}glossary § Prompt injection}
\item \textbf{\texttt{EndConversation}:} a tool whose only effect is to stop the loop; by rule 21 the call is output and the stop is control flow; its un-removability is a safety choice. \texttt{D/\allowbreak{}tools-\allowbreak{}reference § EndConversation tool behavior}
\end{itemize}

\subsubsection*{Summary row}
\begin{longtable}{@{}p{0.18\linewidth}p{0.18\linewidth}p{0.18\linewidth}p{0.18\linewidth}p{0.18\linewidth}@{}}
\toprule \textbf{Input} & \textbf{Context} & \textbf{Model interface} & \textbf{Output} & \textbf{Control flow} \\ \midrule \endhead
Prompts, stdin, queued messages, events, tool results, AskUserQuestion answers & Cached prefix + appended reminders; CLAUDE.md/memory; deferred MCP; compaction & Messages API streaming; 10 retries, watchdogs, fallback chain & Built-in + MCP tools; sandboxed Bash; concurrent read-only calls; checkpoints & Loop until no tool calls; six-step permission gate; Agent/workflow delegation \\
\bottomrule \end{longtable}


\subsection{OpenAI Codex CLI}

\textbf{Source:} https://github.com/openai/codex $\cdot$ \textbf{Version:} commit f0c5077c34ec7a76b4ad9134835925e6ae3d72ec (2026-10-15) $\cdot$ \textbf{Language:} Rust (workspace under \texttt{codex-\allowbreak{}rs/\allowbreak{}})
\textbf{Path conventions:} all paths are relative to \texttt{codex-\allowbreak{}rs/\allowbreak{}}. \texttt{core/\allowbreak{}} = crate \texttt{codex-\allowbreak{}core}; \texttt{session/\allowbreak{}} = \texttt{core/\allowbreak{}src/\allowbreak{}session/\allowbreak{}}; \texttt{tools/\allowbreak{}} under core = \texttt{core/\allowbreak{}src/\allowbreak{}tools/\allowbreak{}}; top-level crates (\texttt{codex-\allowbreak{}api}, \texttt{sandboxing}, \texttt{rollout}, \texttt{otel}, \texttt{prompts}, \texttt{protocol}, \ldots{}) are siblings of \texttt{core/\allowbreak{}}. Line ranges are \texttt{L<start>-\allowbreak{}-\allowbreak{}L<end>} at the pinned commit.
\textbf{What it is:} Codex is OpenAI's terminal/IDE coding agent. A \texttt{Session} (one thread) receives \texttt{Op} submissions over a channel, runs each user turn as a task that loops "sample $\rightarrow$ run tool calls $\rightarrow$ sample again" until the model stops calling tools, and talks to the OpenAI Responses API over HTTP SSE or a Responses WebSocket. It wraps execution in OS sandboxes (Seatbelt / seccomp+bwrap / Windows restricted token), gates risky actions with user or "Guardian" (model-reviewer) approvals, persists every thread as a JSONL rollout, and can spawn sub-agent threads that are full \texttt{Session}s of their own.

\subsubsection*{Control flow}
\begin{itemize}
\item \textbf{Submission loop (session-level dispatcher):} \texttt{submission\_\allowbreak{}loop} \texttt{select!}s on shutdown, the \texttt{rx\_\allowbreak{}sub} channel and mailbox updates, then matches each \texttt{Op} (TurnInput, Interrupt, ExecApproval, UserInputAnswer, Compact, Review, Shutdown\ldots{}). \texttt{core/\allowbreak{}src/\allowbreak{}session/\allowbreak{}handlers.\allowbreak{}rs:\allowbreak{}L440-\allowbreak{}-\allowbreak{}L727}
\item \textbf{Turn admission modes:} \texttt{turn\_\allowbreak{}input:\allowbreak{}:\allowbreak{}handle} routes a TurnInput by mode: \texttt{StartOrSteer}, \texttt{StartIfIdle}, \texttt{ContinueIfIdle} (recovery requires an internal response item), \texttt{Steer \{ expected\_\allowbreak{}turn\_\allowbreak{}id \}}. \texttt{core/\allowbreak{}src/\allowbreak{}session/\allowbreak{}turn\_\allowbreak{}input.\allowbreak{}rs:\allowbreak{}L230-\allowbreak{}-\allowbreak{}L296}
\item \textbf{Task spawn replaces the active task:} \texttt{Session:\allowbreak{}:\allowbreak{}spawn\_\allowbreak{}task} first calls \texttt{abort\_\allowbreak{}all\_\allowbreak{}tasks(\allowbreak{}TurnAbortReason:\allowbreak{}:\allowbreak{}Replaced)} then \texttt{start\_\allowbreak{}task}, which registers the task on \texttt{active\_\allowbreak{}turn}, seeds pending input from the mailbox, and sets a fresh \texttt{CancellationToken}. \texttt{core/\allowbreak{}src/\allowbreak{}tasks/\allowbreak{}mod.\allowbreak{}rs:\allowbreak{}L273-\allowbreak{}-\allowbreak{}L340}
\item \textbf{Abort / interrupt:} \texttt{abort\_\allowbreak{}all\_\allowbreak{}tasks} takes the active turn, cancels the task and clears pending approvals; on \texttt{Interrupted} it may immediately start a new turn for pending work. \texttt{core/\allowbreak{}src/\allowbreak{}tasks/\allowbreak{}mod.\allowbreak{}rs:\allowbreak{}L564-\allowbreak{}-\allowbreak{}L592}
\item \textbf{Outer task loop (RegularTask::run):} calls \texttt{run\_\allowbreak{}turn} repeatedly until there is no pending input or a terminal error was recorded. \texttt{core/\allowbreak{}src/\allowbreak{}tasks/\allowbreak{}regular.\allowbreak{}rs:\allowbreak{}L105-\allowbreak{}-\allowbreak{}L124}
\item \textbf{Inner turn loop (run\_turn):} \texttt{loop \{ drain pending input $\rightarrow$ capture StepContext $\rightarrow$ build sampling input $\rightarrow$ run\_\allowbreak{}sampling\_\allowbreak{}request $\rightarrow$ decide follow-\allowbreak{}up \}}; \texttt{needs\_\allowbreak{}follow\_\allowbreak{}up = model\_\allowbreak{}needs\_\allowbreak{}follow\_\allowbreak{}up || has\_\allowbreak{}pending\_\allowbreak{}input}; stops with \texttt{break} when \texttt{!needs\_\allowbreak{}follow\_\allowbreak{}up} after stop hooks run. \texttt{core/\allowbreak{}src/\allowbreak{}session/\allowbreak{}turn.\allowbreak{}rs:\allowbreak{}L453-\allowbreak{}-\allowbreak{}L578}, \texttt{L604}, \texttt{L718-\allowbreak{}-\allowbreak{}L821}
\item \textbf{Stop condition from the model:} \texttt{try\_\allowbreak{}run\_\allowbreak{}sampling\_\allowbreak{}request} sets \texttt{needs\_\allowbreak{}follow\_\allowbreak{}up = true} when any output item is a tool call, and on \texttt{ResponseEvent:\allowbreak{}:\allowbreak{}Completed \{ end\_\allowbreak{}turn:\allowbreak{} Some(\allowbreak{}false) \}}. \texttt{core/\allowbreak{}src/\allowbreak{}stream\_\allowbreak{}events\_\allowbreak{}utils.\allowbreak{}rs:\allowbreak{}L351-\allowbreak{}-\allowbreak{}L354}; \texttt{core/\allowbreak{}src/\allowbreak{}session/\allowbreak{}turn.\allowbreak{}rs:\allowbreak{}L3112-\allowbreak{}-\allowbreak{}L3118}
\item \textbf{Stop hooks can force another iteration:} \texttt{run\_\allowbreak{}turn\_\allowbreak{}stop\_\allowbreak{}hooks} result \texttt{should\_\allowbreak{}block} with a continuation prompt records a hook message, sets \texttt{stop\_\allowbreak{}hook\_\allowbreak{}active} and \texttt{continue}s the loop; \texttt{should\_\allowbreak{}stop} breaks. \texttt{core/\allowbreak{}src/\allowbreak{}session/\allowbreak{}turn.\allowbreak{}rs:\allowbreak{}L718-\allowbreak{}-\allowbreak{}L768}
\item \textbf{Mid-turn roll-over on context limit:} \texttt{should\_\allowbreak{}roll\_\allowbreak{}over = auto\_\allowbreak{}compaction\_\allowbreak{}enabled \&\& needs\_\allowbreak{}follow\_\allowbreak{}up \&\& (\allowbreak{}new\_\allowbreak{}context\_\allowbreak{}window requested || token\_\allowbreak{}limit\_\allowbreak{}reached || background candidate ready)} triggers \texttt{run\_\allowbreak{}auto\_\allowbreak{}compact(\allowbreak{}.\allowbreak{}.\allowbreak{}MidTurn.\allowbreak{}.\allowbreak{})} then \texttt{continue}. \texttt{core/\allowbreak{}src/\allowbreak{}session/\allowbreak{}turn.\allowbreak{}rs:\allowbreak{}L654-\allowbreak{}-\allowbreak{}L716}
\item \textbf{Steering (user message while a turn runs):} \texttt{steer\_\allowbreak{}input} validates the active task is a Regular task (Review/Compact are non-steerable), then pushes the input onto \texttt{turn\_\allowbreak{}state.\allowbreak{}pending\_\allowbreak{}input} for the next loop iteration. \texttt{core/\allowbreak{}src/\allowbreak{}session/\allowbreak{}turn\_\allowbreak{}input.\allowbreak{}rs:\allowbreak{}L799-\allowbreak{}-\allowbreak{}L889}
\item \textbf{Pending-input drain point:} \texttt{get\_\allowbreak{}pending\_\allowbreak{}input} splits off \texttt{turn\_\allowbreak{}state.\allowbreak{}pending\_\allowbreak{}input.\allowbreak{}items} and appends mailbox items only if the turn currently accepts mailbox delivery. \texttt{core/\allowbreak{}src/\allowbreak{}session/\allowbreak{}input\_\allowbreak{}queue.\allowbreak{}rs:\allowbreak{}L400-\allowbreak{}-\allowbreak{}L432}
\item \textbf{Sampling retry loop:} \texttt{run\_\allowbreak{}sampling\_\allowbreak{}request} wraps \texttt{try\_\allowbreak{}run\_\allowbreak{}sampling\_\allowbreak{}request} in a \texttt{loop} whose continuation is decided by \texttt{handle\_\allowbreak{}response\_\allowbreak{}stream\_\allowbreak{}error} (retry budget \texttt{stream\_\allowbreak{}max\_\allowbreak{}retries}), and returns early with \texttt{needs\_\allowbreak{}follow\_\allowbreak{}up:\allowbreak{} true} on preempt. \texttt{core/\allowbreak{}src/\allowbreak{}session/\allowbreak{}turn.\allowbreak{}rs:\allowbreak{}L1707-\allowbreak{}-\allowbreak{}L1835}
\item \textbf{Approval gate consumed by harness (user):} \texttt{Session:\allowbreak{}:\allowbreak{}request\_\allowbreak{}command\_\allowbreak{}approval} registers a \texttt{oneshot} under the approval id, emits \texttt{EventMsg:\allowbreak{}:\allowbreak{}ExecApprovalRequest}, and awaits \texttt{rx\_\allowbreak{}approve} (defaulting to \texttt{Abort} if dropped); \texttt{exec\_\allowbreak{}approval} resolves it, and \texttt{Abort} calls \texttt{interrupt\_\allowbreak{}task}. \texttt{core/\allowbreak{}src/\allowbreak{}session/\allowbreak{}mod.\allowbreak{}rs:\allowbreak{}L2923-\allowbreak{}-\allowbreak{}L3012}; \texttt{core/\allowbreak{}src/\allowbreak{}session/\allowbreak{}handlers.\allowbreak{}rs:\allowbreak{}L177-\allowbreak{}-\allowbreak{}L207}
\item \textbf{Approval precedence:} hooks first, then Guardian (model reviewer) if enabled/strict, else user. \texttt{core/\allowbreak{}src/\allowbreak{}tools/\allowbreak{}approvals.\allowbreak{}rs:\allowbreak{}L479-\allowbreak{}-\allowbreak{}L498}, \texttt{L541-\allowbreak{}-\allowbreak{}L557}
\item \textbf{Session-scoped approval cache:} \texttt{with\_\allowbreak{}cached\_\allowbreak{}approval} returns \texttt{ApprovedForSession} without asking when every cache key already holds that decision. \texttt{core/\allowbreak{}src/\allowbreak{}tools/\allowbreak{}sandboxing.\allowbreak{}rs:\allowbreak{}L72-\allowbreak{}-\allowbreak{}L110}
\item \textbf{Sub-agent spawn (effect of \texttt{spawn\_\allowbreak{}agent} tool):} \texttt{thread\_\allowbreak{}spawn:\allowbreak{}:\allowbreak{}spawn} builds a \texttt{SessionSource:\allowbreak{}:\allowbreak{}SubAgent(\allowbreak{}ThreadSpawn\{parent\_\allowbreak{}thread\_\allowbreak{}id, depth, \ldots{}\})} and calls \texttt{agent\_\allowbreak{}control.\allowbreak{}spawn(\allowbreak{}SpawnRequest\{ input:\allowbreak{} AgentInput:\allowbreak{}:\allowbreak{}Message\{ mode:\allowbreak{} TriggerTurn \}, \ldots{} \})}. \texttt{core/\allowbreak{}src/\allowbreak{}thread\_\allowbreak{}spawn.\allowbreak{}rs:\allowbreak{}L58-\allowbreak{}-\allowbreak{}L167}
\item \textbf{Depth and fan-out limits:} \texttt{exceeds\_\allowbreak{}thread\_\allowbreak{}spawn\_\allowbreak{}depth\_\allowbreak{}limit(\allowbreak{}depth, max\_\allowbreak{}depth)} (\texttt{depth > max\_\allowbreak{}depth}) rejects a spawn with "Agent depth limit reached. Solve the task yourself."; \texttt{reserve\_\allowbreak{}spawn\_\allowbreak{}slot} returns \texttt{AgentLimitReached \{ max\_\allowbreak{}threads \}} when the registry is full. \texttt{core/\allowbreak{}src/\allowbreak{}agent/\allowbreak{}registry.\allowbreak{}rs:\allowbreak{}L85-\allowbreak{}-\allowbreak{}L109}; \texttt{core/\allowbreak{}src/\allowbreak{}tools/\allowbreak{}handlers/\allowbreak{}multi\_\allowbreak{}agents/\allowbreak{}spawn.\allowbreak{}rs:\allowbreak{}L74-\allowbreak{}-\allowbreak{}L80}
\item \textbf{Sub-agent results return:} \texttt{wait\_\allowbreak{}agent} polls child status futures with \texttt{timeout\_\allowbreak{}at(\allowbreak{}deadline, \ldots{})} (default \texttt{DEFAULT\_\allowbreak{}WAIT\_\allowbreak{}TIMEOUT\_\allowbreak{}MS}) and collects final statuses; \texttt{send\_\allowbreak{}input} steers or starts a child turn through \texttt{start\_\allowbreak{}or\_\allowbreak{}steer\_\allowbreak{}turn}. \texttt{core/\allowbreak{}src/\allowbreak{}tools/\allowbreak{}handlers/\allowbreak{}multi\_\allowbreak{}agents/\allowbreak{}wait.\allowbreak{}rs:\allowbreak{}L95-\allowbreak{}-\allowbreak{}L99}, \texttt{L181-\allowbreak{}-\allowbreak{}L200}; \texttt{core/\allowbreak{}src/\allowbreak{}agent/\allowbreak{}control.\allowbreak{}rs:\allowbreak{}L137-\allowbreak{}-\allowbreak{}L165}
\item \textbf{Delegate (review / guardian) sub-session:} \texttt{run\_\allowbreak{}codex\_\allowbreak{}thread\_\allowbreak{}one\_\allowbreak{}shot} $\rightarrow$ \texttt{run\_\allowbreak{}codex\_\allowbreak{}thread\_\allowbreak{}interactive} $\rightarrow$ \texttt{Session:\allowbreak{}:\allowbreak{}spawn(\allowbreak{}SessionSpawnArgs\{ thread\_\allowbreak{}source:\allowbreak{} Subagent|GuardianReview, agent\_\allowbreak{}control:\allowbreak{} Provided, \ldots{} \})}; child events/ops are forwarded via \texttt{forward\_\allowbreak{}session\_\allowbreak{}io} tasks. \texttt{core/\allowbreak{}src/\allowbreak{}codex\_\allowbreak{}delegate.\allowbreak{}rs:\allowbreak{}L105-\allowbreak{}-\allowbreak{}L228}, \texttt{L231-\allowbreak{}-\allowbreak{}L262}
\item \textbf{Routing keyed to the work (review model):} the review task clones config, disables web search/collab, sets \texttt{base\_\allowbreak{}instructions = REVIEW\_\allowbreak{}PROMPT}, \texttt{approval\_\allowbreak{}policy = Never}, and \texttt{model = config.\allowbreak{}review\_\allowbreak{}model} before delegating. \texttt{core/\allowbreak{}src/\allowbreak{}tasks/\allowbreak{}review.\allowbreak{}rs:\allowbreak{}L105-\allowbreak{}-\allowbreak{}L140}
\item \textbf{Routing keyed to the work (guardian model):} the reviewer session resolves its own \texttt{review\_\allowbreak{}model} and reasoning effort separately from the parent. \texttt{core/\allowbreak{}src/\allowbreak{}guardian/\allowbreak{}review\_\allowbreak{}session.\allowbreak{}rs:\allowbreak{}L330-\allowbreak{}-\allowbreak{}L372}
\item \textbf{Code mode: model-written program sequences tools (output), but can call back into the harness:} nested calls from a cell go through \texttt{CodeModeCellDelegate:\allowbreak{}:\allowbreak{}invoke\_\allowbreak{}tool} $\rightarrow$ \texttt{DispatchMessage:\allowbreak{}:\allowbreak{}InvokeTool} to the turn worker; the program sequences tools only, not model calls. \texttt{core/\allowbreak{}src/\allowbreak{}tools/\allowbreak{}code\_\allowbreak{}mode/\allowbreak{}delegate.\allowbreak{}rs:\allowbreak{}L317-\allowbreak{}-\allowbreak{}L360}
\item \textbf{Session token budget (stop):} \texttt{record\_\allowbreak{}rollout\_\allowbreak{}budget\_\allowbreak{}usage} returns \texttt{SessionBudgetExceeded} once weighted usage reaches \texttt{limit\_\allowbreak{}tokens}. \texttt{core/\allowbreak{}src/\allowbreak{}agent/\allowbreak{}control/\allowbreak{}budget.\allowbreak{}rs:\allowbreak{}L12-\allowbreak{}-\allowbreak{}L17}; \texttt{core/\allowbreak{}src/\allowbreak{}rollout\_\allowbreak{}budget.\allowbreak{}rs:\allowbreak{}L48-\allowbreak{}-\allowbreak{}L65}
\item \textbf{Model-requested context reset (\texttt{new\_\allowbreak{}context\_\allowbreak{}window} tool):} the handler calls \texttt{session.\allowbreak{}request\_\allowbreak{}new\_\allowbreak{}context\_\allowbreak{}window(\allowbreak{})}, which the run\_turn roll-over check consumes via \texttt{take\_\allowbreak{}new\_\allowbreak{}context\_\allowbreak{}window\_\allowbreak{}request}. \texttt{core/\allowbreak{}src/\allowbreak{}tools/\allowbreak{}handlers/\allowbreak{}new\_\allowbreak{}context\_\allowbreak{}window.\allowbreak{}rs:\allowbreak{}L44-\allowbreak{}-\allowbreak{}L56}; \texttt{core/\allowbreak{}src/\allowbreak{}session/\allowbreak{}turn.\allowbreak{}rs:\allowbreak{}L659-\allowbreak{}-\allowbreak{}L663}
\end{itemize}

\subsubsection*{Input}
\begin{itemize}
\item \textbf{Op channel + public thread API:} \texttt{CodexThread:\allowbreak{}:\allowbreak{}submit(\allowbreak{}op)} sends into \texttt{io.\allowbreak{}tx\_\allowbreak{}sub}; \texttt{next\_\allowbreak{}event(\allowbreak{})} reads \texttt{io.\allowbreak{}rx\_\allowbreak{}event}. \texttt{core/\allowbreak{}src/\allowbreak{}codex\_\allowbreak{}thread.\allowbreak{}rs:\allowbreak{}L258-\allowbreak{}-\allowbreak{}L260}, \texttt{L773-\allowbreak{}-\allowbreak{}L775}
\item \textbf{User message submission op:} \texttt{Op:\allowbreak{}:\allowbreak{}TurnInput \{ request, mode, reply \}} is "Submit turn input using the requested routing behavior." \texttt{protocol/\allowbreak{}src/\allowbreak{}protocol.\allowbreak{}rs:\allowbreak{}L703-\allowbreak{}-\allowbreak{}L710}
\item \textbf{User turn ingestion:} \texttt{start\_\allowbreak{}or\_\allowbreak{}steer} builds \texttt{task\_\allowbreak{}input} from merged additional context plus \texttt{pending\_\allowbreak{}turn\_\allowbreak{}input(\allowbreak{}.\allowbreak{}.\allowbreak{}.\allowbreak{})} and calls \texttt{spawn\_\allowbreak{}task(\allowbreak{}turn\_\allowbreak{}context, task\_\allowbreak{}input, RegularTask:\allowbreak{}:\allowbreak{}new(\allowbreak{}))}. \texttt{core/\allowbreak{}src/\allowbreak{}session/\allowbreak{}turn\_\allowbreak{}input.\allowbreak{}rs:\allowbreak{}L428-\allowbreak{}-\allowbreak{}L449}
\item \textbf{Person's answer the MODEL consumes:} \texttt{request\_\allowbreak{}user\_\allowbreak{}input} tool awaits \texttt{session.\allowbreak{}request\_\allowbreak{}user\_\allowbreak{}input(\allowbreak{}.\allowbreak{}.\allowbreak{}.\allowbreak{})} and serializes the \texttt{RequestUserInputResponse} back as the tool result; delivered by \texttt{Op:\allowbreak{}:\allowbreak{}UserInputAnswer \{ id, response \}}. \texttt{core/\allowbreak{}src/\allowbreak{}tools/\allowbreak{}handlers/\allowbreak{}request\_\allowbreak{}user\_\allowbreak{}input.\allowbreak{}rs:\allowbreak{}L78-\allowbreak{}-\allowbreak{}L100}; \texttt{protocol/\allowbreak{}src/\allowbreak{}protocol.\allowbreak{}rs:\allowbreak{}L791-\allowbreak{}-\allowbreak{}L797}
\item \textbf{Tool results as input to the next request:} in-flight tool futures are drained in order and recorded via \texttt{record\_\allowbreak{}annotated\_\allowbreak{}conversation\_\allowbreak{}items}. \texttt{core/\allowbreak{}src/\allowbreak{}session/\allowbreak{}turn.\allowbreak{}rs:\allowbreak{}L2565-\allowbreak{}-\allowbreak{}L2592}
\item \textbf{Inter-agent mail / mailbox:} \texttt{Op:\allowbreak{}:\allowbreak{}InterAgentCommunication} and mailbox updates (\texttt{mailbox\_\allowbreak{}updates}, \texttt{notify\_\allowbreak{}mailbox}) feed pending input; mailbox delivery can preempt a response between items. \texttt{core/\allowbreak{}src/\allowbreak{}session/\allowbreak{}handlers.\allowbreak{}rs:\allowbreak{}L447-\allowbreak{}-\allowbreak{}L479}, \texttt{L649-\allowbreak{}-\allowbreak{}L656}; \texttt{core/\allowbreak{}src/\allowbreak{}session/\allowbreak{}turn.\allowbreak{}rs:\allowbreak{}L2913-\allowbreak{}-\allowbreak{}L2932}
\item \textbf{Hook-contributed input:} \texttt{drain\_\allowbreak{}async\_\allowbreak{}hook\_\allowbreak{}results} records results of async hooks before/after the user prompt. \texttt{core/\allowbreak{}src/\allowbreak{}hook\_\allowbreak{}runtime.\allowbreak{}rs:\allowbreak{}L782}; \texttt{core/\allowbreak{}src/\allowbreak{}session/\allowbreak{}turn.\allowbreak{}rs:\allowbreak{}L185}, \texttt{L591}
\end{itemize}

\subsubsection*{Context}
\begin{itemize}
\item \textbf{Base instructions come from the model catalog:} \texttt{render\_\allowbreak{}model\_\allowbreak{}instructions} reads \texttt{model\_\allowbreak{}info.\allowbreak{}model\_\allowbreak{}messages.\allowbreak{}instructions\_\allowbreak{}template}; the bundled default is \texttt{BASE\_\allowbreak{}INSTRUCTIONS = include\_\allowbreak{}str!(\allowbreak{}".\allowbreak{}.\allowbreak{}/\allowbreak{}prompt.\allowbreak{}md")} in models-manager, overridable by \texttt{config.\allowbreak{}base\_\allowbreak{}instructions}. \texttt{prompts/\allowbreak{}src/\allowbreak{}model\_\allowbreak{}instructions.\allowbreak{}rs:\allowbreak{}L8-\allowbreak{}-\allowbreak{}L17}; \texttt{models-\allowbreak{}manager/\allowbreak{}src/\allowbreak{}model\_\allowbreak{}info.\allowbreak{}rs:\allowbreak{}L16}, \texttt{L46-\allowbreak{}-\allowbreak{}L48}
\item \textbf{Prompt assembly per sampling request:} \texttt{build\_\allowbreak{}prompt} packs \texttt{input} (history), \texttt{tools = tool\_\allowbreak{}router.\allowbreak{}model\_\allowbreak{}visible\_\allowbreak{}specs(\allowbreak{})}, \texttt{parallel\_\allowbreak{}tool\_\allowbreak{}calls:\allowbreak{} true}, \texttt{base\_\allowbreak{}instructions}, \texttt{output\_\allowbreak{}schema}; with incremental tools (WebSocket continuation) tools and instructions are sent empty. \texttt{core/\allowbreak{}src/\allowbreak{}session/\allowbreak{}turn.\allowbreak{}rs:\allowbreak{}L1666-\allowbreak{}-\allowbreak{}L1694}
\item \textbf{History normalization before send:} \texttt{for\_\allowbreak{}prompt} $\rightarrow$ \texttt{normalize\_\allowbreak{}history}: ensure every call has an output, drop orphan outputs, strip images/audio the model can't take. \texttt{core/\allowbreak{}src/\allowbreak{}context\_\allowbreak{}manager/\allowbreak{}history.\allowbreak{}rs:\allowbreak{}L603-\allowbreak{}-\allowbreak{}L617}, \texttt{L968-\allowbreak{}-\allowbreak{}L982}
\item \textbf{Tool-output truncation on record:} \texttt{record\_\allowbreak{}item\_\allowbreak{}with\_\allowbreak{}metadata} applies \texttt{truncate\_\allowbreak{}function\_\allowbreak{}output\_\allowbreak{}payload(\allowbreak{}output, policy, \ldots{})} to Function/CustomToolCallOutput using the model's \texttt{truncation\_\allowbreak{}policy} (live history only; rollout keeps full payload). \texttt{core/\allowbreak{}src/\allowbreak{}context\_\allowbreak{}manager/\allowbreak{}history.\allowbreak{}rs:\allowbreak{}L517-\allowbreak{}-\allowbreak{}L519}, \texttt{L568-\allowbreak{}-\allowbreak{}L577}
\item \textbf{World-state sections (diff-injected context):} \texttt{WorldStateSection} trait --- typed sections (\texttt{agents\_\allowbreak{}md}, \texttt{environment}, \texttt{permissions}, \texttt{tools}, \texttt{model}, \ldots{}) each render a diff vs. the previous snapshot; \texttt{WorldStateSnapshot:\allowbreak{}:\allowbreak{}merge\_\allowbreak{}patch\_\allowbreak{}from} produces an RFC 7386 merge patch. \texttt{core/\allowbreak{}src/\allowbreak{}context/\allowbreak{}world\_\allowbreak{}state/\allowbreak{}mod.\allowbreak{}rs:\allowbreak{}L302-\allowbreak{}-\allowbreak{}L330}, \texttt{L381-\allowbreak{}-\allowbreak{}L425}
\item \textbf{Per-step world-state injection:} \texttt{record\_\allowbreak{}step\_\allowbreak{}world\_\allowbreak{}state\_\allowbreak{}if\_\allowbreak{}changed} renders the step's world state, records only the changed fragments into history, and persists a \texttt{RolloutItem:\allowbreak{}:\allowbreak{}WorldState} patch. \texttt{core/\allowbreak{}src/\allowbreak{}session/\allowbreak{}mod.\allowbreak{}rs:\allowbreak{}L3762-\allowbreak{}-\allowbreak{}L3793}
\item \textbf{AGENTS.md discovery:} walks from project root (marker \texttt{.\allowbreak{}git} by default) down to cwd, concatenating every \texttt{AGENTS.\allowbreak{}md}, bounded by \texttt{project\_\allowbreak{}doc\_\allowbreak{}max\_\allowbreak{}bytes}; skipped for untrusted projects. \texttt{core/\allowbreak{}src/\allowbreak{}agents\_\allowbreak{}md.\allowbreak{}rs:\allowbreak{}L1-\allowbreak{}-\allowbreak{}L18}, \texttt{L58-\allowbreak{}-\allowbreak{}L68}; refreshed per step in \texttt{capture\_\allowbreak{}step\_\allowbreak{}context\_\allowbreak{}inner}. \texttt{core/\allowbreak{}src/\allowbreak{}session/\allowbreak{}mod.\allowbreak{}rs:\allowbreak{}L3915-\allowbreak{}-\allowbreak{}L3924}
\item \textbf{AGENTS.md as a world-state section:} \texttt{AgentsMdState} renders replacement/removal notices when the instructions change. \texttt{core/\allowbreak{}src/\allowbreak{}context/\allowbreak{}world\_\allowbreak{}state/\allowbreak{}agents\_\allowbreak{}md.\allowbreak{}rs:\allowbreak{}L11-\allowbreak{}-\allowbreak{}L60}
\item \textbf{Environment context:} \texttt{FileSystemContext \{ workspace\_\allowbreak{}roots, permission\_\allowbreak{}profile \}} and \texttt{NetworkContext} are rendered into the model-visible environment section. \texttt{core/\allowbreak{}src/\allowbreak{}context/\allowbreak{}environment\_\allowbreak{}context.\allowbreak{}rs:\allowbreak{}L10-\allowbreak{}-\allowbreak{}L30}, \texttt{L213}
\item \textbf{Developer instructions + extension contributors:} \texttt{build\_\allowbreak{}initial\_\allowbreak{}context\_\allowbreak{}with\_\allowbreak{}world\_\allowbreak{}state} collects developer sections, plugin recommendations, and \texttt{context\_\allowbreak{}contributors(\allowbreak{}).\allowbreak{}contribute\_\allowbreak{}thread\_\allowbreak{}context/\allowbreak{}turn\_\allowbreak{}context} fragments into developer/user slots. \texttt{core/\allowbreak{}src/\allowbreak{}session/\allowbreak{}mod.\allowbreak{}rs:\allowbreak{}L4473-\allowbreak{}-\allowbreak{}L4560}
\item \textbf{Skills injected on mention:} \texttt{build\_\allowbreak{}skills\_\allowbreak{}and\_\allowbreak{}plugins} loads prompts for explicitly mentioned skills (\texttt{load\_\allowbreak{}skill\_\allowbreak{}prompts}) and converts each fragment into a response item recorded into history. \texttt{core/\allowbreak{}src/\allowbreak{}session/\allowbreak{}turn.\allowbreak{}rs:\allowbreak{}L1099}, \texttt{L1165-\allowbreak{}-\allowbreak{}L1190}, \texttt{L1222-\allowbreak{}-\allowbreak{}L1236}; recorded at \texttt{L409-\allowbreak{}-\allowbreak{}L416}
\item \textbf{Memory read-path context:} \texttt{MemoryContextFragment:\allowbreak{}:\allowbreak{}\{ReadInstructions (\allowbreak{}developer), ExtractionEvidence (\allowbreak{}user)\}}, capped below 10k tokens; memory root is \texttt{\$CODEX\_\allowbreak{}HOME/\allowbreak{}memories}. \texttt{core/\allowbreak{}src/\allowbreak{}context/\allowbreak{}memory.\allowbreak{}rs:\allowbreak{}L1-\allowbreak{}-\allowbreak{}L30}; \texttt{memories/\allowbreak{}read/\allowbreak{}src/\allowbreak{}lib.\allowbreak{}rs:\allowbreak{}L13}
\item \textbf{Tool catalog for the step:} \texttt{built\_\allowbreak{}tools} builds a \texttt{ToolRouter} from MCP tools, connectors and tool-suggest candidates; \texttt{capture\_\allowbreak{}step\_\allowbreak{}context\_\allowbreak{}inner} joins MCP refresh and tool build into a \texttt{StepContext}. \texttt{core/\allowbreak{}src/\allowbreak{}session/\allowbreak{}turn.\allowbreak{}rs:\allowbreak{}L1884-\allowbreak{}-\allowbreak{}L1973}; \texttt{core/\allowbreak{}src/\allowbreak{}session/\allowbreak{}mod.\allowbreak{}rs:\allowbreak{}L3959-\allowbreak{}-\allowbreak{}L4030}
\item \textbf{Deferred tools (tool search):} \texttt{ToolExposure:\allowbreak{}:\allowbreak{}Deferred} tools are omitted from the initial list and discoverable through the \texttt{tool\_\allowbreak{}search} tool. \texttt{tools/\allowbreak{}src/\allowbreak{}tool\_\allowbreak{}executor.\allowbreak{}rs:\allowbreak{}L48-\allowbreak{}-\allowbreak{}L62}; \texttt{core/\allowbreak{}src/\allowbreak{}tools/\allowbreak{}handlers/\allowbreak{}tool\_\allowbreak{}search.\allowbreak{}rs:\allowbreak{}L68-\allowbreak{}-\allowbreak{}L85}
\item \textbf{Compaction trigger math:} \texttt{context\_\allowbreak{}window\_\allowbreak{}token\_\allowbreak{}status} computes \texttt{token\_\allowbreak{}limit\_\allowbreak{}reached} (auto-compact limit + fallback buffer, or \texttt{effective\_\allowbreak{}context\_\allowbreak{}window\_\allowbreak{}percent} of the window) and \texttt{turn\_\allowbreak{}end\_\allowbreak{}compaction\_\allowbreak{}threshold\_\allowbreak{}reached} (\texttt{model\_\allowbreak{}post\_\allowbreak{}turn\_\allowbreak{}compact\_\allowbreak{}threshold\_\allowbreak{}percent}). \texttt{core/\allowbreak{}src/\allowbreak{}session/\allowbreak{}context\_\allowbreak{}window.\allowbreak{}rs:\allowbreak{}L53-\allowbreak{}-\allowbreak{}L130}
\item \textbf{Pre-turn / post-turn compaction:} \texttt{run\_\allowbreak{}pre\_\allowbreak{}sampling\_\allowbreak{}compact} runs before the user input is recorded; post-turn compaction runs when \texttt{foreground\_\allowbreak{}due || candidate.\allowbreak{}is\_\allowbreak{}some(\allowbreak{})}. \texttt{core/\allowbreak{}src/\allowbreak{}session/\allowbreak{}turn.\allowbreak{}rs:\allowbreak{}L192-\allowbreak{}-\allowbreak{}L212}, \texttt{L779-\allowbreak{}-\allowbreak{}L817}
\item \textbf{Compaction implementation selection:} \texttt{run\_\allowbreak{}auto\_\allowbreak{}compact} chooses token-budget reset (\texttt{Feature:\allowbreak{}:\allowbreak{}TokenBudget}), remote \texttt{ResponsesCompactionV2}, or local summarization. \texttt{core/\allowbreak{}src/\allowbreak{}session/\allowbreak{}turn.\allowbreak{}rs:\allowbreak{}L1507-\allowbreak{}-\allowbreak{}L1570}; \texttt{core/\allowbreak{}src/\allowbreak{}compact\_\allowbreak{}remote\_\allowbreak{}v2.\allowbreak{}rs:\allowbreak{}L112-\allowbreak{}-\allowbreak{}L140}
\item \textbf{Local summarization compaction:} \texttt{run\_\allowbreak{}compact\_\allowbreak{}task\_\allowbreak{}inner\_\allowbreak{}impl} sends history + \texttt{SUMMARIZATION\_\allowbreak{}PROMPT} (\texttt{prompts/\allowbreak{}templates/\allowbreak{}compact/\allowbreak{}prompt.\allowbreak{}md}) to the model, then \texttt{build\_\allowbreak{}compacted\_\allowbreak{}history} keeps initial context + up to \texttt{COMPACT\_\allowbreak{}USER\_\allowbreak{}MESSAGE\_\allowbreak{}MAX\_\allowbreak{}TOKENS = 20\_\allowbreak{}000} tokens of recent user messages + the summary. \texttt{core/\allowbreak{}src/\allowbreak{}compact.\allowbreak{}rs:\allowbreak{}L62-\allowbreak{}-\allowbreak{}L64}, \texttt{L264-\allowbreak{}-\allowbreak{}L330}, \texttt{L702-\allowbreak{}-\allowbreak{}L760}
\item \textbf{Background (speculative) compaction:} \texttt{async\_\allowbreak{}compaction:\allowbreak{}:\allowbreak{}maybe\_\allowbreak{}start} launches a parallel compaction once usage passes \texttt{async\_\allowbreak{}compaction\_\allowbreak{}threshold\_\allowbreak{}percent} of the window, deduped by history version/model/hash. \texttt{core/\allowbreak{}src/\allowbreak{}session/\allowbreak{}async\_\allowbreak{}compaction.\allowbreak{}rs:\allowbreak{}L1-\allowbreak{}-\allowbreak{}L9}, \texttt{L62-\allowbreak{}-\allowbreak{}L95}
\item \textbf{Model-queryable context remaining (\texttt{get\_\allowbreak{}context\_\allowbreak{}remaining} tool):} reports \texttt{base\_\allowbreak{}window\_\allowbreak{}tokens\_\allowbreak{}remaining} from \texttt{context\_\allowbreak{}window\_\allowbreak{}token\_\allowbreak{}status}. \texttt{core/\allowbreak{}src/\allowbreak{}tools/\allowbreak{}handlers/\allowbreak{}get\_\allowbreak{}context\_\allowbreak{}remaining.\allowbreak{}rs:\allowbreak{}L85-\allowbreak{}-\allowbreak{}L92}
\end{itemize}

\subsubsection*{Model interface}
\begin{itemize}
\item \textbf{Client lifetimes:} \texttt{ModelClient} lives for the session; \texttt{ModelClientSession} is created per turn, caches a Responses WebSocket and the \texttt{x-\allowbreak{}codex-\allowbreak{}turn-\allowbreak{}state} sticky-routing token. \texttt{core/\allowbreak{}src/\allowbreak{}client.\allowbreak{}rs:\allowbreak{}L1-\allowbreak{}-\allowbreak{}L26}, \texttt{L295-\allowbreak{}-\allowbreak{}L303}
\item \textbf{Wire format:} \texttt{ResponsesApiRequest \{ model, input, tools, tool\_\allowbreak{}choice:\allowbreak{} "auto", parallel\_\allowbreak{}tool\_\allowbreak{}calls, reasoning, store:\allowbreak{} false, stream:\allowbreak{} true, include:\allowbreak{} ["reasoning.\allowbreak{}encrypted\_\allowbreak{}content", \ldots{}], prompt\_\allowbreak{}cache\_\allowbreak{}key, text, \ldots{} \}} posted to \texttt{/\allowbreak{}responses}. \texttt{core/\allowbreak{}src/\allowbreak{}client.\allowbreak{}rs:\allowbreak{}L1040-\allowbreak{}-\allowbreak{}L1095}; \texttt{codex-\allowbreak{}api/\allowbreak{}src/\allowbreak{}endpoint/\allowbreak{}responses.\allowbreak{}rs:\allowbreak{}L110-\allowbreak{}-\allowbreak{}L139}
\item \textbf{Transport selection + fallback:} \texttt{ModelClientSession:\allowbreak{}:\allowbreak{}stream} tries \texttt{stream\_\allowbreak{}responses\_\allowbreak{}websocket} when the provider supports it and falls back to \texttt{ResponsesHttpExecutor:\allowbreak{}:\allowbreak{}stream\_\allowbreak{}responses\_\allowbreak{}http}; \texttt{try\_\allowbreak{}switch\_\allowbreak{}fallback\_\allowbreak{}transport} permanently disables WebSockets for the session. \texttt{core/\allowbreak{}src/\allowbreak{}client.\allowbreak{}rs:\allowbreak{}L2363-\allowbreak{}-\allowbreak{}L2435}
\item \textbf{SSE event parsing + idle timeout:} \texttt{process\_\allowbreak{}sse} maps \texttt{response.\allowbreak{}output\_\allowbreak{}item.\allowbreak{}done}, \texttt{response.\allowbreak{}output\_\allowbreak{}text.\allowbreak{}delta}, \texttt{response.\allowbreak{}completed|incomplete} to \texttt{ResponseEvent}s; each poll is \texttt{timeout(\allowbreak{}idle\_\allowbreak{}timeout, stream.\allowbreak{}next(\allowbreak{}))} with default \texttt{DEFAULT\_\allowbreak{}STREAM\_\allowbreak{}IDLE\_\allowbreak{}TIMEOUT\_\allowbreak{}MS = 300\_\allowbreak{}000}. \texttt{codex-\allowbreak{}api/\allowbreak{}src/\allowbreak{}sse/\allowbreak{}responses.\allowbreak{}rs:\allowbreak{}L345-\allowbreak{}-\allowbreak{}L470}, \texttt{L554-\allowbreak{}-\allowbreak{}L558}; \texttt{model-\allowbreak{}provider-\allowbreak{}info/\allowbreak{}src/\allowbreak{}lib.\allowbreak{}rs:\allowbreak{}L66}, \texttt{L513-\allowbreak{}-\allowbreak{}L516}
\item \textbf{Retry policy:} \texttt{handle\_\allowbreak{}response\_\allowbreak{}stream\_\allowbreak{}error}: terminal errors (context exceeded, usage limit, invalid request\ldots{}) return \texttt{None} delay; retryable ones use server \texttt{Retry-\allowbreak{}After} else exponential backoff with ±10\% jitter; budget \texttt{DEFAULT\_\allowbreak{}STREAM\_\allowbreak{}MAX\_\allowbreak{}RETRIES = 5}; unbounded connection retries (5s$\rightarrow$60s doubling) behind \texttt{Feature:\allowbreak{}:\allowbreak{}UnboundedConnectionRetries}; after exhausting the budget on WebSocket it switches to HTTPS and resets the counter. \texttt{core/\allowbreak{}src/\allowbreak{}responses\_\allowbreak{}retry.\allowbreak{}rs:\allowbreak{}L23-\allowbreak{}-\allowbreak{}L24}, \texttt{L57-\allowbreak{}-\allowbreak{}L176}; \texttt{protocol/\allowbreak{}src/\allowbreak{}error.\allowbreak{}rs:\allowbreak{}L410-\allowbreak{}-\allowbreak{}L440}; \texttt{codex-\allowbreak{}client/\allowbreak{}src/\allowbreak{}retry.\allowbreak{}rs:\allowbreak{}L43-\allowbreak{}-\allowbreak{}L52}; \texttt{model-\allowbreak{}provider-\allowbreak{}info/\allowbreak{}src/\allowbreak{}lib.\allowbreak{}rs:\allowbreak{}L67}, \texttt{L506-\allowbreak{}-\allowbreak{}L508}
\item \textbf{Content-filter retry injects guidance (context effect):} on \texttt{ContentFilter} the retry path records a \texttt{ContentFilterGuidance} user fragment before retrying. \texttt{core/\allowbreak{}src/\allowbreak{}responses\_\allowbreak{}retry.\allowbreak{}rs:\allowbreak{}L67-\allowbreak{}-\allowbreak{}L82}
\item \textbf{Receiving the stream:} \texttt{try\_\allowbreak{}run\_\allowbreak{}sampling\_\allowbreak{}request} opens \texttt{client\_\allowbreak{}session.\allowbreak{}stream(\allowbreak{}prompt, model\_\allowbreak{}info, \ldots{})} and loops on \texttt{stream.\allowbreak{}next(\allowbreak{})} under preempt/cancellation. \texttt{core/\allowbreak{}src/\allowbreak{}session/\allowbreak{}turn.\allowbreak{}rs:\allowbreak{}L2667-\allowbreak{}-\allowbreak{}L2715}
\item \textbf{Token usage from the wire:} on \texttt{Completed \{ token\_\allowbreak{}usage, \ldots{} \}} the session records usage (\texttt{record\_\allowbreak{}token\_\allowbreak{}usage\_\allowbreak{}info}) and emits a token-count event. \texttt{core/\allowbreak{}src/\allowbreak{}session/\allowbreak{}turn.\allowbreak{}rs:\allowbreak{}L3069-\allowbreak{}-\allowbreak{}L3118}, \texttt{L3295-\allowbreak{}-\allowbreak{}L3301}
\item \textbf{Auth-expiry recovery:} \texttt{PendingUnauthorizedRetry} retries a request once after refreshing credentials on 401. \texttt{core/\allowbreak{}src/\allowbreak{}client.\allowbreak{}rs:\allowbreak{}L1582-\allowbreak{}-\allowbreak{}L1636}, \texttt{L1781-\allowbreak{}-\allowbreak{}L1925}
\item \textbf{Prompt cache key:} \texttt{prompt\_\allowbreak{}cache\_\allowbreak{}key} is sent per request (override or thread-derived) so the server reuses cached prefixes. \texttt{core/\allowbreak{}src/\allowbreak{}client.\allowbreak{}rs:\allowbreak{}L646-\allowbreak{}-\allowbreak{}L668}, \texttt{L1060}, \texttt{L1090}
\end{itemize}

\subsubsection*{Output}
\begin{itemize}
\item \textbf{Tool call recognition:} \texttt{ToolRouter:\allowbreak{}:\allowbreak{}build\_\allowbreak{}tool\_\allowbreak{}call} turns \texttt{FunctionCall}, \texttt{CustomToolCall}, client-executed \texttt{ToolSearchCall}, etc. into a \texttt{ToolCall}. \texttt{core/\allowbreak{}src/\allowbreak{}tools/\allowbreak{}router.\allowbreak{}rs:\allowbreak{}L257-\allowbreak{}-\allowbreak{}L335}
\item \textbf{Dispatch:} \texttt{handle\_\allowbreak{}output\_\allowbreak{}item\_\allowbreak{}done} logs the call, persists the item, and pushes \texttt{tool\_\allowbreak{}runtime.\allowbreak{}handle\_\allowbreak{}tool\_\allowbreak{}call(\allowbreak{}call, token)} into the in-flight queue; non-tool items become turn items and are emitted to the client. \texttt{core/\allowbreak{}src/\allowbreak{}stream\_\allowbreak{}events\_\allowbreak{}utils.\allowbreak{}rs:\allowbreak{}L315-\allowbreak{}-\allowbreak{}L432}
\item \textbf{Concurrency of tool calls from one response:} \texttt{ToolCallRuntime:\allowbreak{}:\allowbreak{}handle\_\allowbreak{}tool\_\allowbreak{}call} spawns each call on tokio; parallel-safe tools take a read lock on \texttt{parallel\_\allowbreak{}execution:\allowbreak{} RwLock<(\allowbreak{})>}, others a write lock (serialized). \texttt{core/\allowbreak{}src/\allowbreak{}tools/\allowbreak{}parallel.\allowbreak{}rs:\allowbreak{}L51}, \texttt{L141-\allowbreak{}-\allowbreak{}L237}; \texttt{core/\allowbreak{}src/\allowbreak{}tools/\allowbreak{}router.\allowbreak{}rs:\allowbreak{}L246-\allowbreak{}-\allowbreak{}L250}
\item \textbf{Streaming text to the person:} \texttt{OutputTextDelta} is parsed and emitted as assistant text deltas; the TUI appends them via \texttt{on\_\allowbreak{}agent\_\allowbreak{}message\_\allowbreak{}delta} $\rightarrow$ \texttt{handle\_\allowbreak{}streaming\_\allowbreak{}delta}. \texttt{core/\allowbreak{}src/\allowbreak{}session/\allowbreak{}turn.\allowbreak{}rs:\allowbreak{}L3120-\allowbreak{}-\allowbreak{}L3135}; \texttt{tui/\allowbreak{}src/\allowbreak{}chatwidget/\allowbreak{}streaming.\allowbreak{}rs:\allowbreak{}L206-\allowbreak{}-\allowbreak{}L208}
\item \textbf{Diff artifact for the user:} \texttt{TurnDiffTracker} accumulates file changes and emits \texttt{EventMsg:\allowbreak{}:\allowbreak{}TurnDiff \{ unified\_\allowbreak{}diff \}} after each response. \texttt{core/\allowbreak{}src/\allowbreak{}session/\allowbreak{}turn.\allowbreak{}rs:\allowbreak{}L441-\allowbreak{}-\allowbreak{}L443}, \texttt{L3307-\allowbreak{}-\allowbreak{}L3317}
\item \textbf{Exec tool (\texttt{exec\_\allowbreak{}command} / \texttt{write\_\allowbreak{}stdin}):} Unified Exec manages PTY processes with output caps (\texttt{DEFAULT\_\allowbreak{}MAX\_\allowbreak{}OUTPUT\_\allowbreak{}TOKENS = 10\_\allowbreak{}000}, \texttt{UNIFIED\_\allowbreak{}EXEC\_\allowbreak{}OUTPUT\_\allowbreak{}MAX\_\allowbreak{}BYTES = 1 MiB}) and yield-time bounds. \texttt{core/\allowbreak{}src/\allowbreak{}unified\_\allowbreak{}exec/\allowbreak{}mod.\allowbreak{}rs:\allowbreak{}L1-\allowbreak{}-\allowbreak{}L25}, \texttt{L68-\allowbreak{}-\allowbreak{}L76}; \texttt{core/\allowbreak{}src/\allowbreak{}tools/\allowbreak{}handlers/\allowbreak{}unified\_\allowbreak{}exec/\allowbreak{}exec\_\allowbreak{}command.\allowbreak{}rs:\allowbreak{}L109}; \texttt{core/\allowbreak{}src/\allowbreak{}tools/\allowbreak{}handlers/\allowbreak{}unified\_\allowbreak{}exec/\allowbreak{}write\_\allowbreak{}stdin.\allowbreak{}rs:\allowbreak{}L39}
\item \textbf{Exec timeout:} \texttt{DEFAULT\_\allowbreak{}EXEC\_\allowbreak{}COMMAND\_\allowbreak{}TIMEOUT\_\allowbreak{}MS = 10\_\allowbreak{}000}; \texttt{ExecExpiration:\allowbreak{}:\allowbreak{}\{Timeout, DefaultTimeout, Cancellation, TimeoutOrCancellation\}}. \texttt{core/\allowbreak{}src/\allowbreak{}exec.\allowbreak{}rs:\allowbreak{}L63}, \texttt{L173-\allowbreak{}-\allowbreak{}L223}
\item \textbf{Orchestrator (approval $\rightarrow$ sandbox $\rightarrow$ run $\rightarrow$ retry):} \texttt{ExecApprovalRequirement:\allowbreak{}:\allowbreak{}\{Skip, Forbidden, NeedsApproval\}} decides whether to ask; on sandbox denial it retries without a sandbox only if policy allows (\texttt{Never}/\texttt{OnRequest} do not retry). \texttt{core/\allowbreak{}src/\allowbreak{}tools/\allowbreak{}orchestrator.\allowbreak{}rs:\allowbreak{}L164-\allowbreak{}-\allowbreak{}L222}, \texttt{L439-\allowbreak{}-\allowbreak{}L470}
\item \textbf{Sandbox technology:} \texttt{get\_\allowbreak{}platform\_\allowbreak{}sandbox} $\rightarrow$ macOS \texttt{MacosSeatbelt}, Linux \texttt{LinuxSeccomp} (bwrap + seccomp + Landlock), Windows \texttt{WindowsRestrictedToken}/\texttt{WindowsMxc}. \texttt{sandboxing/\allowbreak{}src/\allowbreak{}manager.\allowbreak{}rs:\allowbreak{}L50-\allowbreak{}-\allowbreak{}L63}; \texttt{protocol/\allowbreak{}src/\allowbreak{}sandbox.\allowbreak{}rs:\allowbreak{}L29-\allowbreak{}-\allowbreak{}L35}; \texttt{sandboxing/\allowbreak{}src/\allowbreak{}\{seatbelt.\allowbreak{}rs,bwrap.\allowbreak{}rs,landlock.\allowbreak{}rs\}}; \texttt{linux-\allowbreak{}sandbox/\allowbreak{}src/\allowbreak{}seccomp.\allowbreak{}rs}
\item \textbf{Remote executor:} \texttt{codex-\allowbreak{}exec-\allowbreak{}server} is a JSON-RPC server for spawning/controlling subprocesses via \texttt{codex-\allowbreak{}utils-\allowbreak{}pty}, used for non-local environments. \texttt{exec-\allowbreak{}server/\allowbreak{}README.\allowbreak{}md} (top), \texttt{exec-\allowbreak{}server/\allowbreak{}src/\allowbreak{}lib.\allowbreak{}rs:\allowbreak{}L1-\allowbreak{}-\allowbreak{}L12}
\item \textbf{Patch tool safety:} \texttt{assess\_\allowbreak{}patch\_\allowbreak{}safety} auto-approves patches confined to writable roots (still run in a sandbox), rejects under \texttt{Never}, otherwise asks. \texttt{core/\allowbreak{}src/\allowbreak{}safety.\allowbreak{}rs:\allowbreak{}L67-\allowbreak{}-\allowbreak{}L125}, \texttt{L146-\allowbreak{}-\allowbreak{}L180}
\item \textbf{MCP tool calls:} \texttt{handle\_\allowbreak{}mcp\_\allowbreak{}tool\_\allowbreak{}call} prepares the call, applies approval policy, and invokes the server via \texttt{prepared\_\allowbreak{}call.\allowbreak{}call\_\allowbreak{}with\_\allowbreak{}preparation}. \texttt{core/\allowbreak{}src/\allowbreak{}mcp\_\allowbreak{}tool\_\allowbreak{}call.\allowbreak{}rs:\allowbreak{}L131-\allowbreak{}-\allowbreak{}L145}, \texttt{L473-\allowbreak{}-\allowbreak{}L478}
\item \textbf{Code mode executor (\texttt{exec} tool):} the \texttt{exec} tool runs model-written JS in a V8 host process reached over gRPC; nested tool calls are brokered back to the harness. \texttt{code-\allowbreak{}mode-\allowbreak{}protocol/\allowbreak{}src/\allowbreak{}lib.\allowbreak{}rs:\allowbreak{}L54}; \texttt{core/\allowbreak{}src/\allowbreak{}code\_\allowbreak{}mode\_\allowbreak{}host.\allowbreak{}rs:\allowbreak{}L1-\allowbreak{}-\allowbreak{}L4}; \texttt{core/\allowbreak{}src/\allowbreak{}tools/\allowbreak{}code\_\allowbreak{}mode/\allowbreak{}execute\_\allowbreak{}handler.\allowbreak{}rs:\allowbreak{}L176-\allowbreak{}-\allowbreak{}L179}
\item \textbf{Sub-agent tool surface:} \texttt{spawn\_\allowbreak{}agent}, \texttt{send\_\allowbreak{}input}, \texttt{wait\_\allowbreak{}agent}, \texttt{resume\_\allowbreak{}agent}, \texttt{close\_\allowbreak{}agent} handlers. \texttt{core/\allowbreak{}src/\allowbreak{}tools/\allowbreak{}handlers/\allowbreak{}multi\_\allowbreak{}agents/\allowbreak{}\{spawn.\allowbreak{}rs,send\_\allowbreak{}input.\allowbreak{}rs,wait.\allowbreak{}rs,resume\_\allowbreak{}agent.\allowbreak{}rs,close\_\allowbreak{}agent.\allowbreak{}rs\}}
\item \textbf{Review sub-agent as a tool-free delegate:} results come back as events through \texttt{forward\_\allowbreak{}events}. \texttt{core/\allowbreak{}src/\allowbreak{}codex\_\allowbreak{}delegate.\allowbreak{}rs:\allowbreak{}L200-\allowbreak{}-\allowbreak{}L228}
\end{itemize}

\subsubsection*{Production concerns}
\begin{longtable}{@{}p{0.13\linewidth}p{0.17\linewidth}p{0.36\linewidth}p{0.26\linewidth}@{}}
\toprule \textbf{Concern} & \textbf{Where it lives (primitive)} & \textbf{Mechanism} & \textbf{Evidence} \\ \midrule \endhead
Persistence & Input/Context (record of the conversation); Control flow (resume) & Every thread is a JSONL rollout at \texttt{\textasciitilde{}/\allowbreak{}.\allowbreak{}codex/\allowbreak{}sessions/\allowbreak{}.\allowbreak{}.\allowbreak{}.\allowbreak{}/\allowbreak{}rollout-\allowbreak{}<ts>-\allowbreak{}<id>.\allowbreak{}jsonl} written by \texttt{RolloutRecorder}; \texttt{should\_\allowbreak{}persist\_\allowbreak{}response\_\allowbreak{}item} filters items; \texttt{reconstruct\_\allowbreak{}rollout} replays from the latest compaction checkpoint on resume; SQLite mirrors metadata; approvals persist as execpolicy amendments; memories pipeline writes \texttt{\$CODEX\_\allowbreak{}HOME/\allowbreak{}memories} & \texttt{rollout/\allowbreak{}src/\allowbreak{}recorder.\allowbreak{}rs:\allowbreak{}L77-\allowbreak{}-\allowbreak{}L88}; \texttt{rollout/\allowbreak{}src/\allowbreak{}policy.\allowbreak{}rs:\allowbreak{}L80-\allowbreak{}-\allowbreak{}L100}; \texttt{core/\allowbreak{}src/\allowbreak{}session/\allowbreak{}rollout\_\allowbreak{}reconstruction.\allowbreak{}rs:\allowbreak{}L208-\allowbreak{}-\allowbreak{}L250}; \texttt{core/\allowbreak{}src/\allowbreak{}thread\_\allowbreak{}manager.\allowbreak{}rs:\allowbreak{}L1319-\allowbreak{}-\allowbreak{}L1332}; \texttt{state/\allowbreak{}src/\allowbreak{}lib.\allowbreak{}rs:\allowbreak{}L1-\allowbreak{}-\allowbreak{}L5}; \texttt{core/\allowbreak{}src/\allowbreak{}session/\allowbreak{}handlers.\allowbreak{}rs:\allowbreak{}L184-\allowbreak{}-\allowbreak{}L198}; \texttt{memories/\allowbreak{}README.\allowbreak{}md} (When it runs) \\
Safety \& security & Output (sandbox, exec policy); Control flow (approvals, Guardian) & OS sandboxes per platform; execpolicy rules \texttt{Allow/\allowbreak{}Prompt/\allowbreak{}Forbidden} plus dangerous-command matcher; patch writable-root check; approval precedence hooks$\rightarrow$Guardian$\rightarrow$user; untrusted projects skip AGENTS.md; sub-agent depth/thread caps; tool-output metadata cleared before sending unless internal & \texttt{sandboxing/\allowbreak{}src/\allowbreak{}manager.\allowbreak{}rs:\allowbreak{}L50-\allowbreak{}-\allowbreak{}L63}; \texttt{execpolicy/\allowbreak{}src/\allowbreak{}decision.\allowbreak{}rs:\allowbreak{}L9-\allowbreak{}-\allowbreak{}L16}; \texttt{core/\allowbreak{}src/\allowbreak{}exec\_\allowbreak{}policy.\allowbreak{}rs:\allowbreak{}L318-\allowbreak{}-\allowbreak{}L400}; \texttt{core/\allowbreak{}src/\allowbreak{}safety.\allowbreak{}rs:\allowbreak{}L67-\allowbreak{}-\allowbreak{}L125}; \texttt{core/\allowbreak{}src/\allowbreak{}tools/\allowbreak{}approvals.\allowbreak{}rs:\allowbreak{}L479-\allowbreak{}-\allowbreak{}L498}; \texttt{core/\allowbreak{}src/\allowbreak{}agents\_\allowbreak{}md.\allowbreak{}rs:\allowbreak{}L64-\allowbreak{}-\allowbreak{}L66}; \texttt{core/\allowbreak{}src/\allowbreak{}agent/\allowbreak{}registry.\allowbreak{}rs:\allowbreak{}L85-\allowbreak{}-\allowbreak{}L99}; \texttt{core/\allowbreak{}src/\allowbreak{}client.\allowbreak{}rs:\allowbreak{}L1072-\allowbreak{}-\allowbreak{}L1076} \\
Observability & across primitives (what is recorded) & OTel spans per sampling (\texttt{codex.\allowbreak{}sampling}, \texttt{handle\_\allowbreak{}responses} with \texttt{gen\_\allowbreak{}ai.\allowbreak{}usage.\allowbreak{}*} fields), tool-decision and tool-result events, approval counters, per-call \texttt{call\_\allowbreak{}trace}, analytics events, token-count events to the client & \texttt{core/\allowbreak{}src/\allowbreak{}session/\allowbreak{}turn.\allowbreak{}rs:\allowbreak{}L2653-\allowbreak{}-\allowbreak{}L2658}, \texttt{L2715-\allowbreak{}-\allowbreak{}L2729}; \texttt{otel/\allowbreak{}src/\allowbreak{}events/\allowbreak{}session\_\allowbreak{}telemetry.\allowbreak{}rs:\allowbreak{}L637-\allowbreak{}-\allowbreak{}L665}, \texttt{L1195-\allowbreak{}-\allowbreak{}L1204}, \texttt{L1297-\allowbreak{}-\allowbreak{}L1308}; \texttt{core/\allowbreak{}src/\allowbreak{}tools/\allowbreak{}sandboxing.\allowbreak{}rs:\allowbreak{}L100-\allowbreak{}-\allowbreak{}L107}; \texttt{core/\allowbreak{}src/\allowbreak{}stream\_\allowbreak{}events\_\allowbreak{}utils.\allowbreak{}rs:\allowbreak{}L326-\allowbreak{}-\allowbreak{}L331} \\
Resilience & Model interface (retries, fallback, timeouts); Context (compaction); Control flow (resume/recovery) & Stream retry budget with backoff and \texttt{Retry-\allowbreak{}After}; WS$\rightarrow$HTTP fallback; 300 s idle timeout; auth-refresh retry; auto compaction pre/mid/post turn and background; \texttt{Op:\allowbreak{}:\allowbreak{}RecoverTurn}/\texttt{ContinueIfIdle} for recovery; delegate failure cleanup & \texttt{core/\allowbreak{}src/\allowbreak{}responses\_\allowbreak{}retry.\allowbreak{}rs:\allowbreak{}L57-\allowbreak{}-\allowbreak{}L176}; \texttt{core/\allowbreak{}src/\allowbreak{}client.\allowbreak{}rs:\allowbreak{}L2363-\allowbreak{}-\allowbreak{}L2435}; \texttt{codex-\allowbreak{}api/\allowbreak{}src/\allowbreak{}sse/\allowbreak{}responses.\allowbreak{}rs:\allowbreak{}L554-\allowbreak{}-\allowbreak{}L558}; \texttt{core/\allowbreak{}src/\allowbreak{}session/\allowbreak{}turn.\allowbreak{}rs:\allowbreak{}L654-\allowbreak{}-\allowbreak{}L716}; \texttt{core/\allowbreak{}src/\allowbreak{}session/\allowbreak{}handlers.\allowbreak{}rs:\allowbreak{}L574-\allowbreak{}-\allowbreak{}L590}; \texttt{core/\allowbreak{}src/\allowbreak{}codex\_\allowbreak{}delegate.\allowbreak{}rs:\allowbreak{}L170-\allowbreak{}-\allowbreak{}L179} \\
Performance & Model interface (streaming, WS prewarm, cache key); Output (parallel tools); Context (truncation, deferred tools, diff-based world state) & WebSocket connection reuse and prewarm; \texttt{prompt\_\allowbreak{}cache\_\allowbreak{}key}; parallel tool execution gate; tool-output truncation; deferred tools kept out of the initial tool list; world-state merge-patch injection; background compaction & \texttt{core/\allowbreak{}src/\allowbreak{}client.\allowbreak{}rs:\allowbreak{}L17-\allowbreak{}-\allowbreak{}L26}, \texttt{L1060}; \texttt{core/\allowbreak{}src/\allowbreak{}tools/\allowbreak{}parallel.\allowbreak{}rs:\allowbreak{}L213-\allowbreak{}-\allowbreak{}L217}; \texttt{core/\allowbreak{}src/\allowbreak{}context\_\allowbreak{}manager/\allowbreak{}history.\allowbreak{}rs:\allowbreak{}L568-\allowbreak{}-\allowbreak{}L577}; \texttt{tools/\allowbreak{}src/\allowbreak{}tool\_\allowbreak{}executor.\allowbreak{}rs:\allowbreak{}L57-\allowbreak{}-\allowbreak{}L62}; \texttt{core/\allowbreak{}src/\allowbreak{}context/\allowbreak{}world\_\allowbreak{}state/\allowbreak{}mod.\allowbreak{}rs:\allowbreak{}L404-\allowbreak{}-\allowbreak{}L425}; \texttt{core/\allowbreak{}src/\allowbreak{}session/\allowbreak{}async\_\allowbreak{}compaction.\allowbreak{}rs:\allowbreak{}L62-\allowbreak{}-\allowbreak{}L95} \\
Evaluation & --- (test suites) & \textasciitilde{}249 integration suites under \texttt{core/\allowbreak{}tests/\allowbreak{}suite/\allowbreak{}} driven by a wiremock Responses mock (\texttt{sse}, \texttt{ev\_\allowbreak{}function\_\allowbreak{}call}, \texttt{mount\_\allowbreak{}sse\_\allowbreak{}once\_\allowbreak{}match}), plus unit tests colocated as \texttt{*\_\allowbreak{}tests.\allowbreak{}rs} & \texttt{core/\allowbreak{}tests/\allowbreak{}suite/\allowbreak{}} (listing); \texttt{core/\allowbreak{}tests/\allowbreak{}common/\allowbreak{}responses.\allowbreak{}rs:\allowbreak{}L27-\allowbreak{}-\allowbreak{}L31}, \texttt{L753}, \texttt{L971}, \texttt{L1133} \\
\bottomrule \end{longtable}

\subsubsection*{Design decisions (one line each: the choice this harness made within a primitive)}
\begin{itemize}
\item Input: a single \texttt{Op} channel per thread; user text during a running turn is queued (steer) rather than starting a new turn, and the model's questions come back as a typed \texttt{UserInputAnswer} op.
\item Context: history is the Responses item list itself; environment/instructions are typed "world-state sections" injected as diffs, tool outputs are token-truncated in live history only, and compaction is pluggable (local summary, remote v2, or budget reset).
\item Model interface: Responses API only (\texttt{WireApi:\allowbreak{}:\allowbreak{}Responses}), WebSocket-first with session-wide HTTP fallback, server-advised retry timing, \texttt{store:\allowbreak{} false} with encrypted reasoning replay.
\item Output: every risky action goes through one orchestrator (approval $\rightarrow$ platform sandbox $\rightarrow$ run $\rightarrow$ sandbox-denied retry); tools run concurrently under a read/write gate; a V8 "code mode" lets the model batch tool calls in JS.
\item Control flow: one task per turn, inner loop until the model emits no tool call (or \texttt{end\_\allowbreak{}turn=false}), with compaction roll-over, stop hooks, and sub-agents that are full \texttt{Session}s coordinated by \texttt{AgentControl} with depth/thread caps.
\end{itemize}

\subsubsection*{Hard cases}
\begin{itemize}
\item \textbf{Guardian (auto-review of approvals):} a harness-consumed gate (control flow) whose decision is produced by a \emph{model} in a separate \texttt{Session} (\texttt{ThreadSource:\allowbreak{}:\allowbreak{}GuardianReview}). Rule: "an approval the harness consumes is control flow"; the reviewer itself is an inner harness (multi-agent rule). Placed as control flow, with the reviewer's own loop as a nested harness. \texttt{core/\allowbreak{}src/\allowbreak{}guardian/\allowbreak{}mod.\allowbreak{}rs:\allowbreak{}L1-\allowbreak{}-\allowbreak{}L3}; \texttt{core/\allowbreak{}src/\allowbreak{}codex\_\allowbreak{}delegate.\allowbreak{}rs:\allowbreak{}L133-\allowbreak{}-\allowbreak{}L137}
\item \textbf{Content-filter retry that injects guidance:} retrying is model interface, but it also records a user fragment into history (context). Split per rule "assign by function". \texttt{core/\allowbreak{}src/\allowbreak{}responses\_\allowbreak{}retry.\allowbreak{}rs:\allowbreak{}L67-\allowbreak{}-\allowbreak{}L82}
\item \textbf{Incremental tools over WebSocket:} \texttt{build\_\allowbreak{}prompt} deliberately sends empty tools/instructions when the server already holds them (\texttt{previous\_\allowbreak{}response\_\allowbreak{}id}). What the model is shown is unchanged, but the request contents depend on transport state --- placed as model interface with a context dependency. \texttt{core/\allowbreak{}src/\allowbreak{}session/\allowbreak{}turn.\allowbreak{}rs:\allowbreak{}L1675-\allowbreak{}-\allowbreak{}L1689}; \texttt{core/\allowbreak{}src/\allowbreak{}client.\allowbreak{}rs:\allowbreak{}L2137-\allowbreak{}-\allowbreak{}L2160}
\item \textbf{Remote compaction v2:} the \emph{server} summarizes (\texttt{CompactionImplementation:\allowbreak{}:\allowbreak{}ResponsesCompactionV2}); the harness decides when (context) but the summarizer is the model/provider. Placed as context (the fitting step) with the call as model interface. \texttt{core/\allowbreak{}src/\allowbreak{}compact\_\allowbreak{}remote\_\allowbreak{}v2.\allowbreak{}rs:\allowbreak{}L112-\allowbreak{}-\allowbreak{}L140}
\item \textbf{Code mode:} the model-written program runs in V8 and sequences only tools (output per rule), yet nested calls re-enter the harness's approval/sandbox path, so each nested call is a normal output-side tool call. \texttt{core/\allowbreak{}src/\allowbreak{}tools/\allowbreak{}code\_\allowbreak{}mode/\allowbreak{}delegate.\allowbreak{}rs:\allowbreak{}L317-\allowbreak{}-\allowbreak{}L360}
\item \textbf{\texttt{new\_\allowbreak{}context\_\allowbreak{}window} / \texttt{get\_\allowbreak{}context\_\allowbreak{}remaining} tools:} tools whose only effect is on context (reset request) or that expose context state to the model; the call is output, the effect is context/control flow per the "effect" rule. \texttt{core/\allowbreak{}src/\allowbreak{}tools/\allowbreak{}handlers/\allowbreak{}new\_\allowbreak{}context\_\allowbreak{}window.\allowbreak{}rs:\allowbreak{}L44-\allowbreak{}-\allowbreak{}L56}
\item \textbf{World-state diffing:} context mechanism that also writes a \texttt{RolloutItem:\allowbreak{}:\allowbreak{}WorldState} for persistence; the model-facing part is context, the file write is persistence. \texttt{core/\allowbreak{}src/\allowbreak{}session/\allowbreak{}mod.\allowbreak{}rs:\allowbreak{}L3762-\allowbreak{}-\allowbreak{}L3793}
\item \textbf{Steer vs. mailbox vs. hooks as inputs:} three different channels (user steer queue, inter-agent mailbox, async hook results) all merge into \texttt{pending\_\allowbreak{}input} before the next request; all are input, drained by control flow. \texttt{core/\allowbreak{}src/\allowbreak{}session/\allowbreak{}input\_\allowbreak{}queue.\allowbreak{}rs:\allowbreak{}L400-\allowbreak{}-\allowbreak{}L432}
\item \textbf{DOES NOT FIT (partially): plan mode streaming parser.} \texttt{PlanModeStreamState} parses the assistant's streamed text into a "proposed plan" item and can suppress/redirect deltas --- it is output rendering that also changes what is recorded as the last agent message (feeds control flow's stop decision). It straddles output and control flow by design. \texttt{core/\allowbreak{}src/\allowbreak{}session/\allowbreak{}turn.\allowbreak{}rs:\allowbreak{}L1992}, \texttt{L2709}, \texttt{L2846-\allowbreak{}-\allowbreak{}L2858}
\end{itemize}

\subsubsection*{Summary row}
\begin{longtable}{@{}p{0.18\linewidth}p{0.18\linewidth}p{0.18\linewidth}p{0.18\linewidth}p{0.18\linewidth}@{}}
\toprule \textbf{Input} & \textbf{Context} & \textbf{Model interface} & \textbf{Output} & \textbf{Control flow} \\ \midrule \endhead
Op channel; steer queue; mailbox; UserInputAnswer & Responses items + world-state diffs; truncation; 3 compaction modes & Responses API, WS$\rightarrow$HTTP fallback, 5 retries w/ backoff & Sandboxed exec/patch/MCP; parallel gate; V8 code mode & Task per turn; loop until no tool call; approvals; sub-agent Sessions \\
\bottomrule \end{longtable}


\subsection{OpenCode}

\textbf{Source:} https://github.com/anomalyco/opencode $\cdot$ \textbf{Version:} commit 654b79b27e79da735d6b07d766c8d7431a71d1b7 (2026-10-15); \texttt{@opencode/\allowbreak{}core} package version 2.0.29 $\cdot$ \textbf{Language:} TypeScript (Bun runtime, Effect 4 library)
\textbf{Path conventions:} \texttt{core/\allowbreak{}} = \texttt{packages/\allowbreak{}core/\allowbreak{}src/\allowbreak{}}, \texttt{ai/\allowbreak{}} = \texttt{packages/\allowbreak{}ai/\allowbreak{}src/\allowbreak{}}, \texttt{schema/\allowbreak{}} = \texttt{packages/\allowbreak{}schema/\allowbreak{}src/\allowbreak{}}, \texttt{codemode/\allowbreak{}} = \texttt{packages/\allowbreak{}codemode/\allowbreak{}src/\allowbreak{}}, \texttt{server/\allowbreak{}} = \texttt{packages/\allowbreak{}server/\allowbreak{}src/\allowbreak{}}, \texttt{recorder/\allowbreak{}} = \texttt{packages/\allowbreak{}http-\allowbreak{}recorder/\allowbreak{}src/\allowbreak{}}. Line ranges are \texttt{Lstart-\allowbreak{}-\allowbreak{}Lend} at the pinned commit. Note: the brief's expected layout (\texttt{packages/\allowbreak{}opencode/\allowbreak{}src/\allowbreak{}session/\allowbreak{}prompt.\allowbreak{}ts}, \texttt{processor.\allowbreak{}ts}, \texttt{llm.\allowbreak{}ts}, \texttt{system.\allowbreak{}ts}) no longer exists at this commit; the equivalents are \texttt{core/\allowbreak{}session/\allowbreak{}prompt.\allowbreak{}ts} (input preparation), \texttt{core/\allowbreak{}session/\allowbreak{}runner/\allowbreak{}llm.\allowbreak{}ts} (the loop), \texttt{core/\allowbreak{}session/\allowbreak{}runner/\allowbreak{}step.\allowbreak{}ts} (one model call plus its tools), \texttt{core/\allowbreak{}session/\allowbreak{}model-\allowbreak{}request.\allowbreak{}ts} (request assembly), \texttt{core/\allowbreak{}session/\allowbreak{}compaction.\allowbreak{}ts}, \texttt{core/\allowbreak{}session/\allowbreak{}system-\allowbreak{}prompt.\allowbreak{}ts}.
\textbf{What it is:} OpenCode is an open-source coding agent harness by Anomaly (anomalyco) that runs as a local server with TUI, desktop, and web clients; a session is a durable, event-sourced record in SQLite that a per-session "runner" fiber drains step by step. The model is reached through the in-house \texttt{@opencode/\allowbreak{}ai} client, which compiles a canonical \texttt{LLMRequest} into a provider protocol (Anthropic Messages, OpenAI Responses, Bedrock Converse, Gemini, \ldots{}) over HTTP/SSE or a stateful WebSocket, with a Vercel AI SDK bridge for third-party provider packages.

\subsubsection*{Control flow}
\begin{itemize}
\item \textbf{Entry: \texttt{Session.\allowbreak{}prompt} admits input then rings the doorbell:} reconciles the inbox row, prepares the prompt, admits it, and calls \texttt{execution.\allowbreak{}wake(\allowbreak{}sessionID)} unless \texttt{resume === false}. \texttt{core/\allowbreak{}session/\allowbreak{}session.\allowbreak{}ts:\allowbreak{}L145-\allowbreak{}-\allowbreak{}L178}
\item \textbf{Run coordinator (one execution fiber per session):} \texttt{loop} drains, then re-drains while a \texttt{pendingWake} doorbell was rung mid-drain; \texttt{wake} coalesces scopes ("input" subsumes "steer"); \texttt{interrupt} marks \texttt{stopping}, claims recorded wakes and interrupts the owner fiber. \texttt{core/\allowbreak{}session/\allowbreak{}run-\allowbreak{}coordinator.\allowbreak{}ts:\allowbreak{}L74-\allowbreak{}-\allowbreak{}L85}, \texttt{L136-\allowbreak{}-\allowbreak{}L145}, \texttt{L147-\allowbreak{}-\allowbreak{}L167}
\item \textbf{Execution service routes drains and handles \texttt{DrainResult}:} \texttt{Complete} ends, \texttt{Moved} and \texttt{Reloaded} re-enter \texttt{drain} with a continuation step. \texttt{core/\allowbreak{}session/\allowbreak{}execution.\allowbreak{}ts:\allowbreak{}L87-\allowbreak{}-\allowbreak{}L110}; \texttt{DrainResult} enum \texttt{core/\allowbreak{}session/\allowbreak{}runner/\allowbreak{}index.\allowbreak{}ts:\allowbreak{}L23-\allowbreak{}-\allowbreak{}L28}
\item \textbf{The outer loop (\texttt{SessionRunner.\allowbreak{}drain}):} \texttt{while (\allowbreak{}true) \{ next = advanceToStep(\allowbreak{}); if not Ready return; continuing = runStep(\allowbreak{}ctx, step); step++ \}}. \texttt{core/\allowbreak{}session/\allowbreak{}runner/\allowbreak{}llm.\allowbreak{}ts:\allowbreak{}L191-\allowbreak{}-\allowbreak{}L198}
\item \textbf{Step boundary (\texttt{advanceToStep}):} at each boundary the runner checks \texttt{lifecycle.\allowbreak{}isClosed(\allowbreak{})}, pulls the next inbox item under a per-session lock, handles \texttt{compaction} and \texttt{move} controls inline, and otherwise promotes pending steers/queued input before loading context; a promoted input resets \texttt{step = 1}. \texttt{core/\allowbreak{}session/\allowbreak{}runner/\allowbreak{}llm.\allowbreak{}ts:\allowbreak{}L74-\allowbreak{}-\allowbreak{}L189}
\item \textbf{Steer vs queue delivery:} \texttt{Promotable} doc: "steer" promotes steers only (mid-work), "input" also admits one queued item at an idle boundary. \texttt{core/\allowbreak{}session/\allowbreak{}inbox.\allowbreak{}ts:\allowbreak{}L43-\allowbreak{}-\allowbreak{}L48}; \texttt{nextPromotable} \texttt{L407-\allowbreak{}-\allowbreak{}L424}; \texttt{promote} \texttt{L496-\allowbreak{}-\allowbreak{}L531}; a pending compaction is pulled ahead of earlier steers \texttt{L533-\allowbreak{}-\allowbreak{}L550}; \texttt{Delivery = "steer" | "queue"} \texttt{schema/\allowbreak{}session-\allowbreak{}inbox.\allowbreak{}ts:\allowbreak{}L11}
\item \textbf{Inner step policy (\texttt{runStep}):} before every attempt it runs auto-compaction (\texttt{compacted.\allowbreak{}status === "completed"} $\rightarrow$ new assistant id and \texttt{continue}), then prepares the request and dispatches on the attempt \texttt{Outcome}. \texttt{core/\allowbreak{}session/\allowbreak{}runner/\allowbreak{}llm.\allowbreak{}ts:\allowbreak{}L209-\allowbreak{}-\allowbreak{}L330}
\item \textbf{Stop condition:} \texttt{needsContinuation = toolChoice !== "none" \&\& record.\allowbreak{}needsContinuation} (a step with executed tool calls continues; a text-only finish hands back); \texttt{finish === "length"} $\rightarrow$ \texttt{OutputLimit}. \texttt{core/\allowbreak{}session/\allowbreak{}runner/\allowbreak{}step.\allowbreak{}ts:\allowbreak{}L281-\allowbreak{}-\allowbreak{}L283}
\item \textbf{Step budget:} \texttt{stepLimitReached = agent.\allowbreak{}info.\allowbreak{}steps !== undefined \&\& step >= agent.\allowbreak{}info.\allowbreak{}steps}; on the final step the runner appends \texttt{MAX\_\allowbreak{}STEPS\_\allowbreak{}PROMPT} as an assistant message and sets \texttt{toolChoice:\allowbreak{} "none"} while keeping tool definitions to preserve the cached prefix. \texttt{core/\allowbreak{}session/\allowbreak{}runner/\allowbreak{}llm.\allowbreak{}ts:\allowbreak{}L226-\allowbreak{}-\allowbreak{}L247}; prompt text \texttt{core/\allowbreak{}session/\allowbreak{}runner/\allowbreak{}max-\allowbreak{}steps.\allowbreak{}ts:\allowbreak{}L1-\allowbreak{}-\allowbreak{}L16}; \texttt{steps} field \texttt{schema/\allowbreak{}agent.\allowbreak{}ts:\allowbreak{}L33}; tool calls on that step fail with "Tools are disabled after the maximum agent steps" \texttt{core/\allowbreak{}session/\allowbreak{}runner/\allowbreak{}step.\allowbreak{}ts:\allowbreak{}L91-\allowbreak{}-\allowbreak{}L101}
\item \textbf{Output-limit continuation:} after \texttt{finish === "length"} the runner reads the last 3 user/assistant rows; if all three are assistant \texttt{length} finishes it fails with \texttt{output-\allowbreak{}limit}, else publishes the synthetic \texttt{CONTINUE\_\allowbreak{}AFTER\_\allowbreak{}OUTPUT\_\allowbreak{}LIMIT} and goes again. \texttt{core/\allowbreak{}session/\allowbreak{}runner/\allowbreak{}llm.\allowbreak{}ts:\allowbreak{}L290-\allowbreak{}-\allowbreak{}L319}; constant \texttt{L38-\allowbreak{}-\allowbreak{}L39}
\item \textbf{Interrupted-stream continuation:} \texttt{Continue} outcome waits the retry delay, publishes \texttt{CONTINUE\_\allowbreak{}AFTER\_\allowbreak{}INCOMPLETE\_\allowbreak{}STREAM} and starts a fresh assistant message. \texttt{core/\allowbreak{}session/\allowbreak{}runner/\allowbreak{}llm.\allowbreak{}ts:\allowbreak{}L281-\allowbreak{}-\allowbreak{}L289}; constant \texttt{L35-\allowbreak{}-\allowbreak{}L36}
\item \textbf{Retry wait is part of the loop:} \texttt{Retry} outcome $\rightarrow$ \texttt{retry.\allowbreak{}wait} (publishes \texttt{RetryScheduled}, sleeps) then the \texttt{while (\allowbreak{}true)} re-attempts with refreshed context. \texttt{core/\allowbreak{}session/\allowbreak{}runner/\allowbreak{}llm.\allowbreak{}ts:\allowbreak{}L275-\allowbreak{}-\allowbreak{}L280}, \texttt{L216-\allowbreak{}-\allowbreak{}L219}; \texttt{wait} \texttt{core/\allowbreak{}session/\allowbreak{}runner/\allowbreak{}retry.\allowbreak{}ts:\allowbreak{}L97-\allowbreak{}-\allowbreak{}L113}
\item \textbf{Overflow $\rightarrow$ compact $\rightarrow$ go again:} \texttt{Compacted} outcome sets \texttt{recoverOverflow = false} (one overflow compaction per step) and re-attempts. \texttt{core/\allowbreak{}session/\allowbreak{}runner/\allowbreak{}llm.\allowbreak{}ts:\allowbreak{}L265-\allowbreak{}-\allowbreak{}L271}, \texttt{L320-\allowbreak{}-\allowbreak{}L323}
\item \textbf{Approval gate consumed by the harness:} \texttt{Permission.\allowbreak{}assert} evaluates rules (\texttt{deny} $\rightarrow$ \texttt{BlockedError}, \texttt{allow} $\rightarrow$ continue, \texttt{ask} $\rightarrow$ create request and await a \texttt{Deferred}). \texttt{core/\allowbreak{}permission.\allowbreak{}ts:\allowbreak{}L231-\allowbreak{}-\allowbreak{}L263}; the reply endpoint resolves it, \texttt{"always"} persists a rule, a \texttt{reject} with no message fails every pending ask in the session. \texttt{core/\allowbreak{}permission.\allowbreak{}ts:\allowbreak{}L265-\allowbreak{}-\allowbreak{}L325}
\item \textbf{A decline ends the step:} declined tool calls are failed with "The user declined this tool call" and the attempt is interrupted (\texttt{Effect.\allowbreak{}interrupt}) rather than returned to the model. \texttt{core/\allowbreak{}session/\allowbreak{}runner/\allowbreak{}step.\allowbreak{}ts:\allowbreak{}L207-\allowbreak{}-\allowbreak{}L215}, \texttt{L274-\allowbreak{}-\allowbreak{}L277}
\item \textbf{Interrupt with resume:} after interrupting, if the next promotable item is a steer or a control, the coordinator is woken with scope "steer" so queued next-turn prompts stay parked. \texttt{core/\allowbreak{}session/\allowbreak{}execution.\allowbreak{}ts:\allowbreak{}L153-\allowbreak{}-\allowbreak{}L170}
\item \textbf{Sub-agent delegation (outer-harness side):} the \texttt{subagent} tool walks \texttt{parentID} to compute depth and refuses beyond \texttt{experimental.\allowbreak{}subagent\_\allowbreak{}depth} (default 1); it creates a child session with \texttt{parentID}, prompts it, and in foreground blocks on the child's job; background returns immediately. \texttt{core/\allowbreak{}tool/\allowbreak{}plugin/\allowbreak{}subagent.\allowbreak{}ts:\allowbreak{}L117-\allowbreak{}-\allowbreak{}L133}, \texttt{L185-\allowbreak{}-\allowbreak{}L198}, \texttt{L205-\allowbreak{}-\allowbreak{}L218}, \texttt{L229-\allowbreak{}-\allowbreak{}L240}
\item \textbf{Child session is a harness of its own:} the job's \texttt{run} is \texttt{sessions.\allowbreak{}resume(\allowbreak{}childSessionID)} and the result is the text of the last completed assistant message. \texttt{core/\allowbreak{}session/\allowbreak{}subagent-\allowbreak{}job.\allowbreak{}ts:\allowbreak{}L45-\allowbreak{}-\allowbreak{}L63}
\item \textbf{How results return:} foreground $\rightarrow$ tool result text wrapped in \texttt{<subagent sessionID=\ldots{} state="completed">}; background $\rightarrow$ a synthetic user message delivered to the parent when the job settles. \texttt{core/\allowbreak{}tool/\allowbreak{}plugin/\allowbreak{}subagent.\allowbreak{}ts:\allowbreak{}L252-\allowbreak{}-\allowbreak{}L266}; \texttt{core/\allowbreak{}session/\allowbreak{}subagent-\allowbreak{}completion.\allowbreak{}ts:\allowbreak{}L21-\allowbreak{}-\allowbreak{}L46}
\item \textbf{Routing to a cheap model keyed to the work (title):} the title agent uses \texttt{model.\allowbreak{}small(\allowbreak{}primary.\allowbreak{}providerID)} plus a minimal-reasoning variant (\texttt{none|minimal|low}), forked so it never delays the step. \texttt{core/\allowbreak{}session/\allowbreak{}context.\allowbreak{}ts:\allowbreak{}L95-\allowbreak{}-\allowbreak{}L119}, \texttt{L186}; fork \texttt{core/\allowbreak{}session/\allowbreak{}runner/\allowbreak{}llm.\allowbreak{}ts:\allowbreak{}L177-\allowbreak{}-\allowbreak{}L180}
\item \textbf{Compaction model keyed to the work:} the \texttt{compaction} agent may carry its own model; otherwise the session model is used. \texttt{core/\allowbreak{}session/\allowbreak{}compaction.\allowbreak{}ts:\allowbreak{}L213-\allowbreak{}-\allowbreak{}L219}
\item \textbf{Agent profiles (bundles):} \texttt{build} (default, all-allow + question), \texttt{general}/\texttt{explore} sub-agents with deny lists and a custom system prompt, hidden \texttt{compaction}/\texttt{title}/\texttt{summary} agents. \texttt{core/\allowbreak{}plugin/\allowbreak{}agent.\allowbreak{}ts:\allowbreak{}L80-\allowbreak{}-\allowbreak{}L140}; default ruleset \texttt{schema/\allowbreak{}agent.\allowbreak{}ts:\allowbreak{}L46-\allowbreak{}-\allowbreak{}L51}
\item \textbf{Plan mode (bundle):} a \texttt{plan} agent denies \texttt{edit} outside \texttt{\textasciitilde{}/\allowbreak{}.\allowbreak{}opencode/\allowbreak{}plan/\allowbreak{}*} (control flow), rewrites blocked-edit errors (output), and re-injects a \texttt{<system-\allowbreak{}reminder>} per request (context). \texttt{core/\allowbreak{}plugin/\allowbreak{}plan.\allowbreak{}ts:\allowbreak{}L33-\allowbreak{}-\allowbreak{}L42}, \texttt{L44-\allowbreak{}-\allowbreak{}L53}, \texttt{L57-\allowbreak{}-\allowbreak{}L60}
\item \textbf{Crash/restart resume:} at boot, sessions whose execution claim survived are resumed with a synthetic "The server restarted\ldots{}" notice; a durable per-turn counter stops after \texttt{DEFAULT\_\allowbreak{}MAX\_\allowbreak{}ATTEMPTS = 10}. \texttt{core/\allowbreak{}session/\allowbreak{}execution/\allowbreak{}restart.\allowbreak{}ts:\allowbreak{}L77-\allowbreak{}-\allowbreak{}L99}, \texttt{L193-\allowbreak{}-\allowbreak{}L231}; budget \texttt{L33}
\end{itemize}

\subsubsection*{Input}
\begin{itemize}
\item \textbf{HTTP handler $\rightarrow$ \texttt{session.\allowbreak{}prompt}:} \texttt{text}, \texttt{files}, \texttt{agents}, \texttt{skills}, \texttt{metadata}, \texttt{delivery}, \texttt{resume} are passed straight through. \texttt{server/\allowbreak{}handlers/\allowbreak{}session.\allowbreak{}ts:\allowbreak{}L307-\allowbreak{}-\allowbreak{}L322}
\item \textbf{Prompt preparation (\texttt{SessionPrompt.\allowbreak{}prepare}):} fires the plugin hook \texttt{session.\allowbreak{}prompt} (plugins may rewrite text/files/delivery), materializes attachments with concurrency 8, resolves \texttt{skills} to their text, and builds a \texttt{UserPayload}. \texttt{core/\allowbreak{}session/\allowbreak{}prompt.\allowbreak{}ts:\allowbreak{}L29-\allowbreak{}-\allowbreak{}L91}
\item \textbf{Attachments:} \texttt{data:\allowbreak{}} URLs or \texttt{file:\allowbreak{}} URIs; 20 MiB cap; \texttt{?start=\&end=} line ranges for text; a directory becomes a sorted listing. \texttt{core/\allowbreak{}session/\allowbreak{}prompt.\allowbreak{}ts:\allowbreak{}L93-\allowbreak{}-\allowbreak{}L133}, \texttt{L151-\allowbreak{}-\allowbreak{}L201}, \texttt{L203}
\item \textbf{Image normalization on input:} images are resized via \texttt{Image.\allowbreak{}Service} (resizer unavailable $\rightarrow$ passthrough). \texttt{core/\allowbreak{}session/\allowbreak{}prompt.\allowbreak{}ts:\allowbreak{}L135-\allowbreak{}-\allowbreak{}L149}
\item \textbf{Durable admission:} the inbox row is inserted into SQLite (\texttt{admit}, \texttt{reconcile}, \texttt{cancel}, \texttt{steer}, \texttt{queue}). \texttt{core/\allowbreak{}session/\allowbreak{}inbox.\allowbreak{}ts:\allowbreak{}L155}, \texttt{L169}, \texttt{L230-\allowbreak{}-\allowbreak{}L256}, \texttt{L295}
\item \textbf{Harness-originated input (synthetic):} \texttt{Session.\allowbreak{}synthetic} admits text the harness produced (restart notice, job completions) through the same inbox. \texttt{core/\allowbreak{}session/\allowbreak{}session.\allowbreak{}ts:\allowbreak{}L273-\allowbreak{}-\allowbreak{}L314}
\item \textbf{Answer the MODEL consumes (question tool):} the form answers are rendered as \texttt{User has answered your questions:\allowbreak{} \ldots{}} and returned as the tool result. \texttt{core/\allowbreak{}tool/\allowbreak{}plugin/\allowbreak{}question.\allowbreak{}ts:\allowbreak{}L38-\allowbreak{}-\allowbreak{}L46}, \texttt{L62-\allowbreak{}-\allowbreak{}L112}
\item \textbf{Background shell completion event:} \texttt{notifyWhenDone} waits on the job and admits a synthetic message with the output. \texttt{core/\allowbreak{}tool/\allowbreak{}plugin/\allowbreak{}shell.\allowbreak{}ts:\allowbreak{}L157-\allowbreak{}-\allowbreak{}L189}
\item \textbf{Sub-agent completion event:} delivered as a synthetic message tagged \texttt{<subagent \ldots{} state=\ldots{}>}. \texttt{core/\allowbreak{}session/\allowbreak{}subagent-\allowbreak{}completion.\allowbreak{}ts:\allowbreak{}L37-\allowbreak{}-\allowbreak{}L44}
\end{itemize}

\subsubsection*{Context}
\begin{itemize}
\item \textbf{Selection (\texttt{SessionContext.\allowbreak{}select}):} fixes session, agent, tool snapshot (narrowed by merged agent+session permissions) and instruction sources, loaded in parallel. \texttt{core/\allowbreak{}session/\allowbreak{}context.\allowbreak{}ts:\allowbreak{}L121-\allowbreak{}-\allowbreak{}L161}; permission narrowing \texttt{L131}
\item \textbf{Instruction ordering for cache reuse:} "Ordered from most to least shared across sessions so the baseline stays a reusable prompt-cache prefix": Code Mode catalog, MCP, references, skills, AGENTS.md discovery, builtins (date/env), per-session entries. \texttt{core/\allowbreak{}session/\allowbreak{}context.\allowbreak{}ts:\allowbreak{}L147-\allowbreak{}-\allowbreak{}L158}
\item \textbf{Load (\texttt{SessionContext.\allowbreak{}load}):} resolves the model and reads history after the latest replayable compaction. \texttt{core/\allowbreak{}session/\allowbreak{}context.\allowbreak{}ts:\allowbreak{}L163-\allowbreak{}-\allowbreak{}L179}
\item \textbf{System prompt assembly:} \texttt{[agent.\allowbreak{}system ?? default prompt rendered with tool guidance, initial instructions]}; messages via \texttt{toLLMMessages}. \texttt{core/\allowbreak{}session/\allowbreak{}model-\allowbreak{}request.\allowbreak{}ts:\allowbreak{}L100-\allowbreak{}-\allowbreak{}L120}
\item \textbf{Default system prompt text:} 15 lines ("You are an AI agent running in OpenCode\ldots{}"), with \texttt{\$\{OPENCODE\_\allowbreak{}TOOL\_\allowbreak{}GUIDANCE\}} filled per available tool. \texttt{core/\allowbreak{}session/\allowbreak{}runner/\allowbreak{}prompt/\allowbreak{}system.\allowbreak{}txt:\allowbreak{}L1-\allowbreak{}-\allowbreak{}L15}; \texttt{core/\allowbreak{}session/\allowbreak{}system-\allowbreak{}prompt.\allowbreak{}ts:\allowbreak{}L9-\allowbreak{}-\allowbreak{}L28}
\item \textbf{Built-in instructions:} today's date and an \texttt{<env>} block (working dir, workspace root, git, platform, tmp). \texttt{core/\allowbreak{}instructions/\allowbreak{}builtins.\allowbreak{}ts:\allowbreak{}L24-\allowbreak{}-\allowbreak{}L53}
\item \textbf{AGENTS.md discovery:} global \texttt{AGENTS.\allowbreak{}md} plus ancestor-walk from the location to the project root, watched for changes. \texttt{core/\allowbreak{}config/\allowbreak{}plugin/\allowbreak{}instruction.\allowbreak{}ts:\allowbreak{}L36-\allowbreak{}-\allowbreak{}L62}, \texttt{L76-\allowbreak{}-\allowbreak{}L86}
\item \textbf{Catalogs rendered into the baseline:} skills \texttt{<available\_\allowbreak{}skills>} \texttt{core/\allowbreak{}skill/\allowbreak{}instructions.\allowbreak{}ts:\allowbreak{}L16-\allowbreak{}-\allowbreak{}L33}; MCP \texttt{<mcp\_\allowbreak{}instructions>} \texttt{core/\allowbreak{}mcp/\allowbreak{}instructions.\allowbreak{}ts:\allowbreak{}L17-\allowbreak{}-\allowbreak{}L29}; Code Mode catalog \texttt{core/\allowbreak{}codemode/\allowbreak{}instructions.\allowbreak{}ts:\allowbreak{}L9-\allowbreak{}-\allowbreak{}L30}
\item \textbf{Instruction epochs:} changed instruction sources are frozen into a durable \texttt{InstructionsUpdated} event with rendered delta text, so later requests see "Today's date is now\ldots{}" style updates as system messages. \texttt{core/\allowbreak{}session/\allowbreak{}instruction-\allowbreak{}state.\allowbreak{}ts:\allowbreak{}L38-\allowbreak{}-\allowbreak{}L61}, \texttt{L80-\allowbreak{}-\allowbreak{}L87}; rendered as \texttt{Message.\allowbreak{}system} \texttt{core/\allowbreak{}session/\allowbreak{}runner/\allowbreak{}to-\allowbreak{}llm-\allowbreak{}message.\allowbreak{}ts:\allowbreak{}L291-\allowbreak{}-\allowbreak{}L292}
\item \textbf{History read bounded by compaction:} rows from the latest completed compaction \texttt{seq} onward; native checkpoints only when the target model can replay them. \texttt{core/\allowbreak{}session/\allowbreak{}history.\allowbreak{}ts:\allowbreak{}L21-\allowbreak{}-\allowbreak{}L26}, \texttt{L78-\allowbreak{}-\allowbreak{}L108}, \texttt{L118-\allowbreak{}-\allowbreak{}L135}
\item \textbf{Projection of history into model messages:} synthetic/skill/shell messages become \texttt{user} role; a compaction becomes a \texttt{<conversation-\allowbreak{}checkpoint>} user message. \texttt{core/\allowbreak{}session/\allowbreak{}runner/\allowbreak{}to-\allowbreak{}llm-\allowbreak{}message.\allowbreak{}ts:\allowbreak{}L269-\allowbreak{}-\allowbreak{}L327}, \texttt{L330-\allowbreak{}-\allowbreak{}L334}
\item \textbf{Tool definitions snapshot:} tools wholly denied by permission are dropped; \texttt{codemode:\allowbreak{} false} tools are advertised directly, the rest only inside the single \texttt{execute} tool; definitions sorted by name. \texttt{core/\allowbreak{}tool.\allowbreak{}ts:\allowbreak{}L225-\allowbreak{}-\allowbreak{}L287}; JSON Schema generation \texttt{core/\allowbreak{}tool/\allowbreak{}runtime.\allowbreak{}ts:\allowbreak{}L21-\allowbreak{}-\allowbreak{}L27}, \texttt{L174-\allowbreak{}-\allowbreak{}L188}
\item \textbf{Plugin hook that rewrites what the model sees:} \texttt{session.\allowbreak{}context} sees system, messages, tool descriptions/schemas and the default \texttt{maxTokens}; renamed tools are mapped back by object identity. \texttt{core/\allowbreak{}session/\allowbreak{}model-\allowbreak{}request.\allowbreak{}ts:\allowbreak{}L236-\allowbreak{}-\allowbreak{}L274}
\item \textbf{Media fitting:} unsupported modalities are replaced by an ERROR text; inline images are dropped from oldest once the total exceeds 25 MiB until under 15 MiB. \texttt{core/\allowbreak{}session/\allowbreak{}model-\allowbreak{}request.\allowbreak{}ts:\allowbreak{}L43-\allowbreak{}-\allowbreak{}L46}, \texttt{L132-\allowbreak{}-\allowbreak{}L147}, \texttt{L189-\allowbreak{}-\allowbreak{}L216}; applied \texttt{L297}
\item \textbf{Output token budget fitted to the window:} \texttt{outputLimit} = min(catalog limit, 256k), reduced to the room left after \texttt{measured + estimated$\times$1.\allowbreak{}15}, floor 1,024. \texttt{core/\allowbreak{}session/\allowbreak{}model-\allowbreak{}request.\allowbreak{}ts:\allowbreak{}L47-\allowbreak{}-\allowbreak{}L57}, \texttt{L88-\allowbreak{}-\allowbreak{}L98}
\item \textbf{Auto-compaction trigger:} \texttt{due} when \texttt{estimateContext >= ceiling}; ceiling = window - max(10 \%, 16,000) or window - configured \texttt{buffer}; skipped right after a completed compaction or before a measured response. \texttt{core/\allowbreak{}session/\allowbreak{}compaction.\allowbreak{}ts:\allowbreak{}L205-\allowbreak{}-\allowbreak{}L240}, \texttt{L242-\allowbreak{}-\allowbreak{}L251}, \texttt{L911-\allowbreak{}-\allowbreak{}L920}; defaults \texttt{auto:\allowbreak{} true, keep:\allowbreak{} 15\_\allowbreak{}000} \texttt{L197}; settings \texttt{L41-\allowbreak{}-\allowbreak{}L47}
\item \textbf{Prompt size estimate = provider-measured + locally estimated:} anchors on the last assistant message with a usage count from the same provider, estimates everything after at \textasciitilde{}4 chars/token. \texttt{core/\allowbreak{}session/\allowbreak{}compaction.\allowbreak{}ts:\allowbreak{}L880-\allowbreak{}-\allowbreak{}L905}, \texttt{L924-\allowbreak{}-\allowbreak{}L930}, \texttt{L950-\allowbreak{}-\allowbreak{}L953}
\item \textbf{Summary compaction:} splits history into \texttt{older} (summarized) and \texttt{recent} (kept verbatim within \texttt{keep} tokens, starting at a user boundary); sends older + summary prompt; one template nudge before failing. \texttt{core/\allowbreak{}session/\allowbreak{}compaction.\allowbreak{}ts:\allowbreak{}L272-\allowbreak{}-\allowbreak{}L316}, \texttt{L708-\allowbreak{}-\allowbreak{}L744}; template \texttt{L104-\allowbreak{}-\allowbreak{}L135}
\item \textbf{Tool outputs in the verbatim part are cut to 1,250 chars.} \texttt{core/\allowbreak{}session/\allowbreak{}compaction.\allowbreak{}ts:\allowbreak{}L100}, \texttt{L813-\allowbreak{}-\allowbreak{}L826}
\item \textbf{Overflow compaction (provider said too long):} captured before any output starts, then compaction aims at 70 \%/50 \%/35 \% of the rejected request size. \texttt{core/\allowbreak{}session/\allowbreak{}runner/\allowbreak{}step.\allowbreak{}ts:\allowbreak{}L113-\allowbreak{}-\allowbreak{}L120}, \texttt{L155-\allowbreak{}-\allowbreak{}L167}; \texttt{SHRINK\_\allowbreak{}STEPS} \texttt{core/\allowbreak{}session/\allowbreak{}compaction.\allowbreak{}ts:\allowbreak{}L95}, \texttt{L429-\allowbreak{}-\allowbreak{}L448}
\item \textbf{Manual \texttt{/\allowbreak{}compact}:} an inbox control item processed at the step boundary. \texttt{core/\allowbreak{}session/\allowbreak{}session.\allowbreak{}ts:\allowbreak{}L246-\allowbreak{}-\allowbreak{}L264}; \texttt{core/\allowbreak{}session/\allowbreak{}runner/\allowbreak{}llm.\allowbreak{}ts:\allowbreak{}L115-\allowbreak{}-\allowbreak{}L163}
\item \textbf{Native (provider-side) compaction:} for models with \texttt{compaction.\allowbreak{}type === "native"} the provider returns one opaque item that replaces the window; the harness only appends the newest user messages. \texttt{core/\allowbreak{}session/\allowbreak{}compaction.\allowbreak{}ts:\allowbreak{}L318-\allowbreak{}-\allowbreak{}L332}
\item \textbf{Tool result truncation before it enters history:} 2,000 lines / 50 KiB, overflow written to a file with a \texttt{[showing lines 1-\allowbreak{}N of M; full output saved to \ldots{}]} marker. \texttt{core/\allowbreak{}tool-\allowbreak{}output.\allowbreak{}ts:\allowbreak{}L13-\allowbreak{}-\allowbreak{}L16}, \texttt{L67-\allowbreak{}-\allowbreak{}L118}; applied in the step \texttt{core/\allowbreak{}session/\allowbreak{}runner/\allowbreak{}step.\allowbreak{}ts:\allowbreak{}L127}
\item \textbf{Tool input repair (what the model's call is seen as):} stringified JSON, numeric strings, placeholder nulls are repaired against the schema before execution. \texttt{core/\allowbreak{}plugin/\allowbreak{}tool-\allowbreak{}input-\allowbreak{}repair.\allowbreak{}ts:\allowbreak{}L7-\allowbreak{}-\allowbreak{}L36}
\item \textbf{Prompt-cache key / affinity:} \texttt{promptCacheKey} and \texttt{x-\allowbreak{}session-\allowbreak{}affinity} derive from the root session (children and forks share the parent's). \texttt{core/\allowbreak{}session/\allowbreak{}model-\allowbreak{}request.\allowbreak{}ts:\allowbreak{}L278-\allowbreak{}-\allowbreak{}L295}; \texttt{core/\allowbreak{}session/\allowbreak{}affinity.\allowbreak{}ts:\allowbreak{}L6-\allowbreak{}-\allowbreak{}L8}
\item \textbf{Searchable memory for the model:} not found (no model-facing memory store; \texttt{core/\allowbreak{}environment/\allowbreak{}memory.\allowbreak{}ts} is an in-memory filesystem).
\end{itemize}

\subsubsection*{Model interface}
\begin{itemize}
\item \textbf{Streaming request:} the step calls \texttt{llm.\allowbreak{}stream(\allowbreak{}request, options)} and folds every \texttt{LLMEvent} through a publisher; it reads to the end "so the next request can reuse this response". \texttt{core/\allowbreak{}session/\allowbreak{}runner/\allowbreak{}step.\allowbreak{}ts:\allowbreak{}L103-\allowbreak{}-\allowbreak{}L138}; client \texttt{ai/\allowbreak{}route/\allowbreak{}client.\allowbreak{}ts:\allowbreak{}L602-\allowbreak{}-\allowbreak{}L607}, \texttt{L618-\allowbreak{}-\allowbreak{}L622}
\item \textbf{Request construction:} \texttt{LLM.\allowbreak{}request(\allowbreak{}\{model, http.\allowbreak{}headers, promptCacheKey, system, messages, tools, toolChoice, generation, providerOptions\})} with \texttt{x-\allowbreak{}opencode-\allowbreak{}session-\allowbreak{}id}/\texttt{User-\allowbreak{}Agent} headers. \texttt{core/\allowbreak{}session/\allowbreak{}model-\allowbreak{}request.\allowbreak{}ts:\allowbreak{}L279-\allowbreak{}-\allowbreak{}L302}
\item \textbf{Wire format:} per-provider protocol modules (\texttt{anthropic-\allowbreak{}messages.\allowbreak{}ts}, \texttt{openai-\allowbreak{}responses.\allowbreak{}ts}, \texttt{bedrock-\allowbreak{}converse.\allowbreak{}ts}, \texttt{gemini.\allowbreak{}ts}, \ldots{}) and framings (SSE) compile the canonical request and decode the stream. \texttt{ai/\allowbreak{}protocols/\allowbreak{}} (directory), \texttt{ai/\allowbreak{}route/\allowbreak{}client.\allowbreak{}ts:\allowbreak{}L423-\allowbreak{}-\allowbreak{}L497}
\item \textbf{Route/header hooks:} \texttt{session.\allowbreak{}model.\allowbreak{}request} may change \texttt{baseURL}/headers; \texttt{session.\allowbreak{}http.\allowbreak{}request}/\texttt{http.\allowbreak{}response} wrap the raw HTTP exchange. \texttt{core/\allowbreak{}session/\allowbreak{}model-\allowbreak{}request.\allowbreak{}ts:\allowbreak{}L305-\allowbreak{}-\allowbreak{}L320}, \texttt{L337-\allowbreak{}-\allowbreak{}L365}
\item \textbf{HTTP timeouts:} headers 300 s, inter-chunk 900 s, whole-request none by default. \texttt{ai/\allowbreak{}route/\allowbreak{}transport/\allowbreak{}http.\allowbreak{}ts:\allowbreak{}L100-\allowbreak{}-\allowbreak{}L146}; defaults \texttt{ai/\allowbreak{}schema/\allowbreak{}options.\allowbreak{}ts:\allowbreak{}L56}, \texttt{L62}
\item \textbf{Stateful WebSocket transport:} bound per session when \texttt{model.\allowbreak{}transport === "websocket"}; rotates after 55 min, 15 s connect, 30 min idle; after 5 consecutive stream failures the session stays on HTTP. \texttt{core/\allowbreak{}session/\allowbreak{}model-\allowbreak{}request.\allowbreak{}ts:\allowbreak{}L368-\allowbreak{}-\allowbreak{}L392}; \texttt{core/\allowbreak{}session/\allowbreak{}model-\allowbreak{}transport.\allowbreak{}ts:\allowbreak{}L22-\allowbreak{}-\allowbreak{}L27}, \texttt{L186-\allowbreak{}-\allowbreak{}L196}, \texttt{L308}
\item \textbf{Error classification:} \texttt{isContextOverflowFailure} and \texttt{isRetryable} (rate limit, provider internal, transport unless delivery rejected, incomplete stream, unknown; \texttt{x-\allowbreak{}should-\allowbreak{}retry} header overrides). \texttt{ai/\allowbreak{}provider-\allowbreak{}error.\allowbreak{}ts:\allowbreak{}L72-\allowbreak{}-\allowbreak{}L75}, \texttt{L81-\allowbreak{}-\allowbreak{}L110}
\item \textbf{Retry policy:} exponential from 2 s capped at 10 s per gap, 10 retries, jittered; provider \texttt{retry-\allowbreak{}after} honored up to 15 min; at most 3 timeout retries; \texttt{session.\allowbreak{}retry} plugin hook may override the decision. \texttt{core/\allowbreak{}session/\allowbreak{}runner/\allowbreak{}retry.\allowbreak{}ts:\allowbreak{}L31-\allowbreak{}-\allowbreak{}L62}, \texttt{L64-\allowbreak{}-\allowbreak{}L92}
\item \textbf{Retry vs continue vs recover-full:} before output started $\rightarrow$ \texttt{Retry}; after output started on an interrupted stream $\rightarrow$ \texttt{Continue}; WebSocket \texttt{retry-\allowbreak{}full}/\texttt{rotate-\allowbreak{}and-\allowbreak{}retry-\allowbreak{}full} $\rightarrow$ \texttt{RecoverFull}. \texttt{core/\allowbreak{}session/\allowbreak{}runner/\allowbreak{}step.\allowbreak{}ts:\allowbreak{}L182-\allowbreak{}-\allowbreak{}L204}, \texttt{L263-\allowbreak{}-\allowbreak{}L271}, \texttt{L291-\allowbreak{}-\allowbreak{}L295}
\item \textbf{Unknown finish reason is an error:} turned into \texttt{InvalidProviderOutputError(\allowbreak{}"incomplete-\allowbreak{}stream")}. \texttt{core/\allowbreak{}session/\allowbreak{}runner/\allowbreak{}step.\allowbreak{}ts:\allowbreak{}L171-\allowbreak{}-\allowbreak{}L181}
\item \textbf{Model resolution:} session model must be in the available catalog, else \texttt{ModelUnavailableError}. \texttt{core/\allowbreak{}session/\allowbreak{}runner/\allowbreak{}model.\allowbreak{}ts:\allowbreak{}L75-\allowbreak{}-\allowbreak{}L91}
\item \textbf{Token usage and cost:} normalized \texttt{input/\allowbreak{}output/\allowbreak{}reasoning/\allowbreak{}cache.\allowbreak{}read/\allowbreak{}cache.\allowbreak{}write}, cost by context tier. \texttt{core/\allowbreak{}session/\allowbreak{}usage.\allowbreak{}ts:\allowbreak{}L11-\allowbreak{}-\allowbreak{}L19}, \texttt{L22-\allowbreak{}-\allowbreak{}L36}
\item \textbf{Vercel AI SDK bridge:} \texttt{AISDK.\allowbreak{}Service.\allowbreak{}model(\allowbreak{}RuntimeInfo)} wraps a \texttt{LanguageModelV3} as an \texttt{@opencode/\allowbreak{}ai} \texttt{LanguageModel} for third-party provider packages. \texttt{core/\allowbreak{}aisdk.\allowbreak{}ts:\allowbreak{}L237-\allowbreak{}-\allowbreak{}L252}
\item \textbf{Backup model for a request that could not be served (title only):} if the small model's attempt fails, the primary model is tried. \texttt{core/\allowbreak{}session/\allowbreak{}title.\allowbreak{}ts:\allowbreak{}L129-\allowbreak{}-\allowbreak{}L133}
\end{itemize}

\subsubsection*{Output}
\begin{itemize}
\item \textbf{Streaming render path:} provider events are published as durable \texttt{Text.\allowbreak{}Started/\allowbreak{}Delta/\allowbreak{}Ended}, \texttt{Reasoning.\allowbreak{}*}, \texttt{Tool.\allowbreak{}Called/\allowbreak{}Progress/\allowbreak{}Success/\allowbreak{}Failed}, \texttt{Step.\allowbreak{}Started/\allowbreak{}Streamed/\allowbreak{}Ended/\allowbreak{}Failed}; deltas batched every 100 ms. \texttt{core/\allowbreak{}session/\allowbreak{}runner/\allowbreak{}publish-\allowbreak{}llm-\allowbreak{}event.\allowbreak{}ts:\allowbreak{}L443-\allowbreak{}-\allowbreak{}L470}, \texttt{L106}, \texttt{L194-\allowbreak{}-\allowbreak{}L200}
\item \textbf{Concurrent tool execution from one response:} each \texttt{tool-\allowbreak{}call} event forks a fiber immediately; all are awaited after the stream ends. \texttt{core/\allowbreak{}session/\allowbreak{}runner/\allowbreak{}step.\allowbreak{}ts:\allowbreak{}L121-\allowbreak{}-\allowbreak{}L134}, \texttt{L151}
\item \textbf{Tool executor:} \texttt{execute.\allowbreak{}before} hook $\rightarrow$ input decode/validate $\rightarrow$ \texttt{tool.\allowbreak{}execute} $\rightarrow$ \texttt{execute.\allowbreak{}after} hook $\rightarrow$ content normalization (image resize). \texttt{core/\allowbreak{}tool.\allowbreak{}ts:\allowbreak{}L103-\allowbreak{}-\allowbreak{}L156}, \texttt{L263-\allowbreak{}-\allowbreak{}L284}; validation \texttt{core/\allowbreak{}tool/\allowbreak{}runtime.\allowbreak{}ts:\allowbreak{}L29-\allowbreak{}-\allowbreak{}L61}, \texttt{L88-\allowbreak{}-\allowbreak{}L100}
\item \textbf{Shell tool:} 2-minute default timeout, \texttt{AGENT=1 OPENCODE=1} env, jobs with foreground block or background; no OS sandbox found, gating is by permission after tree-sitter parsing of the command into per-command resources and touched directories. \texttt{core/\allowbreak{}tool/\allowbreak{}plugin/\allowbreak{}shell.\allowbreak{}ts:\allowbreak{}L24}, \texttt{L117-\allowbreak{}-\allowbreak{}L155}, \texttt{L199-\allowbreak{}-\allowbreak{}L279}; parser \texttt{core/\allowbreak{}shell/\allowbreak{}parse.\allowbreak{}ts:\allowbreak{}L7}, \texttt{L392}
\item \textbf{Edit tool asks with a diff preview:} \texttt{permission.\allowbreak{}assert(\allowbreak{}\{action:\allowbreak{} "edit", metadata:\allowbreak{} \{files:\allowbreak{} [preview]\}\})} before applying. \texttt{core/\allowbreak{}tool/\allowbreak{}plugin/\allowbreak{}edit.\allowbreak{}ts:\allowbreak{}L271-\allowbreak{}-\allowbreak{}L283}
\item \textbf{Read/external-directory gating:} \texttt{authorizeRead} asks \texttt{external\_\allowbreak{}directory} for paths outside the location, then \texttt{read} per file. \texttt{core/\allowbreak{}file-\allowbreak{}access.\allowbreak{}ts:\allowbreak{}L125-\allowbreak{}-\allowbreak{}L142}, \texttt{L144-\allowbreak{}-\allowbreak{}L161}
\item \textbf{MCP tools as output:} each discovered MCP tool is registered under a server namespace, permission-asserted, then \texttt{callTool}. \texttt{core/\allowbreak{}tool/\allowbreak{}mcp.\allowbreak{}ts:\allowbreak{}L41-\allowbreak{}-\allowbreak{}L77}
\item \textbf{Code Mode \texttt{execute} tool:} the model submits JavaScript; tool calls inside it are routed through the same \texttt{executeTool} (hooks + permissions); \texttt{fetch} only, no imports/fs/timers. \texttt{core/\allowbreak{}codemode/\allowbreak{}tool.\allowbreak{}ts:\allowbreak{}L62-\allowbreak{}-\allowbreak{}L69}, \texttt{L71-\allowbreak{}-\allowbreak{}L108}; sandbox = custom tree-walking interpreter over acorn AST with value limits and optional timeout. \texttt{codemode/\allowbreak{}interpreter/\allowbreak{}execute.\allowbreak{}ts:\allowbreak{}L1}, \texttt{L56-\allowbreak{}-\allowbreak{}L68}; \texttt{codemode/\allowbreak{}interpreter/\allowbreak{}limits.\allowbreak{}ts:\allowbreak{}L6-\allowbreak{}-\allowbreak{}L13}
\item \textbf{File-change artifact for the user:} a git-backed snapshot is captured before and after each step and the changed paths are attached to \texttt{Step.\allowbreak{}Ended}. \texttt{core/\allowbreak{}session/\allowbreak{}runner/\allowbreak{}step.\allowbreak{}ts:\allowbreak{}L74-\allowbreak{}-\allowbreak{}L76}, \texttt{L234-\allowbreak{}-\allowbreak{}L260}; \texttt{core/\allowbreak{}snapshot.\allowbreak{}ts:\allowbreak{}L43-\allowbreak{}-\allowbreak{}L68}, \texttt{L90-\allowbreak{}-\allowbreak{}L105}
\item \textbf{Question tool shows a form:} \texttt{forms.\allowbreak{}ask} with one field per question, "Type your own answer" custom option. \texttt{core/\allowbreak{}tool/\allowbreak{}plugin/\allowbreak{}question.\allowbreak{}ts:\allowbreak{}L62-\allowbreak{}-\allowbreak{}L112}, \texttt{L119-\allowbreak{}-\allowbreak{}L132}
\item \textbf{Tool failure settlement:} interrupted, declined, provider-failed, and result-missing calls are failed with distinct errors so no call stays unsettled. \texttt{core/\allowbreak{}session/\allowbreak{}runner/\allowbreak{}step.\allowbreak{}ts:\allowbreak{}L62-\allowbreak{}-\allowbreak{}L64}, \texttt{L207-\allowbreak{}-\allowbreak{}L232}
\end{itemize}

\subsubsection*{Production concerns}
\begin{longtable}{@{}p{0.13\linewidth}p{0.17\linewidth}p{0.36\linewidth}p{0.26\linewidth}@{}}
\toprule \textbf{Concern} & \textbf{Where it lives (primitive)} & \textbf{Mechanism} & \textbf{Evidence} \\ \midrule \endhead
Persistence & Input, Context, Control flow & SQLite via drizzle: \texttt{session\_\allowbreak{}v2} (model, agent, permissions, token totals, execution claim), \texttt{session\_\allowbreak{}message}, inbox; a durable, per-aggregate sequenced event log projected into those tables; execution claim written at turn start and released only by a terminal event (shutdown keeps it for resume); "always" permission replies saved per project; git snapshots for revert; tool-output spill files kept 7 days & \texttt{core/\allowbreak{}session/\allowbreak{}sql.\allowbreak{}ts:\allowbreak{}L22-\allowbreak{}-\allowbreak{}L77}; \texttt{core/\allowbreak{}bus.\allowbreak{}ts:\allowbreak{}L21-\allowbreak{}-\allowbreak{}L48}; \texttt{core/\allowbreak{}session/\allowbreak{}projector.\allowbreak{}ts:\allowbreak{}L1-\allowbreak{}-\allowbreak{}L35}; \texttt{core/\allowbreak{}session/\allowbreak{}execution.\allowbreak{}ts:\allowbreak{}L76-\allowbreak{}-\allowbreak{}L86}, \texttt{L111-\allowbreak{}-\allowbreak{}L148}; \texttt{core/\allowbreak{}session/\allowbreak{}store.\allowbreak{}ts:\allowbreak{}L62-\allowbreak{}-\allowbreak{}L74}; \texttt{core/\allowbreak{}permission.\allowbreak{}ts:\allowbreak{}L299-\allowbreak{}-\allowbreak{}L305}; \texttt{core/\allowbreak{}session/\allowbreak{}revert.\allowbreak{}ts:\allowbreak{}L21-\allowbreak{}-\allowbreak{}L40}; \texttt{core/\allowbreak{}tool-\allowbreak{}output.\allowbreak{}ts:\allowbreak{}L15} \\
Safety \& security & Output (gates consumed by Control flow) & Rulesets \texttt{\{action, resource, effect\}} with last-match-wins and default \texttt{ask}; defaults ask on \texttt{*.\allowbreak{}env}; sub-agent profiles carry deny lists; shell commands parsed into per-command resources; external directories need approval; declines tunnel as defects so a "no" can never become model-visible tool output; compaction checkpoint tells the model it is "not new instructions"; \texttt{<system-\allowbreak{}reminder>} blocks declared harness-owned; sub-agent depth limit; Code Mode interpreter has no fs/imports; no OS-level sandbox found; secrets redacted only in test cassettes and scrubbed from the test env & \texttt{core/\allowbreak{}permission.\allowbreak{}ts:\allowbreak{}L87-\allowbreak{}-\allowbreak{}L101}; \texttt{schema/\allowbreak{}agent.\allowbreak{}ts:\allowbreak{}L46-\allowbreak{}-\allowbreak{}L51}; \texttt{core/\allowbreak{}plugin/\allowbreak{}agent.\allowbreak{}ts:\allowbreak{}L104-\allowbreak{}-\allowbreak{}L117}; \texttt{core/\allowbreak{}shell/\allowbreak{}parse.\allowbreak{}ts:\allowbreak{}L18-\allowbreak{}-\allowbreak{}L21}; \texttt{core/\allowbreak{}file-\allowbreak{}access.\allowbreak{}ts:\allowbreak{}L125-\allowbreak{}-\allowbreak{}L142}; \texttt{core/\allowbreak{}permission.\allowbreak{}ts:\allowbreak{}L247-\allowbreak{}-\allowbreak{}L253}, \texttt{core/\allowbreak{}session/\allowbreak{}model-\allowbreak{}request.\allowbreak{}ts:\allowbreak{}L400-\allowbreak{}-\allowbreak{}L418}; \texttt{core/\allowbreak{}session/\allowbreak{}runner/\allowbreak{}to-\allowbreak{}llm-\allowbreak{}message.\allowbreak{}ts:\allowbreak{}L316-\allowbreak{}-\allowbreak{}L321}; \texttt{core/\allowbreak{}session/\allowbreak{}runner/\allowbreak{}prompt/\allowbreak{}system.\allowbreak{}txt:\allowbreak{}L5}; \texttt{core/\allowbreak{}tool/\allowbreak{}plugin/\allowbreak{}subagent.\allowbreak{}ts:\allowbreak{}L129-\allowbreak{}-\allowbreak{}L133}; \texttt{core/\allowbreak{}codemode/\allowbreak{}tool.\allowbreak{}ts:\allowbreak{}L63-\allowbreak{}-\allowbreak{}L64}; \texttt{recorder/\allowbreak{}redaction/\allowbreak{}secrets.\allowbreak{}ts:\allowbreak{}L3-\allowbreak{}-\allowbreak{}L14}; \texttt{packages/\allowbreak{}core/\allowbreak{}script/\allowbreak{}test.\allowbreak{}ts:\allowbreak{}L10-\allowbreak{}-\allowbreak{}L22} \\
Observability & Model interface, Output & Durable event log doubles as the trace; per-step tokens/cost on \texttt{Step.\allowbreak{}Ended}; per-session token/cost columns and a stats aggregator; named Effect spans on every service function (\texttt{Effect.\allowbreak{}fn(\allowbreak{}"SessionRunner.\allowbreak{}drain")}); WebSocket lifecycle counter metric; warning logs on overflow and transport fallback; event logger for catalog updates & \texttt{core/\allowbreak{}session/\allowbreak{}runner/\allowbreak{}step.\allowbreak{}ts:\allowbreak{}L246-\allowbreak{}-\allowbreak{}L260}; \texttt{core/\allowbreak{}session/\allowbreak{}sql.\allowbreak{}ts:\allowbreak{}L45-\allowbreak{}-\allowbreak{}L50}; \texttt{core/\allowbreak{}session/\allowbreak{}stats.\allowbreak{}ts:\allowbreak{}L17-\allowbreak{}-\allowbreak{}L46}; \texttt{core/\allowbreak{}session/\allowbreak{}runner/\allowbreak{}llm.\allowbreak{}ts:\allowbreak{}L57}; \texttt{core/\allowbreak{}session/\allowbreak{}model-\allowbreak{}transport.\allowbreak{}ts:\allowbreak{}L28-\allowbreak{}-\allowbreak{}L33}; \texttt{core/\allowbreak{}session/\allowbreak{}runner/\allowbreak{}step.\allowbreak{}ts:\allowbreak{}L161-\allowbreak{}-\allowbreak{}L165}; \texttt{core/\allowbreak{}event-\allowbreak{}logger.\allowbreak{}ts:\allowbreak{}L7-\allowbreak{}-\allowbreak{}L17} \\
Resilience & Model interface, Context, Control flow & Jittered exponential retry with retry-after and timeout caps; overflow-triggered compaction with shrinking targets; output-limit and interrupted-stream continuations; WebSocket$\rightarrow$HTTP fallback after 5 failures; header/chunk timeouts; crash recovery from the execution claim with a 10-attempt budget; stale running compactions and tool calls settled on resume; pending permission asks rejected on shutdown & \texttt{core/\allowbreak{}session/\allowbreak{}runner/\allowbreak{}retry.\allowbreak{}ts:\allowbreak{}L42-\allowbreak{}-\allowbreak{}L62}; \texttt{core/\allowbreak{}session/\allowbreak{}compaction.\allowbreak{}ts:\allowbreak{}L95}, \texttt{L429-\allowbreak{}-\allowbreak{}L448}; \texttt{core/\allowbreak{}session/\allowbreak{}runner/\allowbreak{}llm.\allowbreak{}ts:\allowbreak{}L281-\allowbreak{}-\allowbreak{}L319}; \texttt{core/\allowbreak{}session/\allowbreak{}model-\allowbreak{}transport.\allowbreak{}ts:\allowbreak{}L186-\allowbreak{}-\allowbreak{}L196}; \texttt{ai/\allowbreak{}route/\allowbreak{}transport/\allowbreak{}http.\allowbreak{}ts:\allowbreak{}L100-\allowbreak{}-\allowbreak{}L146}; \texttt{core/\allowbreak{}session/\allowbreak{}execution/\allowbreak{}restart.\allowbreak{}ts:\allowbreak{}L77-\allowbreak{}-\allowbreak{}L99}; \texttt{core/\allowbreak{}session/\allowbreak{}runner/\allowbreak{}llm.\allowbreak{}ts:\allowbreak{}L332-\allowbreak{}-\allowbreak{}L385}; \texttt{core/\allowbreak{}permission.\allowbreak{}ts:\allowbreak{}L133-\allowbreak{}-\allowbreak{}L146} \\
Performance & Context, Model interface, Output & Baseline instructions ordered as a stable cache prefix; \texttt{promptCacheKey} shared across a session tree; tool definitions kept on the budget-exhausted step to keep the cached prefix; parallel context source loading; tool calls run concurrently as they stream; 100 ms delta batching; start snapshot captured while the request is in flight; Code Mode catalog and tool JSON-schema caches; tool output truncation and image byte budget & \texttt{core/\allowbreak{}session/\allowbreak{}context.\allowbreak{}ts:\allowbreak{}L147-\allowbreak{}-\allowbreak{}L158}, \texttt{L142}; \texttt{core/\allowbreak{}session/\allowbreak{}model-\allowbreak{}request.\allowbreak{}ts:\allowbreak{}L295}; \texttt{core/\allowbreak{}session/\allowbreak{}runner/\allowbreak{}llm.\allowbreak{}ts:\allowbreak{}L243-\allowbreak{}-\allowbreak{}L244}; \texttt{core/\allowbreak{}session/\allowbreak{}runner/\allowbreak{}step.\allowbreak{}ts:\allowbreak{}L121-\allowbreak{}-\allowbreak{}L134}, \texttt{L74-\allowbreak{}-\allowbreak{}L76}; \texttt{core/\allowbreak{}session/\allowbreak{}runner/\allowbreak{}publish-\allowbreak{}llm-\allowbreak{}event.\allowbreak{}ts:\allowbreak{}L106}; \texttt{core/\allowbreak{}tool.\allowbreak{}ts:\allowbreak{}L246-\allowbreak{}-\allowbreak{}L254}; \texttt{core/\allowbreak{}tool/\allowbreak{}runtime.\allowbreak{}ts:\allowbreak{}L10-\allowbreak{}-\allowbreak{}L19}; \texttt{core/\allowbreak{}tool-\allowbreak{}output.\allowbreak{}ts:\allowbreak{}L13-\allowbreak{}-\allowbreak{}L14}; \texttt{core/\allowbreak{}session/\allowbreak{}model-\allowbreak{}request.\allowbreak{}ts:\allowbreak{}L43-\allowbreak{}-\allowbreak{}L44} \\
Evaluation & --- (all primitives) & 164 test files under \texttt{packages/\allowbreak{}core/\allowbreak{}test} (e.g. \texttt{session-\allowbreak{}runner-\allowbreak{}recorded.\allowbreak{}test.\allowbreak{}ts} replays provider traffic from \texttt{@opencode/\allowbreak{}http-\allowbreak{}recorder} cassettes; \texttt{session-\allowbreak{}compaction.\allowbreak{}test.\allowbreak{}ts}, \texttt{permission.\allowbreak{}test.\allowbreak{}ts}, \texttt{session-\allowbreak{}run-\allowbreak{}coordinator.\allowbreak{}test.\allowbreak{}ts}); 52 under \texttt{packages/\allowbreak{}ai/\allowbreak{}test}; a simulated provider package; test runner scrubs provider keys from the environment. Model-graded or task-success agent evals: not found & \texttt{packages/\allowbreak{}core/\allowbreak{}test/\allowbreak{}session-\allowbreak{}runner-\allowbreak{}recorded.\allowbreak{}test.\allowbreak{}ts:\allowbreak{}L1-\allowbreak{}-\allowbreak{}L30}; \texttt{packages/\allowbreak{}simulation/\allowbreak{}src/\allowbreak{}backend/\allowbreak{}simulated-\allowbreak{}provider.\allowbreak{}ts}; \texttt{packages/\allowbreak{}core/\allowbreak{}script/\allowbreak{}test.\allowbreak{}ts:\allowbreak{}L10-\allowbreak{}-\allowbreak{}L22} \\
\bottomrule \end{longtable}

\subsubsection*{Design decisions (one line each: the choice this harness made within a primitive)}
\begin{itemize}
\item Input: every input, including harness-generated notices, is a durable inbox row with a \texttt{steer} (interrupt the current step boundary) or \texttt{queue} (next turn) delivery, promoted only at step boundaries under a per-session lock.
\item Context: history is event-sourced and re-projected into model messages on every call; compaction is a summary message in that history (or an opaque provider checkpoint), triggered at \textasciitilde{}90 \% of the window using provider-measured counts plus a local estimate, with the newest 15k tokens kept verbatim.
\item Model interface: an in-house canonical \texttt{LLMRequest} compiled per provider protocol, streamed over HTTP/SSE or a stateful WebSocket, with retry decisions split by whether output had already started.
\item Output: tool calls are executed concurrently as they stream, every tool is permission-asserted inside its own code, most tools are reachable only through a single \texttt{execute} tool that runs model-written JavaScript in a custom interpreter.
\item Control flow: one fiber per session drains a durable inbox; a step continues while tool calls were made, stops on a text-only finish, and is bounded by an optional per-agent step count; sub-agents are full child sessions (depth 1 by default) whose results come back as tool text or a synthetic message.
\end{itemize}

\subsubsection*{Hard cases}
\begin{itemize}
\item \textbf{Code Mode \texttt{execute} tool:} the model writes a program that sequences tool calls (\texttt{tools.\allowbreak{}ns["name"](\allowbreak{}input)}), which looked like control flow; rule "a program that sequences only tools is output" resolves it as output, and the interpreter cannot call the model (\texttt{core/\allowbreak{}codemode/\allowbreak{}tool.\allowbreak{}ts:\allowbreak{}L62-\allowbreak{}-\allowbreak{}L69}). Its per-request catalog in the system prompt is context (\texttt{core/\allowbreak{}codemode/\allowbreak{}instructions.\allowbreak{}ts:\allowbreak{}L19-\allowbreak{}-\allowbreak{}L30}).
\item \textbf{\texttt{Permission.\allowbreak{}assert} lives inside each tool:} by file location it is output; by function it is an approval gate the harness consumes (the model never sees a plain decline, the step is interrupted), so it is control flow; the \texttt{"always"} rule save is persistence.
\item \textbf{Question tool:} the form is output, the answers are input (the model consumes them), and a dismissal is treated like a permission decline (control flow) via the same defect tunnel (\texttt{core/\allowbreak{}tool/\allowbreak{}plugin/\allowbreak{}question.\allowbreak{}ts:\allowbreak{}L90-\allowbreak{}-\allowbreak{}L97}).
\item \textbf{Subagent tool:} the call is output; depth check, child creation, blocking and backgrounding are control flow; the child is a separate harness (\texttt{sessions.\allowbreak{}resume}); foreground results return as tool text, background results as a synthetic user message, which makes the same information input in one mode and output in the other.
\item \textbf{Compaction service calls the model:} \texttt{SessionCompaction} decides when and what to summarize (context) but also issues \texttt{llm.\allowbreak{}stream} itself (\texttt{core/\allowbreak{}session/\allowbreak{}compaction.\allowbreak{}ts:\allowbreak{}L512}) and runs the retry policy (\texttt{L450-\allowbreak{}-\allowbreak{}L459}). The brief says context "never calls the model"; resolved by assigning the trigger/split/budget to context, the summary request itself to model interface, and the \texttt{Compacted $\rightarrow$ re-\allowbreak{}attempt} loop to control flow. Partly DOES NOT FIT: native compaction (\texttt{L318-\allowbreak{}-\allowbreak{}L332}) delegates context assembly to the provider, so the harness's context primitive holds only an opaque item it cannot inspect.
\item \textbf{\texttt{MAX\_\allowbreak{}STEPS\_\allowbreak{}PROMPT} appended as an assistant message:} the step budget is control flow, but it is enforced by editing what the model sees (an assistant-role message plus \texttt{toolChoice:\allowbreak{} "none"}), so the mechanism straddles control flow and context (\texttt{core/\allowbreak{}session/\allowbreak{}runner/\allowbreak{}llm.\allowbreak{}ts:\allowbreak{}L240-\allowbreak{}-\allowbreak{}L244}).
\item \textbf{Instruction epochs:} changed AGENTS.md/date/env are frozen into durable \texttt{InstructionsUpdated} events rendered as system messages in the transcript; this is context, but it is implemented as persistence of history, and history is immutable, so stale instruction text stays in old events.
\item \textbf{Plugin hooks (\texttt{session.\allowbreak{}context}, \texttt{session.\allowbreak{}model.\allowbreak{}request}, \texttt{http.\allowbreak{}request}, \texttt{tool.\allowbreak{}execute.\allowbreak{}before/\allowbreak{}after}, \texttt{permission.\allowbreak{}evaluate}, \texttt{session.\allowbreak{}retry}):} a single hook bus whose entries belong to five different primitives; decomposed by what each hook can change rather than where it is registered.
\item \textbf{Durable event bus:} the same event log is persistence (projected into tables), observability (the trace), and the output channel (clients subscribe); the bus itself is out of scope as transport, its contents are assigned by what they record.
\item \textbf{Title model fallback:} the "try the small model, then the primary" fallback (\texttt{core/\allowbreak{}session/\allowbreak{}title.\allowbreak{}ts:\allowbreak{}L129-\allowbreak{}-\allowbreak{}L133}) is a backup model when a request could not be served (model interface), while choosing the small model in the first place is control flow.
\item \textbf{WebSocket \texttt{RecoverFull}:} the transport reports it cannot replay an incremental exchange (model interface), but the decision to resend the full request is taken by the runner's outcome loop (control flow) (\texttt{core/\allowbreak{}session/\allowbreak{}runner/\allowbreak{}step.\allowbreak{}ts:\allowbreak{}L182-\allowbreak{}-\allowbreak{}L188}, \texttt{core/\allowbreak{}session/\allowbreak{}runner/\allowbreak{}llm.\allowbreak{}ts:\allowbreak{}L324-\allowbreak{}-\allowbreak{}L326}).
\item \textbf{Synthetic messages:} "The server restarted\ldots{}", "Your last response hit the output token limit\ldots{}" are produced by control flow, admitted as input, and shown to the model as user-role messages; assigned as input with control-flow origin.
\end{itemize}

\subsubsection*{Summary row}
\begin{longtable}{@{}p{0.18\linewidth}p{0.18\linewidth}p{0.18\linewidth}p{0.18\linewidth}p{0.18\linewidth}@{}}
\toprule \textbf{Input} & \textbf{Context} & \textbf{Model interface} & \textbf{Output} & \textbf{Control flow} \\ \midrule \endhead
Durable inbox rows: user, synthetic, shell, skill; steer or queue & Event-sourced history re-projected; cached-prefix instructions; summary or native compaction & In-house LLMRequest compiled per protocol; SSE or WebSocket; split retry policy & Concurrent streamed tool fibers; permission-asserted tools; Code Mode interpreter & One fiber per session drains inbox; continue while tools called; child sessions \\
\bottomrule \end{longtable}


\subsection{pi (pi-mono coding agent)}

\textbf{Source:} https://github.com/badlogic/pi-mono $\cdot$ \textbf{Version:} commit 5ab7fe74c400097b386662205cd93900695b8a71 (2026-10-15) $\cdot$ \textbf{Language:} TypeScript (Node/Bun)
\textbf{Path conventions:} \texttt{agent/\allowbreak{}} = \texttt{packages/\allowbreak{}agent/\allowbreak{}src/\allowbreak{}}; \texttt{ai/\allowbreak{}} = \texttt{packages/\allowbreak{}ai/\allowbreak{}src/\allowbreak{}}; \texttt{ca/\allowbreak{}} = \texttt{packages/\allowbreak{}coding-\allowbreak{}agent/\allowbreak{}src/\allowbreak{}}; \texttt{ex/\allowbreak{}} = \texttt{packages/\allowbreak{}coding-\allowbreak{}agent/\allowbreak{}examples/\allowbreak{}extensions/\allowbreak{}}; \texttt{ev/\allowbreak{}} = \texttt{packages/\allowbreak{}evals/\allowbreak{}}. All paths relative to the repo root. Package scope at this commit is \texttt{@earendil-\allowbreak{}works/\allowbreak{}*} (pi-agent-core, pi-ai, pi-coding-agent).
\textbf{What it is:} pi is a terminal coding agent by Mario Zechner (badlogic) with a layered design: \texttt{pi-\allowbreak{}agent-\allowbreak{}core} holds a minimal agent loop, \texttt{pi-\allowbreak{}ai} holds provider adapters (Anthropic, OpenAI, Google, Bedrock, Mistral, \ldots{}) behind one \texttt{StreamFunction}, and \texttt{pi-\allowbreak{}coding-\allowbreak{}agent} wraps both in an \texttt{AgentSession} with JSONL session trees, compaction, retries, built-in tools, and an extension system; interactive TUI, print, JSON, and RPC modes sit on top. The model is reached over HTTPS/SSE (or WebSocket for some providers) through provider SDKs or raw fetch, with \texttt{maxRetries:\allowbreak{} 0} so pi owns retry policy.

\subsubsection*{Control flow}
\begin{itemize}
\item \textbf{Inner/outer loop (\texttt{runLoop}):} inner \texttt{while (\allowbreak{}hasMoreToolCalls || pendingMessages.\allowbreak{}length > 0)} issues one model call per iteration; outer \texttt{while (\allowbreak{}true)} re-enters only if follow-up messages arrived or \texttt{finishTurn} asked for one more context-only turn, else breaks and emits \texttt{agent\_\allowbreak{}end}. \texttt{agent/\allowbreak{}agent-\allowbreak{}loop.\allowbreak{}ts:\allowbreak{}L163-\allowbreak{}-\allowbreak{}L321}
\item \textbf{Stop on error/abort:} an assistant \texttt{stopReason} of \texttt{"error"} or \texttt{"aborted"} is a hard exit: \texttt{finishTurn} is called, \texttt{turn\_\allowbreak{}end}/\texttt{agent\_\allowbreak{}end} emitted, loop returns. \texttt{agent/\allowbreak{}agent-\allowbreak{}loop.\allowbreak{}ts:\allowbreak{}L245-\allowbreak{}-\allowbreak{}L256}
\item \textbf{Go-again decision:} \texttt{hasMoreToolCalls = !executedToolBatch.\allowbreak{}terminate}; if the response had no tool calls and no steering/follow-up is queued, the loop ends. \texttt{agent/\allowbreak{}agent-\allowbreak{}loop.\allowbreak{}ts:\allowbreak{}L258-\allowbreak{}-\allowbreak{}L278}, \texttt{L295-\allowbreak{}-\allowbreak{}L318}
\item \textbf{\texttt{finishTurn} decision hook:} returns \texttt{\{action:\allowbreak{}"end"\}} (stop now) or \texttt{\{action:\allowbreak{}"continue"\}} (force one more request); undefined keeps normal scheduling. \texttt{agent/\allowbreak{}agent-\allowbreak{}loop.\allowbreak{}ts:\allowbreak{}L286-\allowbreak{}-\allowbreak{}L298}, \texttt{agent/\allowbreak{}types.\allowbreak{}ts:\allowbreak{}L146-\allowbreak{}-\allowbreak{}L158}
\item \textbf{Early termination from tools:} a batch ends the run only when every finalized tool result sets \texttt{terminate:\allowbreak{} true} (\texttt{shouldTerminateToolBatch}); a blocked \texttt{beforeToolCall} may set it too. \texttt{agent/\allowbreak{}agent-\allowbreak{}loop.\allowbreak{}ts:\allowbreak{}L690-\allowbreak{}-\allowbreak{}L692}, \texttt{L745-\allowbreak{}-\allowbreak{}L755}, \texttt{agent/\allowbreak{}types.\allowbreak{}ts:\allowbreak{}L66-\allowbreak{}-\allowbreak{}L74}
\item \textbf{Steering vs. follow-up queues:} \texttt{Agent.\allowbreak{}steer(\allowbreak{})} injects after the current turn's tools finish; \texttt{followUp(\allowbreak{})} waits until the agent would otherwise stop; drain modes \texttt{"one-\allowbreak{}at-\allowbreak{}a-\allowbreak{}time"} (default) or \texttt{"all"}. \texttt{agent/\allowbreak{}agent.\allowbreak{}ts:\allowbreak{}L143-\allowbreak{}-\allowbreak{}L174}, \texttt{L247-\allowbreak{}-\allowbreak{}L248}, \texttt{L299-\allowbreak{}-\allowbreak{}L306}, \texttt{L496-\allowbreak{}-\allowbreak{}L503}; polled in the loop at \texttt{agent/\allowbreak{}agent-\allowbreak{}loop.\allowbreak{}ts:\allowbreak{}L176}, \texttt{L204-\allowbreak{}-\allowbreak{}L206}, \texttt{L295}, \texttt{L302-\allowbreak{}-\allowbreak{}L308}
\item \textbf{One run at a time:} \texttt{prompt(\allowbreak{})} throws if a run is active; \texttt{continue(\allowbreak{})} resumes from a user/toolResult tail, else drains steering/follow-up queues. \texttt{agent/\allowbreak{}agent.\allowbreak{}ts:\allowbreak{}L373-\allowbreak{}-\allowbreak{}L381}, \texttt{L384-\allowbreak{}-\allowbreak{}L411}, \texttt{L507-\allowbreak{}-\allowbreak{}L530}
\item \textbf{Abort:} \texttt{Agent.\allowbreak{}abort(\allowbreak{})} fires the run's \texttt{AbortController}; \texttt{AgentSession.\allowbreak{}abort(\allowbreak{})} also cancels retry backoff, compaction, branch summary, then waits for idle. \texttt{agent/\allowbreak{}agent.\allowbreak{}ts:\allowbreak{}L341-\allowbreak{}-\allowbreak{}L343}; \texttt{ca/\allowbreak{}core/\allowbreak{}agent-\allowbreak{}session.\allowbreak{}ts:\allowbreak{}L2464-\allowbreak{}-\allowbreak{}L2474}
\item \textbf{Session post-run loop:} \texttt{\_\allowbreak{}startRun} calls \texttt{agent.\allowbreak{}prompt(\allowbreak{}messages)} then loops \texttt{\_\allowbreak{}handlePostAgentRun(\allowbreak{})} $\rightarrow$ \texttt{agent.\allowbreak{}continue(\allowbreak{})} while retry/overflow-recovery/queued messages/extension \texttt{agent\_\allowbreak{}before\_\allowbreak{}settle} ask for more. \texttt{ca/\allowbreak{}core/\allowbreak{}agent-\allowbreak{}session.\allowbreak{}ts:\allowbreak{}L1837-\allowbreak{}-\allowbreak{}L1882}, \texttt{L1943-\allowbreak{}-\allowbreak{}L1981}, \texttt{L1983-\allowbreak{}-\allowbreak{}L2005}
\item \textbf{Auto-retry budget (session level):} \texttt{\_\allowbreak{}prepareRetry} increments \texttt{\_\allowbreak{}retryAttempt}, stops past \texttt{maxRetries}, sleeps \texttt{retryDelayMs(\allowbreak{})} (abortable), omits the failed assistant message from projection, and returns true so the loop re-issues the request. Defaults \texttt{enabled:\allowbreak{} true, maxRetries:\allowbreak{} 3, baseDelayMs:\allowbreak{} 2000}. \texttt{ca/\allowbreak{}core/\allowbreak{}agent-\allowbreak{}session.\allowbreak{}ts:\allowbreak{}L3792-\allowbreak{}-\allowbreak{}L3832}; \texttt{ca/\allowbreak{}core/\allowbreak{}settings-\allowbreak{}defaults.\allowbreak{}ts:\allowbreak{}L18-\allowbreak{}-\allowbreak{}L26}
\item \textbf{Overflow recovery is one-shot:} \texttt{\_\allowbreak{}checkCompaction} case 1 compacts and retries the turn once (\texttt{\_\allowbreak{}overflowRecoveryAttempted} guard); case 2 (completed response) compacts without retry; case 3 threshold compacts without retry. \texttt{ca/\allowbreak{}core/\allowbreak{}agent-\allowbreak{}session.\allowbreak{}ts:\allowbreak{}L2977-\allowbreak{}-\allowbreak{}L3115}
\item \textbf{Virtual-model routing keyed to the work:} \texttt{prepareRequest} resolves a virtual model to a physical one per request with \texttt{reason:\allowbreak{} "user" | "continuation" | "retry"}, storing router state as a session entry; a router (extension-registered) picks the model. \texttt{ca/\allowbreak{}core/\allowbreak{}agent-\allowbreak{}session.\allowbreak{}ts:\allowbreak{}L815-\allowbreak{}-\allowbreak{}L843}; \texttt{ca/\allowbreak{}core/\allowbreak{}model-\allowbreak{}runtime.\allowbreak{}ts:\allowbreak{}L1003-\allowbreak{}-\allowbreak{}L1034}; \texttt{ca/\allowbreak{}core/\allowbreak{}virtual-\allowbreak{}models.\allowbreak{}ts:\allowbreak{}L1-\allowbreak{}-\allowbreak{}L11}, \texttt{L43-\allowbreak{}-\allowbreak{}L50}
\item \textbf{Extension boundary hooks that can force another turn:} \texttt{turn\_\allowbreak{}end} and \texttt{agent\_\allowbreak{}before\_\allowbreak{}settle} handlers return \texttt{continue:\allowbreak{} true} and may commit boundary drafts; the session validates \texttt{canContinue} first. \texttt{ca/\allowbreak{}core/\allowbreak{}agent-\allowbreak{}session.\allowbreak{}ts:\allowbreak{}L847-\allowbreak{}-\allowbreak{}L897}, \texttt{L1983-\allowbreak{}-\allowbreak{}L2005}
\item \textbf{\texttt{tool\_\allowbreak{}call} approval gate (extension):} \texttt{beforeToolCall} $\rightarrow$ \texttt{runner.\allowbreak{}emitToolCall}; first handler returning \texttt{\{block:\allowbreak{}true\}} wins; the \texttt{permission-\allowbreak{}gate} example asks \texttt{ctx.\allowbreak{}ui.\allowbreak{}select} and blocks when no UI. \texttt{ca/\allowbreak{}core/\allowbreak{}agent-\allowbreak{}session.\allowbreak{}ts:\allowbreak{}L659-\allowbreak{}-\allowbreak{}L682}; \texttt{ca/\allowbreak{}core/\allowbreak{}extensions/\allowbreak{}runner.\allowbreak{}ts:\allowbreak{}L1249-\allowbreak{}-\allowbreak{}L1267}; \texttt{ex/\allowbreak{}permission-\allowbreak{}gate.\allowbreak{}ts:\allowbreak{}L13-\allowbreak{}-\allowbreak{}L33}
\item \textbf{Project trust gate:} \texttt{resolveProjectTrusted} checks a persisted trust store, an extension \texttt{project\_\allowbreak{}trust} handler, a \texttt{defaultProjectTrust} setting ("ask" default), then prompts; untrusted cwd means project \texttt{.\allowbreak{}pi/\allowbreak{}} extensions/skills/prompts/SYSTEM.md are not loaded. \texttt{ca/\allowbreak{}core/\allowbreak{}project-\allowbreak{}trust.\allowbreak{}ts:\allowbreak{}L47-\allowbreak{}-\allowbreak{}L97}; \texttt{ca/\allowbreak{}core/\allowbreak{}trust-\allowbreak{}manager.\allowbreak{}ts:\allowbreak{}L30-\allowbreak{}-\allowbreak{}L39}, \texttt{L45-\allowbreak{}-\allowbreak{}L59}, \texttt{L67-\allowbreak{}-\allowbreak{}L97}
\item \textbf{Sub-agents (example extension, not core):} the \texttt{subagent} tool spawns a separate \texttt{pi -\allowbreak{}-\allowbreak{}mode json -\allowbreak{}p -\allowbreak{}-\allowbreak{}no-\allowbreak{}session} process per agent, in \texttt{single}, \texttt{parallel}, or \texttt{chain} modes, and returns the child's final output as the tool result (depth 1; results return as tool-result text). \texttt{ex/\allowbreak{}subagent/\allowbreak{}index.\allowbreak{}ts:\allowbreak{}L298}, \texttt{L342-\allowbreak{}-\allowbreak{}L348}, \texttt{L470-\allowbreak{}-\allowbreak{}L475}; \texttt{ex/\allowbreak{}subagent/\allowbreak{}README.\allowbreak{}md}
\item \textbf{Nested tool calls from tools:} \texttt{ctx.\allowbreak{}executeTool(\allowbreak{})} runs through \texttt{runToolCall} with the session's hooks; ids are \texttt{<parent>/\allowbreak{}<n>} and never enter the transcript (sequences only tools $\rightarrow$ output by rule, but concurrency is serialized by a queue). \texttt{ca/\allowbreak{}core/\allowbreak{}nested-\allowbreak{}tool-\allowbreak{}calls.\allowbreak{}ts:\allowbreak{}L161-\allowbreak{}-\allowbreak{}L200}; \texttt{agent/\allowbreak{}agent-\allowbreak{}loop.\allowbreak{}ts:\allowbreak{}L811-\allowbreak{}-\allowbreak{}L819}; \texttt{ca/\allowbreak{}core/\allowbreak{}extensions/\allowbreak{}types.\allowbreak{}ts:\allowbreak{}L384-\allowbreak{}-\allowbreak{}L395}
\end{itemize}

\subsubsection*{Input}
\begin{itemize}
\item \textbf{TUI submit:} the editor's \texttt{onSubmit} trims text, handles slash commands locally, then \texttt{session.\allowbreak{}prompt(\allowbreak{}text)}; while streaming it calls \texttt{prompt(\allowbreak{}text, \{streamingBehavior:\allowbreak{}"steer"\})}; Alt+Enter queues \texttt{followUp}. \texttt{ca/\allowbreak{}modes/\allowbreak{}interactive/\allowbreak{}interactive-\allowbreak{}mode.\allowbreak{}ts:\allowbreak{}L3200-\allowbreak{}-\allowbreak{}L3202}, \texttt{L3364-\allowbreak{}-\allowbreak{}L3385}, \texttt{L4508}
\item \textbf{Print/JSON mode:} initial message plus \texttt{messages[]} are sent via \texttt{session.\allowbreak{}prompt(\allowbreak{})} sequentially. \texttt{ca/\allowbreak{}modes/\allowbreak{}print-\allowbreak{}mode.\allowbreak{}ts:\allowbreak{}L131-\allowbreak{}-\allowbreak{}L137}
\item \textbf{RPC mode:} JSONL commands \texttt{prompt}, \texttt{steer}, \texttt{follow\_\allowbreak{}up}, \texttt{abort}, \texttt{clear\_\allowbreak{}queue} map onto session methods with \texttt{source:\allowbreak{} "rpc"}. \texttt{ca/\allowbreak{}modes/\allowbreak{}rpc/\allowbreak{}rpc-\allowbreak{}mode.\allowbreak{}ts:\allowbreak{}L394-\allowbreak{}-\allowbreak{}L433}
\item \textbf{\texttt{@file} arguments:} CLI \texttt{@path} args become text content and image attachments before the first prompt. \texttt{ca/\allowbreak{}cli/\allowbreak{}file-\allowbreak{}processor.\allowbreak{}ts:\allowbreak{}L24-\allowbreak{}-\allowbreak{}L25}
\item \textbf{\texttt{AgentSession.\allowbreak{}prompt(\allowbreak{})} capture path:} extension-command check $\rightarrow$ \texttt{input} extension handlers (transform/handled) $\rightarrow$ \texttt{/\allowbreak{}skill:\allowbreak{}} and prompt-template expansion $\rightarrow$ queue if streaming, else auth check and \texttt{\_\allowbreak{}startRun}. \texttt{ca/\allowbreak{}core/\allowbreak{}agent-\allowbreak{}session.\allowbreak{}ts:\allowbreak{}L2058-\allowbreak{}-\allowbreak{}L2141}, \texttt{L2007-\allowbreak{}-\allowbreak{}L2025}; \texttt{ca/\allowbreak{}core/\allowbreak{}extensions/\allowbreak{}runner.\allowbreak{}ts:\allowbreak{}L1525-\allowbreak{}-\allowbreak{}L1564}; \texttt{ca/\allowbreak{}core/\allowbreak{}extensions/\allowbreak{}types.\allowbreak{}ts:\allowbreak{}L1171-\allowbreak{}-\allowbreak{}L1187}
\item \textbf{Queue capture:} \texttt{\_\allowbreak{}queueSteer}/\texttt{\_\allowbreak{}queueFollowUp} build a \texttt{user} message (text + images) and hand it to \texttt{agent.\allowbreak{}steer}/\texttt{followUp}. \texttt{ca/\allowbreak{}core/\allowbreak{}agent-\allowbreak{}session.\allowbreak{}ts:\allowbreak{}L2241-\allowbreak{}-\allowbreak{}L2293}
\item \textbf{Prompt normalization in core:} a string prompt becomes \texttt{\{role:\allowbreak{}"user", content:\allowbreak{}[text, .\allowbreak{}.\allowbreak{}.\allowbreak{}images]\}}. \texttt{agent/\allowbreak{}agent.\allowbreak{}ts:\allowbreak{}L413-\allowbreak{}-\allowbreak{}L430}
\item \textbf{Tool results as input:} each executed call becomes a \texttt{toolResult} message (\texttt{content}, \texttt{details}, \texttt{isError}, \texttt{durationMs}) pushed into context. \texttt{agent/\allowbreak{}agent-\allowbreak{}loop.\allowbreak{}ts:\allowbreak{}L930-\allowbreak{}-\allowbreak{}L944}, \texttt{L274-\allowbreak{}-\allowbreak{}L277}
\item \textbf{Streamed tool output:} \texttt{onUpdate} partial results flow out as \texttt{tool\_\allowbreak{}execution\_\allowbreak{}update} events (bash emits snapshots of its \texttt{OutputAccumulator}). \texttt{agent/\allowbreak{}agent-\allowbreak{}loop.\allowbreak{}ts:\allowbreak{}L779-\allowbreak{}-\allowbreak{}L788}, \texttt{L821-\allowbreak{}-\allowbreak{}L857}; \texttt{ca/\allowbreak{}core/\allowbreak{}tools/\allowbreak{}bash.\allowbreak{}ts:\allowbreak{}L276-\allowbreak{}-\allowbreak{}L294}
\item \textbf{User \texttt{!bash} results:} \texttt{executeBash} runs a user command outside the model loop and records a \texttt{bashExecution} message that \texttt{convertToLlm} renders as a user message (unless \texttt{!!} excluded). \texttt{ca/\allowbreak{}core/\allowbreak{}agent-\allowbreak{}session.\allowbreak{}ts:\allowbreak{}L3871-\allowbreak{}-\allowbreak{}L3903}; \texttt{ca/\allowbreak{}core/\allowbreak{}messages.\allowbreak{}ts:\allowbreak{}L152-\allowbreak{}-\allowbreak{}L161}
\item \textbf{Extension-injected messages:} \texttt{before\_\allowbreak{}agent\_\allowbreak{}start} handlers may return a \texttt{custom} message appended after the user prompt; \texttt{pi.\allowbreak{}sendUserMessage} starts a run. \texttt{ca/\allowbreak{}core/\allowbreak{}agent-\allowbreak{}session.\allowbreak{}ts:\allowbreak{}L1926-\allowbreak{}-\allowbreak{}L1936}; \texttt{ca/\allowbreak{}core/\allowbreak{}extensions/\allowbreak{}runner.\allowbreak{}ts:\allowbreak{}L1427-\allowbreak{}-\allowbreak{}L1477}
\item \textbf{Image inputs:} prompt images are resized per the model's \texttt{inputLimits} before being attached. \texttt{ca/\allowbreak{}core/\allowbreak{}agent-\allowbreak{}session.\allowbreak{}ts:\allowbreak{}L2027-\allowbreak{}-\allowbreak{}L2040}
\end{itemize}

\subsubsection*{Context}
\begin{itemize}
\item \textbf{System prompt sections:} \texttt{buildSystemPromptSections} produces ordered sections \texttt{preamble}, \texttt{tools}, \texttt{rules}, \texttt{docs}, \texttt{addendum}, \texttt{project\_\allowbreak{}context}, \texttt{skills}, \texttt{cwd}, plus custom sections, each wrapped in a same-named XML tag so later patches can replace one section. \texttt{ca/\allowbreak{}core/\allowbreak{}system-\allowbreak{}prompt.\allowbreak{}ts:\allowbreak{}L128-\allowbreak{}-\allowbreak{}L193}
\item \textbf{Default preamble/rules text:} "You are an expert coding assistant operating inside pi\ldots{}"; rules merge per-tool guidelines, extension guidelines, and two fixed rules. \texttt{ca/\allowbreak{}core/\allowbreak{}system-\allowbreak{}prompt.\allowbreak{}ts:\allowbreak{}L155-\allowbreak{}-\allowbreak{}L169}, \texttt{L88-\allowbreak{}-\allowbreak{}L125}
\item \textbf{Prompt diffing into the transcript:} \texttt{diffSystemPromptSections} yields a \texttt{SystemMessage.\allowbreak{}sections} patch only when changed; \texttt{\_\allowbreak{}preparePromptAndToolLoadout} appends it before the next request. \texttt{ca/\allowbreak{}core/\allowbreak{}system-\allowbreak{}prompt.\allowbreak{}ts:\allowbreak{}L217-\allowbreak{}-\allowbreak{}L229}; \texttt{ca/\allowbreak{}core/\allowbreak{}agent-\allowbreak{}session.\allowbreak{}ts:\allowbreak{}L1743-\allowbreak{}-\allowbreak{}L1755}, \texttt{L899-\allowbreak{}-\allowbreak{}L933}
\item \textbf{Project instructions:} \texttt{AGENTS.\allowbreak{}override.\allowbreak{}md}/\texttt{AGENTS.\allowbreak{}md}/\texttt{CLAUDE.\allowbreak{}md} from \texttt{\textasciitilde{}/\allowbreak{}.\allowbreak{}pi/\allowbreak{}agent} and every ancestor of cwd (worktree shadowing handled), rendered as \texttt{<project\_\allowbreak{}instructions path=\ldots{}>}. \texttt{ca/\allowbreak{}core/\allowbreak{}resource-\allowbreak{}loader.\allowbreak{}ts:\allowbreak{}L184-\allowbreak{}-\allowbreak{}L203}, \texttt{L232-\allowbreak{}-\allowbreak{}L274}; \texttt{ca/\allowbreak{}core/\allowbreak{}system-\allowbreak{}prompt.\allowbreak{}ts:\allowbreak{}L79-\allowbreak{}-\allowbreak{}L86}
\item \textbf{SYSTEM.md / APPEND\_SYSTEM.md overrides:} project \texttt{.\allowbreak{}pi/\allowbreak{}SYSTEM.\allowbreak{}md} replaces the preamble; \texttt{APPEND\_\allowbreak{}SYSTEM.\allowbreak{}md} becomes \texttt{addendum}. \texttt{ca/\allowbreak{}core/\allowbreak{}resource-\allowbreak{}loader.\allowbreak{}ts:\allowbreak{}L1214-\allowbreak{}-\allowbreak{}L1233}; \texttt{ca/\allowbreak{}core/\allowbreak{}agent-\allowbreak{}session.\allowbreak{}ts:\allowbreak{}L1714-\allowbreak{}-\allowbreak{}L1730}
\item \textbf{Skills in prompt:} only name/description/location of each \texttt{SKILL.\allowbreak{}md} is listed under \texttt{<available\_\allowbreak{}skills>}; the model loads the file with \texttt{read}. \texttt{ca/\allowbreak{}core/\allowbreak{}skills.\allowbreak{}ts:\allowbreak{}L358-\allowbreak{}-\allowbreak{}L388}
\item \textbf{Tool declarations as transcript deltas:} \texttt{declareToolChanges} computes \texttt{toolsAdded}/\texttt{toolsRemoved} on a system message before each request so replay yields exactly \texttt{context.\allowbreak{}tools}. \texttt{agent/\allowbreak{}agent-\allowbreak{}loop.\allowbreak{}ts:\allowbreak{}L333-\allowbreak{}-\allowbreak{}L363}
\item \textbf{Tool loadout/exposure:} tools have \texttt{exposure} \texttt{direct | model-\allowbreak{}only | codemode | deferred | hidden}; \texttt{tool\_\allowbreak{}search} (BM25) activates deferred tools for the next request; hidden declarations are stripped from every request. \texttt{ca/\allowbreak{}core/\allowbreak{}extensions/\allowbreak{}types.\allowbreak{}ts:\allowbreak{}L516}; \texttt{ca/\allowbreak{}extensions/\allowbreak{}tool-\allowbreak{}search/\allowbreak{}index.\allowbreak{}ts:\allowbreak{}L14}, \texttt{ca/\allowbreak{}extensions/\allowbreak{}tool-\allowbreak{}search/\allowbreak{}tool.\allowbreak{}ts:\allowbreak{}L192-\allowbreak{}-\allowbreak{}L206}; \texttt{ca/\allowbreak{}core/\allowbreak{}agent-\allowbreak{}session.\allowbreak{}ts:\allowbreak{}L1772-\allowbreak{}-\allowbreak{}L1790}
\item \textbf{Canonical context = session projection:} \texttt{prepareRequest} replaces \texttt{context.\allowbreak{}messages} with \texttt{sessionManager.\allowbreak{}buildSessionProjection(\allowbreak{}).\allowbreak{}messages}, which walks the tree path, applies the newest compaction (\texttt{firstKeptEntryId}) and \texttt{context\_\allowbreak{}edit} replacements. \texttt{ca/\allowbreak{}core/\allowbreak{}agent-\allowbreak{}session.\allowbreak{}ts:\allowbreak{}L788-\allowbreak{}-\allowbreak{}L811}; \texttt{ca/\allowbreak{}core/\allowbreak{}session-\allowbreak{}manager.\allowbreak{}ts:\allowbreak{}L476-\allowbreak{}-\allowbreak{}L512}, \texttt{L543-\allowbreak{}-\allowbreak{}L573}
\item \textbf{\texttt{convertToLlm}:} maps \texttt{bashExecution}, \texttt{custom}, \texttt{branchSummary}, \texttt{compactionSummary} roles to \texttt{user} messages with prefixes and passes system/user/assistant/toolResult through. \texttt{ca/\allowbreak{}core/\allowbreak{}messages.\allowbreak{}ts:\allowbreak{}L148-\allowbreak{}-\allowbreak{}L196}
\item \textbf{Extension \texttt{context} / \texttt{context\_\allowbreak{}with\_\allowbreak{}system} hooks:} run inside \texttt{transformContext} before every call; \texttt{context} sees non-system messages (clone), \texttt{context\_\allowbreak{}with\_\allowbreak{}system} owns the full list. \texttt{ca/\allowbreak{}core/\allowbreak{}sdk.\allowbreak{}ts:\allowbreak{}L436-\allowbreak{}-\allowbreak{}L440}; \texttt{ca/\allowbreak{}core/\allowbreak{}extensions/\allowbreak{}runner.\allowbreak{}ts:\allowbreak{}L1305-\allowbreak{}-\allowbreak{}L1366}; \texttt{agent/\allowbreak{}agent-\allowbreak{}loop.\allowbreak{}ts:\allowbreak{}L388-\allowbreak{}-\allowbreak{}L397}
\item \textbf{Forced prompt projection:} a \texttt{before\_\allowbreak{}agent\_\allowbreak{}start} \texttt{systemPrompt} override collapses system messages into one head for the request only; the transcript keeps structured sections. \texttt{ca/\allowbreak{}core/\allowbreak{}agent-\allowbreak{}session.\allowbreak{}ts:\allowbreak{}L1792-\allowbreak{}-\allowbreak{}L1807}
\item \textbf{Compaction trigger:} \texttt{shouldCompact} fires when \texttt{contextTokens > contextWindow -\allowbreak{} reserveTokens} (defaults \texttt{reserveTokens 16384}, \texttt{keepRecentTokens 20000}); checked before the next assistant response and after the run. \texttt{ca/\allowbreak{}core/\allowbreak{}compaction/\allowbreak{}compaction.\allowbreak{}ts:\allowbreak{}L120-\allowbreak{}-\allowbreak{}L130}, \texttt{L267-\allowbreak{}-\allowbreak{}L270}; \texttt{ca/\allowbreak{}core/\allowbreak{}agent-\allowbreak{}session.\allowbreak{}ts:\allowbreak{}L768-\allowbreak{}-\allowbreak{}L786}, \texttt{L3079-\allowbreak{}-\allowbreak{}L3113}
\item \textbf{Cut point:} walks back from newest accumulating estimated tokens until \texttt{keepRecentTokens}, then picks the nearest valid cut (turn boundary; can split a turn and summarize its prefix). \texttt{ca/\allowbreak{}core/\allowbreak{}compaction/\allowbreak{}compaction.\allowbreak{}ts:\allowbreak{}L446-\allowbreak{}-\allowbreak{}L480}, \texttt{L872-\allowbreak{}-\allowbreak{}L905}
\item \textbf{Summary generation:} a separate model call with a fixed structured prompt (\texttt{\#\# Goal}, \texttt{\#\# Constraints \& Preferences}, \texttt{\#\# Progress} \ldots{}), update prompt when a previous summary exists, \texttt{cacheRetention:\allowbreak{} "none"}, wrapped in \texttt{retryAssistantCall}. \texttt{ca/\allowbreak{}core/\allowbreak{}compaction/\allowbreak{}compaction.\allowbreak{}ts:\allowbreak{}L507-\allowbreak{}-\allowbreak{}L520}, \texttt{L619-\allowbreak{}-\allowbreak{}L639}, \texttt{L682-\allowbreak{}-\allowbreak{}L693}
\item \textbf{Compaction persisted as a tree entry:} \texttt{appendCompaction} stores \texttt{summary}, \texttt{firstKeptEntryId}, \texttt{tokensBefore}, usage, and the current system message. \texttt{ca/\allowbreak{}core/\allowbreak{}session-\allowbreak{}manager.\allowbreak{}ts:\allowbreak{}L1261-\allowbreak{}-\allowbreak{}L1287}
\item \textbf{Token estimation:} provider \texttt{usage.\allowbreak{}totalTokens} when present, else \texttt{chars /\allowbreak{} 3.\allowbreak{}5} and 4800 chars per image. \texttt{ai/\allowbreak{}utils/\allowbreak{}estimate.\allowbreak{}ts:\allowbreak{}L15-\allowbreak{}-\allowbreak{}L19}, \texttt{L38-\allowbreak{}-\allowbreak{}L40}
\item \textbf{Tool-output trimming:} read/bash truncate to 2000 lines or 50 KB (head for read, tail for bash) with full output spilled to a temp file. \texttt{ca/\allowbreak{}core/\allowbreak{}tools/\allowbreak{}truncate.\allowbreak{}ts:\allowbreak{}L11-\allowbreak{}-\allowbreak{}L12}; \texttt{ca/\allowbreak{}core/\allowbreak{}tools/\allowbreak{}bash.\allowbreak{}ts:\allowbreak{}L255}; \texttt{ca/\allowbreak{}core/\allowbreak{}tools/\allowbreak{}read.\allowbreak{}ts:\allowbreak{}L96}; \texttt{ca/\allowbreak{}core/\allowbreak{}bash-\allowbreak{}executor.\allowbreak{}ts:\allowbreak{}L56-\allowbreak{}-\allowbreak{}L104}, \texttt{L130-\allowbreak{}-\allowbreak{}L145}
\item \textbf{Image policy in context:} \texttt{blockImages} setting strips images in \texttt{convertToLlm}; tool-result images are resized per model limits after \texttt{tool\_\allowbreak{}result} hooks. \texttt{ca/\allowbreak{}core/\allowbreak{}sdk.\allowbreak{}ts:\allowbreak{}L296-\allowbreak{}-\allowbreak{}L310}; \texttt{ca/\allowbreak{}core/\allowbreak{}agent-\allowbreak{}session.\allowbreak{}ts:\allowbreak{}L705-\allowbreak{}-\allowbreak{}L711}
\item \textbf{Branch summaries:} navigating the tree back can append a \texttt{branch\_\allowbreak{}summary} entry rendered as a user message. \texttt{ca/\allowbreak{}core/\allowbreak{}session-\allowbreak{}manager.\allowbreak{}ts:\allowbreak{}L1600-\allowbreak{}-\allowbreak{}L1615}; \texttt{ca/\allowbreak{}core/\allowbreak{}messages.\allowbreak{}ts:\allowbreak{}L170-\allowbreak{}-\allowbreak{}L175}
\end{itemize}

\subsubsection*{Model interface}
\begin{itemize}
\item \textbf{\texttt{StreamFunction} contract:} \texttt{(\allowbreak{}model, normalizedTranscript, options) => AssistantMessageEventStream}; errors must terminate the stream with \texttt{stopReason "error"|"aborted"}. \texttt{ai/\allowbreak{}types.\allowbreak{}ts:\allowbreak{}L379-\allowbreak{}-\allowbreak{}L394}
\item \textbf{Loop $\rightarrow$ stream boundary:} \texttt{streamAssistantResponse} resolves the API key via \texttt{getApiKey} (expiring tokens), calls \texttt{streamFunction}, and folds \texttt{start}/\texttt{*\_\allowbreak{}delta}/\texttt{done}/\texttt{error} events into a partial \texttt{AssistantMessage} in context. \texttt{agent/\allowbreak{}agent-\allowbreak{}loop.\allowbreak{}ts:\allowbreak{}L399-\allowbreak{}-\allowbreak{}L407}, \texttt{L414-\allowbreak{}-\allowbreak{}L458}
\item \textbf{Session stream function:} \texttt{modelRuntime.\allowbreak{}streamSimple} $\rightarrow$ \texttt{lazyStream} (async setup behind a sync stream) $\rightarrow$ \texttt{prepareRequest} (auth/headers) $\rightarrow$ provider \texttt{streamSimple}. \texttt{ca/\allowbreak{}core/\allowbreak{}sdk.\allowbreak{}ts:\allowbreak{}L420-\allowbreak{}-\allowbreak{}L431}; \texttt{ca/\allowbreak{}core/\allowbreak{}model-\allowbreak{}runtime.\allowbreak{}ts:\allowbreak{}L723-\allowbreak{}-\allowbreak{}L749}; \texttt{ai/\allowbreak{}api/\allowbreak{}lazy.\allowbreak{}ts:\allowbreak{}L46-\allowbreak{}-\allowbreak{}L62}; \texttt{ai/\allowbreak{}models.\allowbreak{}ts:\allowbreak{}L907-\allowbreak{}-\allowbreak{}L914}
\item \textbf{Anthropic adapter:} builds a \texttt{@anthropic-\allowbreak{}ai/\allowbreak{}sdk} client (API key or OAuth bearer with Claude Code identity headers), calls \texttt{client.\allowbreak{}beta.\allowbreak{}messages.\allowbreak{}create(\allowbreak{}params, \{maxRetries:\allowbreak{} 0, signal, timeout\})}, decodes SSE by hand. \texttt{ai/\allowbreak{}api/\allowbreak{}anthropic-\allowbreak{}messages.\allowbreak{}ts:\allowbreak{}L603-\allowbreak{}-\allowbreak{}L645}, \texttt{L1018-\allowbreak{}-\allowbreak{}L1060}, \texttt{L475-\allowbreak{}-\allowbreak{}L500}, \texttt{L534}
\item \textbf{Provider-level retry:} \texttt{retryProviderRequest} mirrors SDK policy (408/409/429/5xx or \texttt{x-\allowbreak{}should-\allowbreak{}retry}), honors \texttt{retry-\allowbreak{}after(\allowbreak{}-\allowbreak{}ms)} up to \texttt{maxRetryDelayMs} (default 60 s, else fail fast so the outer policy sees it), otherwise \texttt{min(\allowbreak{}0.\allowbreak{}5$\cdot$2\textasciicircum{}n, 8)s} with 25 \% jitter, abortable. \texttt{ai/\allowbreak{}utils/\allowbreak{}provider-\allowbreak{}retry.\allowbreak{}ts:\allowbreak{}L25-\allowbreak{}-\allowbreak{}L37}, \texttt{L53-\allowbreak{}-\allowbreak{}L69}, \texttt{L107-\allowbreak{}-\allowbreak{}L128}; \texttt{ai/\allowbreak{}api/\allowbreak{}anthropic-\allowbreak{}messages.\allowbreak{}ts:\allowbreak{}L658-\allowbreak{}-\allowbreak{}L666}
\item \textbf{Assistant-level retry classifier:} \texttt{isRetryableAssistantError} matches overloaded/429/5xx/network text but excludes quota/billing patterns; \texttt{retryAssistantCall} does bounded exponential backoff \texttt{baseDelayMs$\cdot$2\textasciicircum{}(\allowbreak{}n-\allowbreak{}1)} capped at 60 s. \texttt{ai/\allowbreak{}utils/\allowbreak{}retry.\allowbreak{}ts:\allowbreak{}L7-\allowbreak{}-\allowbreak{}L28}, \texttt{L30-\allowbreak{}-\allowbreak{}L60}, \texttt{L121-\allowbreak{}-\allowbreak{}L137}, \texttt{L196-\allowbreak{}-\allowbreak{}L246}, \texttt{L257-\allowbreak{}-\allowbreak{}L262}
\item \textbf{Overflow detection:} regex table over provider error texts (\texttt{prompt is too long}, \texttt{request\_\allowbreak{}too\_\allowbreak{}large}, \ldots{}) with non-overflow exclusions (rate limit), used to divert to compaction instead of retry. \texttt{ai/\allowbreak{}utils/\allowbreak{}overflow.\allowbreak{}ts:\allowbreak{}L37-\allowbreak{}-\allowbreak{}L63}, \texttt{L76-\allowbreak{}-\allowbreak{}L80}; \texttt{ca/\allowbreak{}core/\allowbreak{}agent-\allowbreak{}session.\allowbreak{}ts:\allowbreak{}L3739-\allowbreak{}-\allowbreak{}L3743}
\item \textbf{Strict-tools fallback within the adapter:} a 400 "compiled grammar too large"/"Too many strict tools" rebuilds params without preferred strict schemas and retries once. \texttt{ai/\allowbreak{}api/\allowbreak{}anthropic-\allowbreak{}messages.\allowbreak{}ts:\allowbreak{}L667-\allowbreak{}-\allowbreak{}L691}
\item \textbf{Stream termination:} abort $\rightarrow$ \texttt{stopReason "aborted"}, exceptions $\rightarrow$ \texttt{"error"} with \texttt{errorMessage}; \texttt{done} carries the final message with usage and cost. \texttt{ai/\allowbreak{}api/\allowbreak{}anthropic-\allowbreak{}messages.\allowbreak{}ts:\allowbreak{}L897-\allowbreak{}-\allowbreak{}L933}; \texttt{ai/\allowbreak{}utils/\allowbreak{}event-\allowbreak{}stream.\allowbreak{}ts:\allowbreak{}L26-\allowbreak{}-\allowbreak{}L58}
\item \textbf{Prompt caching:} \texttt{cache\_\allowbreak{}control} on the last user/system block and last tool declaration; retention \texttt{short}/\texttt{long}/\texttt{none}; \texttt{sessionId} forwarded for cache-aware backends. \texttt{ai/\allowbreak{}api/\allowbreak{}anthropic-\allowbreak{}messages.\allowbreak{}ts:\allowbreak{}L1529-\allowbreak{}-\allowbreak{}L1555}, \texttt{L1659}; \texttt{ai/\allowbreak{}types.\allowbreak{}ts:\allowbreak{}L239-\allowbreak{}-\allowbreak{}L249}
\item \textbf{Thinking mapping:} pi \texttt{ThinkingLevel} maps to Anthropic effort or \texttt{budget\_\allowbreak{}tokens} (adaptive vs. budget thinking), interleaved-thinking beta header. \texttt{ai/\allowbreak{}api/\allowbreak{}anthropic-\allowbreak{}messages.\allowbreak{}ts:\allowbreak{}L944-\allowbreak{}-\allowbreak{}L1012}, \texttt{L1296-\allowbreak{}-\allowbreak{}L1305}, \texttt{L195}
\item \textbf{Wire-level hooks:} \texttt{onPayload} (extension \texttt{before\_\allowbreak{}provider\_\allowbreak{}request} may replace the body), \texttt{onResponse} (\texttt{after\_\allowbreak{}provider\_\allowbreak{}response}), \texttt{onProviderStreamEvent}. \texttt{ai/\allowbreak{}api/\allowbreak{}anthropic-\allowbreak{}messages.\allowbreak{}ts:\allowbreak{}L653-\allowbreak{}-\allowbreak{}L656}, \texttt{L693}, \texttt{L700}; \texttt{ca/\allowbreak{}core/\allowbreak{}sdk.\allowbreak{}ts:\allowbreak{}L382-\allowbreak{}-\allowbreak{}L409}; \texttt{ai/\allowbreak{}types.\allowbreak{}ts:\allowbreak{}L174-\allowbreak{}-\allowbreak{}L182}, \texttt{L218-\allowbreak{}-\allowbreak{}L223}
\item \textbf{Transport/options:} \texttt{transport:\allowbreak{} "sse"|"websocket"|\ldots{}}, \texttt{timeoutMs}, \texttt{maxRetries}, \texttt{maxRetryDelayMs}, custom \texttt{fetch}, custom \texttt{headers}. \texttt{ai/\allowbreak{}types.\allowbreak{}ts:\allowbreak{}L144}, \texttt{L157-\allowbreak{}-\allowbreak{}L210}, \texttt{L235-\allowbreak{}-\allowbreak{}L238}
\item \textbf{Credential handling:} OAuth tokens auto-refresh under a file lock with a timeout; \texttt{auth.\allowbreak{}json} written with mode \texttt{0o600} in a \texttt{0o700} dir. \texttt{ai/\allowbreak{}auth/\allowbreak{}resolve.\allowbreak{}ts:\allowbreak{}L119-\allowbreak{}-\allowbreak{}L142}; \texttt{ca/\allowbreak{}core/\allowbreak{}auth-\allowbreak{}storage.\allowbreak{}ts:\allowbreak{}L37}, \texttt{L71}
\item \textbf{Direct-route virtual model outside the loop:} \texttt{streamSimple} on a virtual model routes with \texttt{reason:\allowbreak{} "direct"} and withholds the caller's credentials when the provider differs. \texttt{ca/\allowbreak{}core/\allowbreak{}model-\allowbreak{}runtime.\allowbreak{}ts:\allowbreak{}L725-\allowbreak{}-\allowbreak{}L743}
\end{itemize}

\subsubsection*{Output}
\begin{itemize}
\item \textbf{Tool execution pipeline:} \texttt{prepareToolCall} (lookup, \texttt{prepareArguments}, schema validation, \texttt{beforeToolCall}) $\rightarrow$ \texttt{executePreparedToolCall} (\texttt{tool.\allowbreak{}execute(\allowbreak{}id, args, signal, onUpdate)} in try/catch, duration measured) $\rightarrow$ \texttt{finalizeExecutedToolCall} (\texttt{afterToolCall} may replace content/details/isError). \texttt{agent/\allowbreak{}agent-\allowbreak{}loop.\allowbreak{}ts:\allowbreak{}L708-\allowbreak{}-\allowbreak{}L777}, \texttt{L821-\allowbreak{}-\allowbreak{}L857}, \texttt{L859-\allowbreak{}-\allowbreak{}L910}
\item \textbf{Parallel vs. sequential:} default \texttt{"parallel"}: preflight sequentially, run allowed tools concurrently with \texttt{Promise.\allowbreak{}all}, emit results in source order; \texttt{"sequential"} if configured or any called tool has \texttt{executionMode:\allowbreak{} "sequential"}. \texttt{agent/\allowbreak{}agent-\allowbreak{}loop.\allowbreak{}ts:\allowbreak{}L508-\allowbreak{}-\allowbreak{}L523}, \texttt{L586-\allowbreak{}-\allowbreak{}L660}; \texttt{agent/\allowbreak{}types.\allowbreak{}ts:\allowbreak{}L308-\allowbreak{}-\allowbreak{}L317}
\item \textbf{Truncated-output guard:} if the response \texttt{stopReason} is \texttt{"length"}, every tool call is failed with an error result instead of executed. \texttt{agent/\allowbreak{}agent-\allowbreak{}loop.\allowbreak{}ts:\allowbreak{}L263-\allowbreak{}-\allowbreak{}L270}, \texttt{L478-\allowbreak{}-\allowbreak{}L503}
\item \textbf{Built-in tool set:} \texttt{read, bash, powershell, edit, write, grep, find, ls} (default four: read, bash, edit, write). \texttt{ca/\allowbreak{}core/\allowbreak{}tools/\allowbreak{}index.\allowbreak{}ts:\allowbreak{}L95-\allowbreak{}-\allowbreak{}L105}; \texttt{ca/\allowbreak{}core/\allowbreak{}system-\allowbreak{}prompt.\allowbreak{}ts:\allowbreak{}L64}
\item \textbf{bash executor:} \texttt{spawn(\allowbreak{}shell, args, \{cwd, detached, env:\allowbreak{} process.\allowbreak{}env + bin dir\})}, stdout/stderr streamed, optional per-call timeout, abort kills the whole process group with \texttt{SIGKILL} (\texttt{taskkill /\allowbreak{}T} on Windows). \texttt{ca/\allowbreak{}core/\allowbreak{}tools/\allowbreak{}bash.\allowbreak{}ts:\allowbreak{}L97-\allowbreak{}-\allowbreak{}L166}; \texttt{ca/\allowbreak{}utils/\allowbreak{}shell.\allowbreak{}ts:\allowbreak{}L138-\allowbreak{}-\allowbreak{}L150}, \texttt{L187-\allowbreak{}-\allowbreak{}L215}
\item \textbf{edit tool:} exact unique \texttt{oldText}$\rightarrow$\texttt{newText} replacements against the original file, BOM/LF normalization, fuzzy-match fallback with diff details; file mutations serialized per real path. \texttt{ca/\allowbreak{}core/\allowbreak{}tools/\allowbreak{}edit.\allowbreak{}ts:\allowbreak{}L152}, \texttt{L190-\allowbreak{}-\allowbreak{}L194}; \texttt{ca/\allowbreak{}core/\allowbreak{}tools/\allowbreak{}edit-\allowbreak{}diff.\allowbreak{}ts:\allowbreak{}L28-\allowbreak{}-\allowbreak{}L37}; \texttt{ca/\allowbreak{}core/\allowbreak{}tools/\allowbreak{}file-\allowbreak{}mutation-\allowbreak{}queue.\allowbreak{}ts:\allowbreak{}L32-\allowbreak{}-\allowbreak{}L61}
\item \textbf{read tool:} text with \texttt{offset}/\texttt{limit} and head truncation; images returned as image blocks (resized) or omitted for non-vision models. \texttt{ca/\allowbreak{}core/\allowbreak{}tools/\allowbreak{}read.\allowbreak{}ts:\allowbreak{}L96}, \texttt{L159-\allowbreak{}-\allowbreak{}L177}, \texttt{L80-\allowbreak{}-\allowbreak{}L83}
\item \textbf{grep/find:} spawn ripgrep (\texttt{rg}, downloaded if missing). \texttt{ca/\allowbreak{}core/\allowbreak{}tools/\allowbreak{}grep.\allowbreak{}ts:\allowbreak{}L119-\allowbreak{}-\allowbreak{}L121}, \texttt{L168}
\item \textbf{MCP tools:} the built-in \texttt{mcp} extension registers each connected server's tools via \texttt{pi.\allowbreak{}registerTool} (hidden exposure when only reachable from codemode) and forwards calls with \texttt{client.\allowbreak{}callTool}. \texttt{ca/\allowbreak{}extensions/\allowbreak{}mcp/\allowbreak{}index.\allowbreak{}ts:\allowbreak{}L387-\allowbreak{}-\allowbreak{}L460}; \texttt{ca/\allowbreak{}extensions/\allowbreak{}mcp/\allowbreak{}tools.\allowbreak{}ts:\allowbreak{}L337-\allowbreak{}-\allowbreak{}L339}
\item \textbf{codemode tool:} the model writes JavaScript that calls other tools; it runs in a QuickJS WASM sandbox (256 MiB heap, 2 workers max, 10k-token output cap, optional \texttt{timeout\_\allowbreak{}ms}), with no Node/fs/network; nested calls use \texttt{ctx.\allowbreak{}executeTool}. \texttt{ca/\allowbreak{}extensions/\allowbreak{}codemode/\allowbreak{}tool.\allowbreak{}ts:\allowbreak{}L1-\allowbreak{}-\allowbreak{}L10}, \texttt{L118-\allowbreak{}-\allowbreak{}L120}; \texttt{ca/\allowbreak{}extensions/\allowbreak{}codemode/\allowbreak{}execute.\allowbreak{}ts:\allowbreak{}L50-\allowbreak{}-\allowbreak{}L59}, \texttt{L254}, \texttt{L510-\allowbreak{}-\allowbreak{}L522}
\item \textbf{Rendering streamed fragments:} the TUI updates a streaming assistant component on each \texttt{message\_\allowbreak{}update} and creates tool components as tool calls appear; JSON mode writes each event as one JSON line; print mode prints final assistant text. \texttt{ca/\allowbreak{}modes/\allowbreak{}interactive/\allowbreak{}interactive-\allowbreak{}mode.\allowbreak{}ts:\allowbreak{}L3524-\allowbreak{}-\allowbreak{}L3536}; \texttt{ca/\allowbreak{}modes/\allowbreak{}print-\allowbreak{}mode.\allowbreak{}ts:\allowbreak{}L108-\allowbreak{}-\allowbreak{}L112}, \texttt{L139-\allowbreak{}-\allowbreak{}L158}
\item \textbf{Extension UI as output:} \texttt{ctx.\allowbreak{}ui.\allowbreak{}select/\allowbreak{}confirm/\allowbreak{}input/\allowbreak{}editor} render dialogs (also used by approval gates). \texttt{ca/\allowbreak{}core/\allowbreak{}extensions/\allowbreak{}types.\allowbreak{}ts:\allowbreak{}L149-\allowbreak{}-\allowbreak{}L157}, \texttt{L240}
\end{itemize}

\subsubsection*{Production concerns}
\begin{longtable}{@{}p{0.13\linewidth}p{0.17\linewidth}p{0.36\linewidth}p{0.26\linewidth}@{}}
\toprule \textbf{Concern} & \textbf{Where it lives (primitive)} & \textbf{Mechanism} & \textbf{Evidence} \\ \midrule \endhead
Persistence & Context (session transcript) & Append-only JSONL under \texttt{\textasciitilde{}/\allowbreak{}.\allowbreak{}pi/\allowbreak{}agent/\allowbreak{}sessions/\allowbreak{}-\allowbreak{}-\allowbreak{}<cwd>-\allowbreak{}-\allowbreak{}/\allowbreak{}<ts>\_\allowbreak{}<id>.\allowbreak{}jsonl}, v3 entries with \texttt{id}/\texttt{parentId} tree; file created on first user/assistant message; branches via \texttt{branch(\allowbreak{})}; model/thinking changes and usage persisted as entries; \texttt{context\_\allowbreak{}edit} entries persist recovery omissions & \texttt{ca/\allowbreak{}core/\allowbreak{}session-\allowbreak{}manager.\allowbreak{}ts:\allowbreak{}L41}, \texttt{L589-\allowbreak{}-\allowbreak{}L594}, \texttt{L1057-\allowbreak{}-\allowbreak{}L1083}, \texttt{L1166-\allowbreak{}-\allowbreak{}L1214}, \texttt{L1244}, \texttt{L1360}, \texttt{L1579-\allowbreak{}-\allowbreak{}L1584}; \texttt{ca/\allowbreak{}core/\allowbreak{}agent-\allowbreak{}session.\allowbreak{}ts:\allowbreak{}L1147-\allowbreak{}-\allowbreak{}L1168}, \texttt{L1241-\allowbreak{}-\allowbreak{}L1257} \\
Safety \& security & Control flow (gates), Output (sandbox), Model interface (secrets) & Project trust gate before loading repo-controlled extensions/skills; \texttt{tool\_\allowbreak{}call} block hook (example permission gate); OS sandbox only via example extension (\texttt{@anthropic-\allowbreak{}ai/\allowbreak{}sandbox-\allowbreak{}runtime}: sandbox-exec/bubblewrap); codemode scripts in QuickJS with no fs/network; auth file 0600; bash inherits full \texttt{process.\allowbreak{}env} (no secret scrubbing found); nothing found for prompt-injection defenses beyond trust gating & \texttt{ca/\allowbreak{}core/\allowbreak{}project-\allowbreak{}trust.\allowbreak{}ts:\allowbreak{}L47-\allowbreak{}-\allowbreak{}L97}; \texttt{ca/\allowbreak{}core/\allowbreak{}trust-\allowbreak{}manager.\allowbreak{}ts:\allowbreak{}L30-\allowbreak{}-\allowbreak{}L39}; \texttt{ex/\allowbreak{}permission-\allowbreak{}gate.\allowbreak{}ts:\allowbreak{}L13-\allowbreak{}-\allowbreak{}L33}; \texttt{ex/\allowbreak{}sandbox/\allowbreak{}index.\allowbreak{}ts:\allowbreak{}L1-\allowbreak{}-\allowbreak{}L10}; \texttt{ca/\allowbreak{}extensions/\allowbreak{}codemode/\allowbreak{}tool.\allowbreak{}ts:\allowbreak{}L118}; \texttt{ca/\allowbreak{}core/\allowbreak{}auth-\allowbreak{}storage.\allowbreak{}ts:\allowbreak{}L37}; \texttt{ca/\allowbreak{}utils/\allowbreak{}shell.\allowbreak{}ts:\allowbreak{}L138-\allowbreak{}-\allowbreak{}L150} \\
Observability & (observability) & Lifecycle events (\texttt{agent\_\allowbreak{}start/\allowbreak{}turn\_\allowbreak{}*/\allowbreak{}message\_\allowbreak{}*/\allowbreak{}tool\_\allowbreak{}execution\_\allowbreak{}*}) to listeners/extensions; \texttt{usage} (input/output/cache/cost) on every assistant message and \texttt{usage} session entries; \texttt{getContextUsage}/\texttt{getSessionStats}; crash log; redacted bug-report bundle; \texttt{TelemetryContext} span interface passed through options but no \texttt{startSpan} call found in \texttt{ai/\allowbreak{}} & \texttt{agent/\allowbreak{}agent.\allowbreak{}ts:\allowbreak{}L565-\allowbreak{}-\allowbreak{}L612}; \texttt{ai/\allowbreak{}types.\allowbreak{}ts:\allowbreak{}L451-\allowbreak{}-\allowbreak{}L470}, \texttt{L583}; \texttt{ca/\allowbreak{}core/\allowbreak{}session-\allowbreak{}manager.\allowbreak{}ts:\allowbreak{}L1244}; \texttt{ca/\allowbreak{}core/\allowbreak{}agent-\allowbreak{}session.\allowbreak{}ts:\allowbreak{}L4267-\allowbreak{}-\allowbreak{}L4290}; \texttt{ca/\allowbreak{}core/\allowbreak{}crash-\allowbreak{}log.\allowbreak{}ts:\allowbreak{}L124-\allowbreak{}-\allowbreak{}L144}; \texttt{ca/\allowbreak{}core/\allowbreak{}bug-\allowbreak{}report.\allowbreak{}ts:\allowbreak{}L27}, \texttt{L51}; \texttt{packages/\allowbreak{}telemetry/\allowbreak{}src/\allowbreak{}index.\allowbreak{}ts:\allowbreak{}L1-\allowbreak{}-\allowbreak{}L22}; \texttt{ai/\allowbreak{}types.\allowbreak{}ts:\allowbreak{}L159-\allowbreak{}-\allowbreak{}L160} \\
Resilience & Model interface (retries), Context (compaction), Control flow (recovery loop) & Two retry tiers (provider HTTP retry with retry-after cap; assistant-level exponential backoff, 3 attempts); overflow $\rightarrow$ compact-and-retry once; \texttt{length} stop fails tool calls; abortable sleeps; resume from JSONL on restart; compaction itself retried & \texttt{ai/\allowbreak{}utils/\allowbreak{}provider-\allowbreak{}retry.\allowbreak{}ts:\allowbreak{}L107-\allowbreak{}-\allowbreak{}L128}; \texttt{ai/\allowbreak{}utils/\allowbreak{}retry.\allowbreak{}ts:\allowbreak{}L196-\allowbreak{}-\allowbreak{}L246}; \texttt{ca/\allowbreak{}core/\allowbreak{}agent-\allowbreak{}session.\allowbreak{}ts:\allowbreak{}L3040-\allowbreak{}-\allowbreak{}L3077}, \texttt{L3792-\allowbreak{}-\allowbreak{}L3832}; \texttt{agent/\allowbreak{}agent-\allowbreak{}loop.\allowbreak{}ts:\allowbreak{}L263-\allowbreak{}-\allowbreak{}L270}; \texttt{ca/\allowbreak{}core/\allowbreak{}compaction/\allowbreak{}compaction.\allowbreak{}ts:\allowbreak{}L619-\allowbreak{}-\allowbreak{}L639} \\
Performance & Model interface (cache), Output (parallel tools), Context (trimming) & Prompt \texttt{cache\_\allowbreak{}control} + \texttt{sessionId}; \texttt{CacheWarmer} re-sends a 1-token request at 90 \% of cache TTL when expected savings $\geq$ \$0.05; parallel tool execution; streaming; 50 KB/2000-line tool output caps; lazy provider module loading & \texttt{ai/\allowbreak{}api/\allowbreak{}anthropic-\allowbreak{}messages.\allowbreak{}ts:\allowbreak{}L1529-\allowbreak{}-\allowbreak{}L1555}; \texttt{ca/\allowbreak{}core/\allowbreak{}cache-\allowbreak{}warmer.\allowbreak{}ts:\allowbreak{}L15-\allowbreak{}-\allowbreak{}L32}, \texttt{L284-\allowbreak{}-\allowbreak{}L296}; \texttt{agent/\allowbreak{}agent-\allowbreak{}loop.\allowbreak{}ts:\allowbreak{}L586-\allowbreak{}-\allowbreak{}L660}; \texttt{ca/\allowbreak{}core/\allowbreak{}tools/\allowbreak{}truncate.\allowbreak{}ts:\allowbreak{}L11-\allowbreak{}-\allowbreak{}L12}; \texttt{ai/\allowbreak{}api/\allowbreak{}lazy.\allowbreak{}ts:\allowbreak{}L74-\allowbreak{}-\allowbreak{}L80} \\
Evaluation & (evaluation) & 651 \texttt{*.\allowbreak{}test.\allowbreak{}ts} files (vitest) across packages (coding-agent 207, ai 180, agent 5), incl. agent-session retry/compaction/branching suites; \texttt{packages/\allowbreak{}evals} runs the real agent via \texttt{vitest-\allowbreak{}evals} harness with Docker variants and LLM-graded docs evals (\texttt{autoevals}) & \texttt{packages/\allowbreak{}coding-\allowbreak{}agent/\allowbreak{}test/\allowbreak{}} (listing); \texttt{ev/\allowbreak{}package.\allowbreak{}json:\allowbreak{}L1-\allowbreak{}-\allowbreak{}L23}; \texttt{ev/\allowbreak{}evals/\allowbreak{}smoke.\allowbreak{}eval.\allowbreak{}ts:\allowbreak{}L5-\allowbreak{}-\allowbreak{}L18}; \texttt{ev/\allowbreak{}src/\allowbreak{}harness.\allowbreak{}ts:\allowbreak{}L1-\allowbreak{}-\allowbreak{}L50} \\
\bottomrule \end{longtable}

\subsubsection*{Design decisions (one line each: the choice this harness made within a primitive)}
\begin{itemize}
\item Input: three capture surfaces (TUI editor, print/JSON CLI, JSONL RPC) all funnel into \texttt{AgentSession.\allowbreak{}prompt(\allowbreak{})}, with mid-run input split into "steer" (after this turn) vs. "follow-up" (after the agent would stop), one message at a time by default.
\item Context: the transcript is a persisted tree; every request is a projection of the active branch (compaction entry + kept tail + context edits), the system prompt is a set of named XML sections diffed into system-message patches, and tool declarations are transcript deltas, so replaying the file reproduces the exact request.
\item Model interface: one \texttt{StreamFunction} contract across \textasciitilde{}15 provider adapters with SDK retries disabled (\texttt{maxRetries:\allowbreak{} 0}) and pi's own two-tier retry (HTTP-level with retry-after cap, assistant-level with error-text classification), overflow routed to compaction rather than retry.
\item Output: tools are plain \texttt{\{name, parameters (\allowbreak{}TypeBox), execute\}} objects; parallel execution by default with per-file mutation serialization; no sandbox in core (OS sandboxing and permission prompts are extensions); a QuickJS-sandboxed \texttt{codemode} tool lets the model script tool calls.
\item Control flow: a two-level loop in \texttt{pi-\allowbreak{}agent-\allowbreak{}core} (tool calls/steering inner, follow-ups outer) plus a session-level outer loop for retries/overflow recovery/extension continuations; sub-agents are not in core but an example spawns child \texttt{pi} processes; per-request model routing via virtual models.
\end{itemize}

\subsubsection*{Hard cases}
\begin{itemize}
\item \textbf{\texttt{before\_\allowbreak{}provider\_\allowbreak{}request} payload hook:} lives in the adapter (\texttt{onPayload}) at the wire boundary, but an extension can rewrite the request body, i.e. decide what the model is shown. By the function rule it is Context when it changes content; the transport callback itself stays Model interface. \texttt{ai/\allowbreak{}api/\allowbreak{}anthropic-\allowbreak{}messages.\allowbreak{}ts:\allowbreak{}L653-\allowbreak{}-\allowbreak{}L656}; \texttt{ca/\allowbreak{}core/\allowbreak{}sdk.\allowbreak{}ts:\allowbreak{}L382-\allowbreak{}-\allowbreak{}L386}
\item \textbf{Virtual-model routing:} \texttt{prepareRequest} chooses the physical model per request by reason (user/continuation/retry) --- keyed to the work, so Control flow --- even though it is implemented inside the model-runtime layer; the \texttt{reason:\allowbreak{} "retry"} branch after a failed request overlaps Model interface's "backup model" (resolved: the router decides based on work state, so Control flow). \texttt{ca/\allowbreak{}core/\allowbreak{}agent-\allowbreak{}session.\allowbreak{}ts:\allowbreak{}L815-\allowbreak{}-\allowbreak{}L843}; \texttt{ca/\allowbreak{}core/\allowbreak{}model-\allowbreak{}runtime.\allowbreak{}ts:\allowbreak{}L1003-\allowbreak{}-\allowbreak{}L1034}
\item \textbf{Strict-tools 400 fallback:} an in-adapter request rebuild on a specific 400 (drops strict schemas). It alters the request contents, but the trigger is "request could not be served" --- assigned Model interface by the unavailability rule. \texttt{ai/\allowbreak{}api/\allowbreak{}anthropic-\allowbreak{}messages.\allowbreak{}ts:\allowbreak{}L667-\allowbreak{}-\allowbreak{}L691}
\item \textbf{Compaction summary call:} a second model call issued by the Context primitive (summary generation). Rule: Context never calls the model; here it does, through the same \texttt{StreamFunction}. Resolved as Context mechanism using the Model interface as a service; noted as a borderline.
\item \textbf{Nested tool calls / codemode:} a script the model writes sequences only tools (Output by rule), yet it also reaches \texttt{models.\allowbreak{}*} (classifiers, image models) from inside the sandbox --- a model call issued from inside a tool. \texttt{ca/\allowbreak{}extensions/\allowbreak{}codemode/\allowbreak{}tool.\allowbreak{}ts:\allowbreak{}L1-\allowbreak{}-\allowbreak{}L10}; \texttt{ca/\allowbreak{}extensions/\allowbreak{}codemode/\allowbreak{}execute.\allowbreak{}ts:\allowbreak{}L50}. Partially DOES NOT FIT: a tool that calls models is neither pure Output nor Control flow in the given scheme.
\item \textbf{\texttt{agent\_\allowbreak{}before\_\allowbreak{}settle} / \texttt{turn\_\allowbreak{}end} boundaries:} extensions may commit "boundary drafts" (entries appended to the session: Context) and return \texttt{continue} (Control flow) in one hook; decomposed per rule 28. \texttt{ca/\allowbreak{}core/\allowbreak{}agent-\allowbreak{}session.\allowbreak{}ts:\allowbreak{}L847-\allowbreak{}-\allowbreak{}L897}, \texttt{L1983-\allowbreak{}-\allowbreak{}L2005}
\item \textbf{\texttt{pi-\allowbreak{}durable} package:} a second, experimental harness in the same repo (durable task graph, child tasks, commit-before-show) sharing \texttt{pi-\allowbreak{}ai}; not deconstructed here beyond noting it exists. \texttt{packages/\allowbreak{}durable/\allowbreak{}README.\allowbreak{}md:\allowbreak{}L1-\allowbreak{}-\allowbreak{}L6}, \texttt{L581}
\item \textbf{Telemetry package:} span interface exists and \texttt{telemetryContext} is passed through options, but no span is started in \texttt{ai/\allowbreak{}}, \texttt{agent/\allowbreak{}}, or \texttt{ca/\allowbreak{}} at this commit --- recorded as "interface present, emission not found". \texttt{packages/\allowbreak{}telemetry/\allowbreak{}src/\allowbreak{}index.\allowbreak{}ts:\allowbreak{}L1-\allowbreak{}-\allowbreak{}L22}; \texttt{ai/\allowbreak{}api/\allowbreak{}simple-\allowbreak{}options.\allowbreak{}ts:\allowbreak{}L48}
\end{itemize}

\subsubsection*{Summary row}
\begin{longtable}{@{}p{0.18\linewidth}p{0.18\linewidth}p{0.18\linewidth}p{0.18\linewidth}p{0.18\linewidth}@{}}
\toprule \textbf{Input} & \textbf{Context} & \textbf{Model interface} & \textbf{Output} & \textbf{Control flow} \\ \midrule \endhead
TUI/CLI/RPC prompts, steer/follow-up queues, tool results & Tree-projected JSONL transcript, sectioned prompt, threshold compaction & One StreamFunction, SSE, two-tier retry, overflow classifier & Parallel TypeBox tools, bash/edit/read, QuickJS codemode & Inner tool loop, outer follow-up loop, retry/overflow re-entry, routing \\
\bottomrule \end{longtable}


\subsection{Hermes Agent}

\textbf{Source:} https://github.com/NousResearch/hermes-agent $\cdot$ \textbf{Version:} commit 7bf3766f3db72d4643d1a78d18e51c1d9ddf1d6d (2026-10-15) $\cdot$ \textbf{Language:} Python (TypeScript for the TUI/desktop front-ends, out of scope)
\textbf{Path conventions:} paths are relative to the repo root (\texttt{hermes-\allowbreak{}agent/\allowbreak{}}). \texttt{L} = line. Clone read as untrusted data; nothing executed.
\textbf{What it is:} A personal agent by Nous Research that runs one agent core (\texttt{AIAgent}) behind a prompt\_toolkit CLI/TUI, \textasciitilde{}25 messaging-platform gateway adapters (Telegram, Discord, Slack, \ldots{}), an OpenAI-compatible HTTP server, an ACP adapter, a cron scheduler and a desktop app. The model is reached through an OpenAI-SDK chat-completions client by default, with alternative wire transports for Anthropic Messages, OpenAI Responses ("codex") and Bedrock Converse selected by \texttt{api\_\allowbreak{}mode}; providers are plugins and the model slug is config.

\subsubsection*{Control flow}
\begin{itemize}
\item \textbf{Turn entry + lease:} \texttt{AIAgent.\allowbreak{}run\_\allowbreak{}conversation} builds a task/turn id and admits a durable turn lease before forwarding to the loop (one turn per session at a time). \texttt{agent/\allowbreak{}turn\_\allowbreak{}facade.\allowbreak{}py:\allowbreak{}L22}, \texttt{agent/\allowbreak{}turn\_\allowbreak{}facade.\allowbreak{}py:\allowbreak{}L83}
\item \textbf{Main loop:} \texttt{\_\allowbreak{}run\_\allowbreak{}conversation\_\allowbreak{}turn} runs \texttt{while (\allowbreak{}api\_\allowbreak{}call\_\allowbreak{}count < max\_\allowbreak{}iterations and iteration\_\allowbreak{}budget.\allowbreak{}remaining > 0) or \_\allowbreak{}budget\_\allowbreak{}grace\_\allowbreak{}call}; each iteration = begin $\rightarrow$ prepare $\rightarrow$ assemble $\rightarrow$ preflight gate $\rightarrow$ API retry loop $\rightarrow$ normalize $\rightarrow$ \texttt{run\_\allowbreak{}tool\_\allowbreak{}round} if tool calls else \texttt{finish\_\allowbreak{}text\_\allowbreak{}response}; post-loop \texttt{finalize\_\allowbreak{}turn}. \texttt{agent/\allowbreak{}conversation\_\allowbreak{}loop.\allowbreak{}py:\allowbreak{}L1656-\allowbreak{}-\allowbreak{}L1706}
\item \textbf{Phase dispatch:} each phase is a sibling function returning a verdict with \texttt{action $\in$ \{fallthrough, continue, break, return\}}; \texttt{\_\allowbreak{}run\_\allowbreak{}phase} copies verdict fields back into \texttt{\_\allowbreak{}LoopState}. \texttt{agent/\allowbreak{}conversation\_\allowbreak{}loop.\allowbreak{}py:\allowbreak{}L1389-\allowbreak{}-\allowbreak{}L1502}
\item \textbf{Inner retry loop:} \texttt{\_\allowbreak{}run\_\allowbreak{}api\_\allowbreak{}retry\_\allowbreak{}loop} loops \texttt{while retry\_\allowbreak{}count < max\_\allowbreak{}retries} over guard $\rightarrow$ build $\rightarrow$ call $\rightarrow$ check, with \texttt{handle\_\allowbreak{}api\_\allowbreak{}interrupt} / \texttt{handle\_\allowbreak{}api\_\allowbreak{}error} branches. \texttt{agent/\allowbreak{}conversation\_\allowbreak{}loop.\allowbreak{}py:\allowbreak{}L1505-\allowbreak{}-\allowbreak{}L1534}
\item \textbf{Step budget:} \texttt{IterationBudget} is a thread-safe counter with \texttt{consume(\allowbreak{})}/\texttt{refund(\allowbreak{})}; \texttt{max\_\allowbreak{}iterations} defaults to \texttt{sys.\allowbreak{}maxsize} (unlimited, shared with subagents); a tool round whose only tool was \texttt{execute\_\allowbreak{}code} refunds its iteration; a reasoning-only stall that triggers a model fallback sets \texttt{\_\allowbreak{}budget\_\allowbreak{}grace\_\allowbreak{}call} for one extra call. \texttt{agent/\allowbreak{}iteration\_\allowbreak{}budget.\allowbreak{}py:\allowbreak{}L25-\allowbreak{}-\allowbreak{}L56}, \texttt{run\_\allowbreak{}agent.\allowbreak{}py:\allowbreak{}L241}, \texttt{agent/\allowbreak{}turn\_\allowbreak{}tool\_\allowbreak{}round.\allowbreak{}py:\allowbreak{}L199-\allowbreak{}-\allowbreak{}L200}, \texttt{agent/\allowbreak{}turn\_\allowbreak{}truncation.\allowbreak{}py:\allowbreak{}L630-\allowbreak{}-\allowbreak{}L635}
\item \textbf{Stop gates (go again vs. hand back):} on a text answer, \texttt{apply\_\allowbreak{}stop\_\allowbreak{}gates} runs verify-on-stop $\rightarrow$ \texttt{pre\_\allowbreak{}verify} plugin hook $\rightarrow$ kanban stop guard; a nudge appends a user-role row and continues the loop, holding the answer as a budget-exhaustion fallback. \texttt{agent/\allowbreak{}turn\_\allowbreak{}stop\_\allowbreak{}gates.\allowbreak{}py:\allowbreak{}L105-\allowbreak{}-\allowbreak{}L140}, \texttt{agent/\allowbreak{}turn\_\allowbreak{}stop\_\allowbreak{}gates.\allowbreak{}py:\allowbreak{}L35-\allowbreak{}-\allowbreak{}L50}
\item \textbf{Tool-call guardrail halt:} a \texttt{ToolCallGuardrailController} decision can end the turn (\texttt{guardrail\_\allowbreak{}halt}) after a tool round; identical-call stall threshold is 3. \texttt{agent/\allowbreak{}turn\_\allowbreak{}tool\_\allowbreak{}round.\allowbreak{}py:\allowbreak{}L179-\allowbreak{}-\allowbreak{}L193}, \texttt{agent/\allowbreak{}tool\_\allowbreak{}guardrails.\allowbreak{}py:\allowbreak{}L30}, \texttt{agent/\allowbreak{}tool\_\allowbreak{}guardrails.\allowbreak{}py:\allowbreak{}L307}
\item \textbf{Approval gate (harness consumes the answer):} \texttt{check\_\allowbreak{}all\_\allowbreak{}command\_\allowbreak{}guards} runs floor blocks $\rightarrow$ yolo/\texttt{off} bypass (process \texttt{-\allowbreak{}-\allowbreak{}yolo}, session \texttt{/\allowbreak{}yolo}, \texttt{approvals.\allowbreak{}mode}) $\rightarrow$ permanent allowlist $\rightarrow$ unattended-context policy $\rightarrow$ \texttt{detect\_\allowbreak{}dangerous\_\allowbreak{}command} $\rightarrow$ \texttt{\_\allowbreak{}human\_\allowbreak{}decision}; answers are \texttt{once/\allowbreak{}session/\allowbreak{}always/\allowbreak{}deny/\allowbreak{}timeout}; in gateway ask mode the tool returns \texttt{status="pending\_\allowbreak{}approval"} instead of blocking. \texttt{tools/\allowbreak{}approval.\allowbreak{}py:\allowbreak{}L1151-\allowbreak{}-\allowbreak{}L1200}, \texttt{tools/\allowbreak{}approval.\allowbreak{}py:\allowbreak{}L506-\allowbreak{}-\allowbreak{}L510}, \texttt{tools/\allowbreak{}approval.\allowbreak{}py:\allowbreak{}L796-\allowbreak{}-\allowbreak{}L802}, \texttt{tools/\allowbreak{}approval\_\allowbreak{}prompt.\allowbreak{}py:\allowbreak{}L37}, \texttt{tools/\allowbreak{}terminal\_\allowbreak{}tool.\allowbreak{}py:\allowbreak{}L1019-\allowbreak{}-\allowbreak{}L1028}
\item \textbf{Interrupt / redirect crossing a response:} if a redirect arrived while the request was in flight the response is discarded and a restart is armed; \texttt{\_\allowbreak{}apply\_\allowbreak{}active\_\allowbreak{}turn\_\allowbreak{}redirect} then appends a checkpoint + correction row so role alternation holds. \texttt{agent/\allowbreak{}turn\_\allowbreak{}api\_\allowbreak{}call.\allowbreak{}py:\allowbreak{}L148-\allowbreak{}-\allowbreak{}L167}, \texttt{agent/\allowbreak{}conversation\_\allowbreak{}loop.\allowbreak{}py:\allowbreak{}L322-\allowbreak{}-\allowbreak{}L340}
\item \textbf{Gateway steer:} a message arriving while a turn runs is injected via \texttt{agent.\allowbreak{}steer(\allowbreak{})} or queued; after a tool batch, pending steer is folded onto tool results once the per-turn budget is enforced. \texttt{gateway/\allowbreak{}run\_\allowbreak{}inbound.\allowbreak{}py:\allowbreak{}L649-\allowbreak{}-\allowbreak{}L664}, \texttt{agent/\allowbreak{}tool\_\allowbreak{}executor.\allowbreak{}py:\allowbreak{}L1209-\allowbreak{}-\allowbreak{}L1220}
\item \textbf{Delegation (sub-agents):} \texttt{delegate\_\allowbreak{}task} refuses when \texttt{\_\allowbreak{}delegate\_\allowbreak{}depth >= delegation.\allowbreak{}max\_\allowbreak{}spawn\_\allowbreak{}depth} (default 1 = flat); children are full \texttt{AIAgent} instances built with \texttt{child\_\allowbreak{}depth = parent + 1} and role orchestrator/leaf derived from depth. \texttt{tools/\allowbreak{}delegate\_\allowbreak{}tool.\allowbreak{}py:\allowbreak{}L520-\allowbreak{}-\allowbreak{}L527}, \texttt{tools/\allowbreak{}delegate\_\allowbreak{}tool\_\allowbreak{}config.\allowbreak{}py:\allowbreak{}L151-\allowbreak{}-\allowbreak{}L159}, \texttt{tools/\allowbreak{}delegate\_\allowbreak{}tool.\allowbreak{}py:\allowbreak{}L199-\allowbreak{}-\allowbreak{}L234}
\item \textbf{Child run, parallelism and result return:} \texttt{\_\allowbreak{}run\_\allowbreak{}single\_\allowbreak{}child} runs a pre-built child on a worker thread, heartbeats the parent, and returns \texttt{\{status $\in$ completed|interrupted|failed, exit\_\allowbreak{}reason\}}; a batch runs on \texttt{DaemonThreadPoolExecutor(\allowbreak{}max\_\allowbreak{}workers=batch.\allowbreak{}max\_\allowbreak{}children)} polled with short waits; \texttt{background=True} returns a dispatch handle. \texttt{tools/\allowbreak{}delegate\_\allowbreak{}tool.\allowbreak{}py:\allowbreak{}L342-\allowbreak{}-\allowbreak{}L406}, \texttt{tools/\allowbreak{}delegate\_\allowbreak{}tool\_\allowbreak{}dispatch.\allowbreak{}py:\allowbreak{}L118-\allowbreak{}-\allowbreak{}L145}, \texttt{tools/\allowbreak{}delegate\_\allowbreak{}tool\_\allowbreak{}dispatch.\allowbreak{}py:\allowbreak{}L470-\allowbreak{}-\allowbreak{}L478}
\item \textbf{Work-keyed model routing (auxiliary):} titling, compression summaries, vision etc. resolve via \texttt{\_\allowbreak{}resolve\_\allowbreak{}auto\_\allowbreak{}route} (main model $\rightarrow$ per-task fallback chain $\rightarrow$ discovery chain); titles come from one small-model call off the critical path. \texttt{agent/\allowbreak{}auxiliary\_\allowbreak{}client.\allowbreak{}py:\allowbreak{}L4659-\allowbreak{}-\allowbreak{}L4684}, \texttt{agent/\allowbreak{}title\_\allowbreak{}generator.\allowbreak{}py:\allowbreak{}L1-\allowbreak{}-\allowbreak{}L5}
\item \textbf{Mixture-of-Agents turn:} \texttt{/\allowbreak{}moa} gathers reference-model context before each iteration while the normal loop keeps tool calling. \texttt{agent/\allowbreak{}moa\_\allowbreak{}loop.\allowbreak{}py:\allowbreak{}L1-\allowbreak{}-\allowbreak{}L5}
\item \textbf{Cron turns:} the scheduler runs \texttt{agent.\allowbreak{}run\_\allowbreak{}conversation} on a worker thread under an inactivity watchdog (default 600 s) on a fresh \texttt{AIAgent(\allowbreak{}platform="cron")}. \texttt{cron/\allowbreak{}scheduler.\allowbreak{}py:\allowbreak{}L1961-\allowbreak{}-\allowbreak{}L1963}, \texttt{cron/\allowbreak{}scheduler.\allowbreak{}py:\allowbreak{}L2496-\allowbreak{}-\allowbreak{}L2525}
\item \textbf{Codex app-server handoff:} when \texttt{api\_\allowbreak{}mode == "codex\_\allowbreak{}app\_\allowbreak{}server"} the whole turn is handed to an external codex subprocess; on failure a fallback is activated and the generic loop retries the same user turn. \texttt{agent/\allowbreak{}conversation\_\allowbreak{}loop.\allowbreak{}py:\allowbreak{}L1650-\allowbreak{}-\allowbreak{}L1651}, \texttt{agent/\allowbreak{}conversation\_\allowbreak{}loop.\allowbreak{}py:\allowbreak{}L1537-\allowbreak{}-\allowbreak{}L1552}
\end{itemize}

\subsubsection*{Input}
\begin{itemize}
\item \textbf{CLI REPL:} a worker thread drains \texttt{\_\allowbreak{}pending\_\allowbreak{}input} (queue) and dispatches each input; \texttt{\_\allowbreak{}looks\_\allowbreak{}like\_\allowbreak{}slash\_\allowbreak{}command} routes to \texttt{process\_\allowbreak{}command}, else \texttt{self.\allowbreak{}chat(\allowbreak{}.\allowbreak{}.\allowbreak{}.\allowbreak{})}; \texttt{/\allowbreak{}skill-\allowbreak{}name} is rewritten into a user message. \texttt{hermes\_\allowbreak{}cli/\allowbreak{}cli\_\allowbreak{}tui\_\allowbreak{}runtime\_\allowbreak{}mixin.\allowbreak{}py:\allowbreak{}L31-\allowbreak{}-\allowbreak{}L47}, \texttt{hermes\_\allowbreak{}cli/\allowbreak{}cli\_\allowbreak{}tui\_\allowbreak{}runtime\_\allowbreak{}mixin.\allowbreak{}py:\allowbreak{}L139}, \texttt{cli.\allowbreak{}py:\allowbreak{}L710}, \texttt{cli.\allowbreak{}py:\allowbreak{}L1300}
\item \textbf{CLI chat turn:} \texttt{chat(\allowbreak{})} runs the agent on a daemon thread; input typed meanwhile goes to \texttt{\_\allowbreak{}interrupt\_\allowbreak{}queue} and is re-queued as the next turn; the call is \texttt{run\_\allowbreak{}conversation(\allowbreak{}user\_\allowbreak{}message, conversation\_\allowbreak{}history[:\allowbreak{}-\allowbreak{}1], stream\_\allowbreak{}callback, task\_\allowbreak{}id, persist\_\allowbreak{}user\_\allowbreak{}message)}. \texttt{hermes\_\allowbreak{}cli/\allowbreak{}cli\_\allowbreak{}chat\_\allowbreak{}turn\_\allowbreak{}mixin.\allowbreak{}py:\allowbreak{}L49-\allowbreak{}-\allowbreak{}L60}, \texttt{hermes\_\allowbreak{}cli/\allowbreak{}cli\_\allowbreak{}chat\_\allowbreak{}turn\_\allowbreak{}mixin.\allowbreak{}py:\allowbreak{}L115-\allowbreak{}-\allowbreak{}L118}, \texttt{hermes\_\allowbreak{}cli/\allowbreak{}cli\_\allowbreak{}chat\_\allowbreak{}turn\_\allowbreak{}mixin.\allowbreak{}py:\allowbreak{}L326-\allowbreak{}-\allowbreak{}L330}
\item \textbf{Gateway event shape:} every adapter normalizes to \texttt{MessageEvent(\allowbreak{}text, message\_\allowbreak{}type, user\_\allowbreak{}id, \ldots{})}; adapters subclass \texttt{BasePlatformAdapter}; in-tree adapters (signal, webhook, bluebubbles, weixin, yuanbao, whatsapp\_cloud, api\_server) plus plugin adapters (telegram, discord, slack, matrix, teams, email, irc, \ldots{}). \texttt{gateway/\allowbreak{}platforms/\allowbreak{}event.\allowbreak{}py:\allowbreak{}L45-\allowbreak{}-\allowbreak{}L51}, \texttt{gateway/\allowbreak{}platforms/\allowbreak{}base.\allowbreak{}py:\allowbreak{}L1876}, \texttt{gateway/\allowbreak{}platforms/\allowbreak{}} (listing), \texttt{plugins/\allowbreak{}platforms/\allowbreak{}} (listing)
\item \textbf{Gateway inbound pipeline:} \texttt{\_\allowbreak{}handle\_\allowbreak{}message} = admit $\rightarrow$ pending-reply intercepts $\rightarrow$ authority admission $\rightarrow$ running-session handling (steer/interrupt/preempt) $\rightarrow$ idle commands $\rightarrow$ \texttt{\_\allowbreak{}handle\_\allowbreak{}message\_\allowbreak{}with\_\allowbreak{}agent}. \texttt{gateway/\allowbreak{}run\_\allowbreak{}inbound.\allowbreak{}py:\allowbreak{}L1196-\allowbreak{}-\allowbreak{}L1300}
\item \textbf{Authorization and bot-loop budget at admission:} unauthorized users get a pairing code (DM) or are ignored before any session state is touched; bot-authored messages are capped by a sliding window per conversation. \texttt{gateway/\allowbreak{}run\_\allowbreak{}inbound.\allowbreak{}py:\allowbreak{}L201-\allowbreak{}-\allowbreak{}L225}, \texttt{gateway/\allowbreak{}bot\_\allowbreak{}loop\_\allowbreak{}guard.\allowbreak{}py:\allowbreak{}L1-\allowbreak{}-\allowbreak{}L5}
\item \textbf{Gateway $\rightarrow$ agent call:} the turn runner calls \texttt{agent.\allowbreak{}run\_\allowbreak{}conversation(\allowbreak{}api\_\allowbreak{}message, **kwargs)} with \texttt{persist\_\allowbreak{}user\_\allowbreak{}message}, display kind and the raw platform message id. \texttt{gateway/\allowbreak{}run\_\allowbreak{}turn\_\allowbreak{}runner.\allowbreak{}py:\allowbreak{}L934-\allowbreak{}-\allowbreak{}L936}
\item \textbf{HTTP and webhook input:} OpenAI-compatible aiohttp server (\texttt{/\allowbreak{}v1/\allowbreak{}chat/\allowbreak{}completions}, \texttt{/\allowbreak{}v1/\allowbreak{}responses}, \texttt{/\allowbreak{}v1/\allowbreak{}runs}, \ldots{}); webhook payloads are treated as untrusted and get no file/system execution toolset. \texttt{gateway/\allowbreak{}platforms/\allowbreak{}api\_\allowbreak{}server.\allowbreak{}py:\allowbreak{}L1-\allowbreak{}-\allowbreak{}L5}, \texttt{toolsets.\allowbreak{}py:\allowbreak{}L41}
\item \textbf{ACP input:} the ACP agent calls \texttt{agent.\allowbreak{}run\_\allowbreak{}conversation(\allowbreak{}user\_\allowbreak{}message=\ldots{}, conversation\_\allowbreak{}history=state.\allowbreak{}history)}. \texttt{acp\_\allowbreak{}adapter/\allowbreak{}server.\allowbreak{}py:\allowbreak{}L816-\allowbreak{}-\allowbreak{}L819}
\item \textbf{Model asks the user (clarify):} \texttt{clarify} is an inline agent tool whose UI is a platform callback; the gateway intercepts the reply and its result returns to the model. \texttt{tools/\allowbreak{}clarify\_\allowbreak{}tool.\allowbreak{}py:\allowbreak{}L1-\allowbreak{}-\allowbreak{}L6}, \texttt{tools/\allowbreak{}clarify\_\allowbreak{}tool.\allowbreak{}py:\allowbreak{}L111}, \texttt{gateway/\allowbreak{}run\_\allowbreak{}inbound.\allowbreak{}py:\allowbreak{}L339}
\item \textbf{Tool results captured:} every executed tool's result is committed as a \texttt{role:\allowbreak{} tool} row (\texttt{\_\allowbreak{}commit\_\allowbreak{}tool\_\allowbreak{}result}); invalid names get an error tool row so each call has a matching result. \texttt{agent/\allowbreak{}tool\_\allowbreak{}executor.\allowbreak{}py:\allowbreak{}L1066}, \texttt{agent/\allowbreak{}turn\_\allowbreak{}tool\_\allowbreak{}round.\allowbreak{}py:\allowbreak{}L116-\allowbreak{}-\allowbreak{}L121}
\item \textbf{Multimodal:} native image parts attached to the user turn stay visible to \texttt{vision\_\allowbreak{}analyze} for the turn. \texttt{agent/\allowbreak{}conversation\_\allowbreak{}loop.\allowbreak{}py:\allowbreak{}L1749}
\end{itemize}

\subsubsection*{Context}
\begin{itemize}
\item \textbf{User message staging:} \texttt{messages = list(\allowbreak{}conversation\_\allowbreak{}history)} + staged user dict appended; \texttt{current\_\allowbreak{}turn\_\allowbreak{}user\_\allowbreak{}idx} recorded. \texttt{agent/\allowbreak{}turn\_\allowbreak{}context.\allowbreak{}py:\allowbreak{}L1104-\allowbreak{}-\allowbreak{}L1115}
\item \textbf{System prompt as three cache tiers:} \texttt{stable} (SOUL.md from HERMES\_HOME or default identity, guidance, pinned skills, coding brief) $\rightarrow$ \texttt{context} (caller \texttt{system\_\allowbreak{}message}, project context files, workspace) $\rightarrow$ \texttt{volatile} (skills index, memory, plugin sections, profile line, timestamp, runtime environment). Never re-rendered mid-session. \texttt{agent/\allowbreak{}system\_\allowbreak{}prompt.\allowbreak{}py:\allowbreak{}L730-\allowbreak{}-\allowbreak{}L793}, \texttt{agent/\allowbreak{}system\_\allowbreak{}prompt.\allowbreak{}py:\allowbreak{}L796-\allowbreak{}-\allowbreak{}L805}, \texttt{agent/\allowbreak{}system\_\allowbreak{}prompt.\allowbreak{}py:\allowbreak{}L544-\allowbreak{}-\allowbreak{}L550}, \texttt{agent/\allowbreak{}prompt\_\allowbreak{}builder.\allowbreak{}py:\allowbreak{}L1652-\allowbreak{}-\allowbreak{}L1700}
\item \textbf{Project context files:} \texttt{.\allowbreak{}hermes.\allowbreak{}md} $\rightarrow$ \texttt{AGENTS.\allowbreak{}md} chain $\rightarrow$ \texttt{CLAUDE.\allowbreak{}md} $\rightarrow$ \texttt{.\allowbreak{}cursorrules}, loaded from cwd at startup only; cap 20 000 chars floor, dynamic 6 \% of the window up to 500 000, overridable by \texttt{context\_\allowbreak{}file\_\allowbreak{}max\_\allowbreak{}chars}. \texttt{agent/\allowbreak{}prompt\_\allowbreak{}builder.\allowbreak{}py:\allowbreak{}L1783-\allowbreak{}-\allowbreak{}L1788}, \texttt{agent/\allowbreak{}system\_\allowbreak{}prompt.\allowbreak{}py:\allowbreak{}L709-\allowbreak{}-\allowbreak{}L722}, \texttt{agent/\allowbreak{}prompt\_\allowbreak{}builder.\allowbreak{}py:\allowbreak{}L1156-\allowbreak{}-\allowbreak{}L1171}
\item \textbf{Memory files:} \texttt{MEMORY.\allowbreak{}md} (2 200 chars) and \texttt{USER.\allowbreak{}md} (1 375 chars) are rendered into the volatile tier; the \texttt{memory} tool edits them. \texttt{agent/\allowbreak{}system\_\allowbreak{}prompt.\allowbreak{}py:\allowbreak{}L516-\allowbreak{}-\allowbreak{}L541}, \texttt{tools/\allowbreak{}memory\_\allowbreak{}tool\_\allowbreak{}store.\allowbreak{}py:\allowbreak{}L186}, \texttt{tools/\allowbreak{}memory\_\allowbreak{}tool\_\allowbreak{}store.\allowbreak{}py:\allowbreak{}L618}
\item \textbf{External memory providers:} one plugin provider at a time; \texttt{build\_\allowbreak{}system\_\allowbreak{}prompt(\allowbreak{})} block plus per-turn \texttt{prefetch\_\allowbreak{}all(\allowbreak{}query)} text injected into the request. \texttt{agent/\allowbreak{}memory\_\allowbreak{}manager.\allowbreak{}py:\allowbreak{}L1-\allowbreak{}-\allowbreak{}L4}, \texttt{agent/\allowbreak{}memory\_\allowbreak{}manager.\allowbreak{}py:\allowbreak{}L501-\allowbreak{}-\allowbreak{}L511}, \texttt{agent/\allowbreak{}turn\_\allowbreak{}context.\allowbreak{}py:\allowbreak{}L903-\allowbreak{}-\allowbreak{}L931}
\item \textbf{Skills index + progressive disclosure:} names/descriptions in the prompt; \texttt{skill\_\allowbreak{}view} loads a SKILL.md on demand as a tool result. \texttt{agent/\allowbreak{}system\_\allowbreak{}prompt.\allowbreak{}py:\allowbreak{}L301-\allowbreak{}-\allowbreak{}L314}, \texttt{tools/\allowbreak{}skills\_\allowbreak{}tool.\allowbreak{}py:\allowbreak{}L2-\allowbreak{}-\allowbreak{}L5}
\item \textbf{Byte-stable prompt restored from DB:} on resume the stored \texttt{system\_\allowbreak{}prompt} row is reused if it matches the runtime; rebuilt only after compression (\texttt{invalidate\_\allowbreak{}system\_\allowbreak{}prompt}). \texttt{agent/\allowbreak{}conversation\_\allowbreak{}loop.\allowbreak{}py:\allowbreak{}L749-\allowbreak{}-\allowbreak{}L772}, \texttt{agent/\allowbreak{}system\_\allowbreak{}prompt.\allowbreak{}py:\allowbreak{}L808-\allowbreak{}-\allowbreak{}L826}, \texttt{hermes\_\allowbreak{}state\_\allowbreak{}common.\allowbreak{}py:\allowbreak{}L377-\allowbreak{}-\allowbreak{}L380}
\item \textbf{Tool definitions:} \texttt{get\_\allowbreak{}tool\_\allowbreak{}definitions(\allowbreak{}enabled/\allowbreak{}disabled toolsets)} at init; \texttt{\_\allowbreak{}HERMES\_\allowbreak{}CORE\_\allowbreak{}TOOLS} is the always-present core. \texttt{agent/\allowbreak{}agent\_\allowbreak{}init.\allowbreak{}py:\allowbreak{}L1118}, \texttt{model\_\allowbreak{}tools.\allowbreak{}py:\allowbreak{}L215}, \texttt{toolsets.\allowbreak{}py:\allowbreak{}L11-\allowbreak{}-\allowbreak{}L39}
\item \textbf{Tool Search deferral:} MCP/plugin and selected core tools are replaced in the model-visible array by \texttt{tool\_\allowbreak{}search /\allowbreak{} tool\_\allowbreak{}describe /\allowbreak{} tool\_\allowbreak{}call} bridge tools, catalog rebuilt per assembly. \texttt{tools/\allowbreak{}tool\_\allowbreak{}search.\allowbreak{}py:\allowbreak{}L1-\allowbreak{}-\allowbreak{}L7}
\item \textbf{Request assembly order:} \texttt{build\_\allowbreak{}api\_\allowbreak{}messages} $\rightarrow$ MoA context $\rightarrow$ canonicalization; "ORDER IS LOAD-BEARING" for cache breakpoints; system row prepended. \texttt{agent/\allowbreak{}turn\_\allowbreak{}request\_\allowbreak{}assembly.\allowbreak{}py:\allowbreak{}L106-\allowbreak{}-\allowbreak{}L130}, \texttt{agent/\allowbreak{}turn\_\allowbreak{}context.\allowbreak{}py:\allowbreak{}L1344}
\item \textbf{Compaction trigger:} \texttt{ContextCompressor(\allowbreak{}threshold\_\allowbreak{}percent=0.\allowbreak{}50, protect\_\allowbreak{}first\_\allowbreak{}n=3, protect\_\allowbreak{}last\_\allowbreak{}n=20)}; \texttt{should\_\allowbreak{}compress\_\allowbreak{}info} fires when \texttt{tokens >= threshold\_\allowbreak{}tokens} unless in cooldown or the ineffective breaker; per-model threshold overrides. \texttt{agent/\allowbreak{}context\_\allowbreak{}compressor.\allowbreak{}py:\allowbreak{}L2811-\allowbreak{}-\allowbreak{}L2832}, \texttt{agent/\allowbreak{}context\_\allowbreak{}compressor.\allowbreak{}py:\allowbreak{}L3006-\allowbreak{}-\allowbreak{}L3015}, \texttt{agent/\allowbreak{}context\_\allowbreak{}compressor.\allowbreak{}py:\allowbreak{}L2659-\allowbreak{}-\allowbreak{}L2671}
\item \textbf{Compaction algorithm:} prune old tool-result bodies (no LLM) $\rightarrow$ pick boundaries on the pruned copy $\rightarrow$ auxiliary-model summary $\rightarrow$ clean orphaned tool pairs. \texttt{agent/\allowbreak{}context\_\allowbreak{}compressor.\allowbreak{}py:\allowbreak{}L5528-\allowbreak{}-\allowbreak{}L5575}, \texttt{agent/\allowbreak{}context\_\allowbreak{}compressor.\allowbreak{}py:\allowbreak{}L3285-\allowbreak{}-\allowbreak{}L3300}
\item \textbf{Where compaction runs:} turn-start compaction; a pre-API preflight gate comparing the assembled request with the threshold (insufficient progress blocks further passes); 413 / context-length provider errors compress and restart. \texttt{agent/\allowbreak{}turn\_\allowbreak{}context\_\allowbreak{}compaction.\allowbreak{}py:\allowbreak{}L129}, \texttt{agent/\allowbreak{}turn\_\allowbreak{}preflight\_\allowbreak{}gate.\allowbreak{}py:\allowbreak{}L20-\allowbreak{}-\allowbreak{}L120}, \texttt{agent/\allowbreak{}turn\_\allowbreak{}overflow.\allowbreak{}py:\allowbreak{}L1-\allowbreak{}-\allowbreak{}L6}, \texttt{agent/\allowbreak{}turn\_\allowbreak{}overflow.\allowbreak{}py:\allowbreak{}L463}
\item \textbf{Gateway session hygiene:} a second compaction layer at 0.85 of the window between turns. \texttt{gateway/\allowbreak{}run\_\allowbreak{}turn\_\allowbreak{}hygiene.\allowbreak{}py:\allowbreak{}L94-\allowbreak{}-\allowbreak{}L102}
\item \textbf{Tool-output truncation:} 40 \% head / 60 \% tail with one marker; limits 50 000 bytes, 2 000 lines, 2 000 chars/line; \texttt{read\_\allowbreak{}file} pages by char budget and reports \texttt{next\_\allowbreak{}offset}. \texttt{tools/\allowbreak{}tool\_\allowbreak{}output\_\allowbreak{}truncate.\allowbreak{}py:\allowbreak{}L1-\allowbreak{}-\allowbreak{}L12}, \texttt{tools/\allowbreak{}tool\_\allowbreak{}output\_\allowbreak{}limits.\allowbreak{}py:\allowbreak{}L13-\allowbreak{}-\allowbreak{}L15}, \texttt{tools/\allowbreak{}file\_\allowbreak{}tools.\allowbreak{}py:\allowbreak{}L96-\allowbreak{}-\allowbreak{}L113}
\item \textbf{Spill-to-file instead of truncating:} results over the per-tool threshold are written to \texttt{\$HERMES\_\allowbreak{}HOME/\allowbreak{}cache/\allowbreak{}spillover/\allowbreak{}\{id\}.\allowbreak{}txt} and replaced by preview + path; \texttt{enforce\_\allowbreak{}turn\_\allowbreak{}budget} caps the batch. \texttt{tools/\allowbreak{}tool\_\allowbreak{}result\_\allowbreak{}storage.\allowbreak{}py:\allowbreak{}L1-\allowbreak{}-\allowbreak{}L6}, \texttt{tools/\allowbreak{}tool\_\allowbreak{}result\_\allowbreak{}storage.\allowbreak{}py:\allowbreak{}L293-\allowbreak{}-\allowbreak{}L300}, \texttt{tools/\allowbreak{}tool\_\allowbreak{}result\_\allowbreak{}storage.\allowbreak{}py:\allowbreak{}L359}
\item \textbf{Subdirectory hints:} AGENTS.md/CLAUDE.md found while navigating are appended to the tool result, never to the system prompt. \texttt{agent/\allowbreak{}subdirectory\_\allowbreak{}hints.\allowbreak{}py:\allowbreak{}L1-\allowbreak{}-\allowbreak{}L5}
\item \textbf{Untrusted-result wrapping:} web/browser/MCP results $\geq$ 32 chars are wrapped in \texttt{<untrusted\_\allowbreak{}tool\_\allowbreak{}result>} with delimiter neutralization. \texttt{agent/\allowbreak{}tool\_\allowbreak{}dispatch\_\allowbreak{}helpers.\allowbreak{}py:\allowbreak{}L484-\allowbreak{}-\allowbreak{}L495}, \texttt{agent/\allowbreak{}tool\_\allowbreak{}dispatch\_\allowbreak{}helpers.\allowbreak{}py:\allowbreak{}L572-\allowbreak{}-\allowbreak{}L595}
\item \textbf{History source on resume:} CLI loads \texttt{get\_\allowbreak{}messages\_\allowbreak{}as\_\allowbreak{}conversation(\allowbreak{}session\_\allowbreak{}id, repair\_\allowbreak{}alternation=True)}; gateway loads \texttt{load\_\allowbreak{}transcript} then hygiene. \texttt{hermes\_\allowbreak{}cli/\allowbreak{}cli\_\allowbreak{}agent\_\allowbreak{}setup\_\allowbreak{}mixin.\allowbreak{}py:\allowbreak{}L593-\allowbreak{}-\allowbreak{}L596}, \texttt{hermes\_\allowbreak{}state\_\allowbreak{}messages.\allowbreak{}py:\allowbreak{}L1537-\allowbreak{}-\allowbreak{}L1545}, \texttt{gateway/\allowbreak{}run\_\allowbreak{}turn\_\allowbreak{}prepare.\allowbreak{}py:\allowbreak{}L693-\allowbreak{}-\allowbreak{}L700}
\item \textbf{Pluggable context engine:} \texttt{context.\allowbreak{}engine} selects a plugin that owns \texttt{should\_\allowbreak{}compress/\allowbreak{}compress} and may expose tools. \texttt{agent/\allowbreak{}context\_\allowbreak{}engine.\allowbreak{}py:\allowbreak{}L1-\allowbreak{}-\allowbreak{}L7}
\end{itemize}

\subsubsection*{Model interface}
\begin{itemize}
\item \textbf{Transport registry:} \texttt{anthropic}, \texttt{codex}, \texttt{chat\_\allowbreak{}completions}, \texttt{bedrock} transports behind \texttt{ProviderTransport} (convert\_messages/convert\_tools/build\_kwargs/normalize\_response/validate\_response). \texttt{agent/\allowbreak{}transports/\allowbreak{}\_\allowbreak{}\_\allowbreak{}init\_\allowbreak{}\_\allowbreak{}.\allowbreak{}py:\allowbreak{}L18-\allowbreak{}-\allowbreak{}L31}, \texttt{agent/\allowbreak{}transports/\allowbreak{}base.\allowbreak{}py:\allowbreak{}L12-\allowbreak{}-\allowbreak{}L51}
\item \textbf{Client construction:} \texttt{\_\allowbreak{}build\_\allowbreak{}client} picks the Anthropic SDK, the MoA facade, boto3 Bedrock, or the OpenAI SDK client (shared httpx) by \texttt{api\_\allowbreak{}mode}; GPT-5.x on a direct OpenAI URL auto-upgrades \texttt{chat\_\allowbreak{}completions} $\rightarrow$ \texttt{codex\_\allowbreak{}responses}. \texttt{agent/\allowbreak{}agent\_\allowbreak{}init.\allowbreak{}py:\allowbreak{}L1029-\allowbreak{}-\allowbreak{}L1042}, \texttt{agent/\allowbreak{}agent\_\allowbreak{}init.\allowbreak{}py:\allowbreak{}L990-\allowbreak{}-\allowbreak{}L1026}, \texttt{agent/\allowbreak{}agent\_\allowbreak{}init.\allowbreak{}py:\allowbreak{}L482-\allowbreak{}-\allowbreak{}L499}
\item \textbf{Request kwargs:} \texttt{\_\allowbreak{}build\_\allowbreak{}api\_\allowbreak{}kwargs\_\allowbreak{}for\_\allowbreak{}mode} branches anthropic / bedrock / codex / chat-completions builders. \texttt{agent/\allowbreak{}chat\_\allowbreak{}completion\_\allowbreak{}helpers.\allowbreak{}py:\allowbreak{}L1555-\allowbreak{}-\allowbreak{}L1572}
\item \textbf{Per-attempt request build:} re-apply reasoning echo for the provider, re-render prompt-cache decoration, strip images for rejecting models, LLM request middleware (trace kept), \texttt{pre\_\allowbreak{}api\_\allowbreak{}request} hook, optional \texttt{HERMES\_\allowbreak{}DUMP\_\allowbreak{}REQUESTS} dump. \texttt{agent/\allowbreak{}turn\_\allowbreak{}api\_\allowbreak{}request.\allowbreak{}py:\allowbreak{}L93-\allowbreak{}-\allowbreak{}L205}, \texttt{agent/\allowbreak{}turn\_\allowbreak{}api\_\allowbreak{}request.\allowbreak{}py:\allowbreak{}L152-\allowbreak{}-\allowbreak{}L153}
\item \textbf{Prompt-cache decoration:} \texttt{build\_\allowbreak{}prompt\_\allowbreak{}cache\_\allowbreak{}plan} places Anthropic \texttt{cache\_\allowbreak{}control} markers; OpenAI-style history is converted to Anthropic content blocks. \texttt{agent/\allowbreak{}prompt\_\allowbreak{}caching.\allowbreak{}py:\allowbreak{}L298-\allowbreak{}-\allowbreak{}L308}, \texttt{agent/\allowbreak{}anthropic\_\allowbreak{}message\_\allowbreak{}convert.\allowbreak{}py:\allowbreak{}L1-\allowbreak{}-\allowbreak{}L4}
\item \textbf{Call issue:} \texttt{perform\_\allowbreak{}api\_\allowbreak{}call} chooses streaming (\texttt{\_\allowbreak{}should\_\allowbreak{}stream}: preferred even with no consumer, off for ACP/external-process providers) and runs through \texttt{run\_\allowbreak{}llm\_\allowbreak{}execution\_\allowbreak{}middleware}. \texttt{agent/\allowbreak{}turn\_\allowbreak{}api\_\allowbreak{}call.\allowbreak{}py:\allowbreak{}L67-\allowbreak{}-\allowbreak{}L167}, \texttt{agent/\allowbreak{}turn\_\allowbreak{}api\_\allowbreak{}call.\allowbreak{}py:\allowbreak{}L46-\allowbreak{}-\allowbreak{}L64}
\item \textbf{Streaming receive:} \texttt{\_\allowbreak{}StreamingCall} issues \texttt{chat.\allowbreak{}completions.\allowbreak{}create(\allowbreak{}**stream\_\allowbreak{}kwargs)} with \texttt{stream\_\allowbreak{}options.\allowbreak{}include\_\allowbreak{}usage}; the Anthropic path uses \texttt{messages.\allowbreak{}stream}; a cross-turn stale-stream breaker runs first. \texttt{agent/\allowbreak{}chat\_\allowbreak{}completion\_\allowbreak{}helpers.\allowbreak{}py:\allowbreak{}L3993-\allowbreak{}-\allowbreak{}L4006}, \texttt{agent/\allowbreak{}chat\_\allowbreak{}completion\_\allowbreak{}helpers.\allowbreak{}py:\allowbreak{}L2734}, \texttt{agent/\allowbreak{}chat\_\allowbreak{}completion\_\allowbreak{}helpers.\allowbreak{}py:\allowbreak{}L2957-\allowbreak{}-\allowbreak{}L2966}, \texttt{agent/\allowbreak{}chat\_\allowbreak{}completion\_\allowbreak{}helpers.\allowbreak{}py:\allowbreak{}L3412}
\item \textbf{Non-stream call on a worker thread:} per-request client so interrupts close only that connection; a stale-call detector kills and raises. \texttt{agent/\allowbreak{}chat\_\allowbreak{}completion\_\allowbreak{}helpers.\allowbreak{}py:\allowbreak{}L1292-\allowbreak{}-\allowbreak{}L1305}
\item \textbf{Retry policy:} \texttt{api\_\allowbreak{}max\_\allowbreak{}retries} default 3; wait = provider \texttt{Retry-\allowbreak{}After} (capped) else \texttt{jittered\_\allowbreak{}backoff(\allowbreak{}retry, base 2.\allowbreak{}0, max 60)} = min(base$\cdot$2\textasciicircum{}(n-1), max) + jitter; adaptive policy for 429. \texttt{agent/\allowbreak{}agent\_\allowbreak{}init.\allowbreak{}py:\allowbreak{}L1441-\allowbreak{}-\allowbreak{}L1446}, \texttt{agent/\allowbreak{}turn\_\allowbreak{}recovery.\allowbreak{}py:\allowbreak{}L1401-\allowbreak{}-\allowbreak{}L1430}, \texttt{agent/\allowbreak{}retry\_\allowbreak{}utils.\allowbreak{}py:\allowbreak{}L121-\allowbreak{}-\allowbreak{}L130}, \texttt{agent/\allowbreak{}turn\_\allowbreak{}api\_\allowbreak{}error.\allowbreak{}py:\allowbreak{}L406-\allowbreak{}-\allowbreak{}L417}
\item \textbf{Error classification:} \texttt{classify\_\allowbreak{}api\_\allowbreak{}error} $\rightarrow$ structured \texttt{ClassifiedError} (retryable / should\_compress / rotate credential / fallback); plugins may override via \texttt{transform\_\allowbreak{}api\_\allowbreak{}error\_\allowbreak{}classification}. \texttt{agent/\allowbreak{}error\_\allowbreak{}classifier.\allowbreak{}py:\allowbreak{}L974-\allowbreak{}-\allowbreak{}L980}, \texttt{hermes\_\allowbreak{}cli/\allowbreak{}plugins.\allowbreak{}py:\allowbreak{}L126-\allowbreak{}-\allowbreak{}L130}
\item \textbf{Backup model when unservable:} \texttt{try\_\allowbreak{}activate\_\allowbreak{}fallback} walks \texttt{\_\allowbreak{}fallback\_\allowbreak{}chain}, swapping client/model/provider in place, on classified failover or after retries exhaust; a persistent credential pool gives same-provider failover. \texttt{agent/\allowbreak{}chat\_\allowbreak{}completion\_\allowbreak{}helpers.\allowbreak{}py:\allowbreak{}L2086-\allowbreak{}-\allowbreak{}L2121}, \texttt{agent/\allowbreak{}turn\_\allowbreak{}api\_\allowbreak{}error.\allowbreak{}py:\allowbreak{}L341}, \texttt{agent/\allowbreak{}turn\_\allowbreak{}api\_\allowbreak{}error.\allowbreak{}py:\allowbreak{}L380}, \texttt{agent/\allowbreak{}credential\_\allowbreak{}pool.\allowbreak{}py:\allowbreak{}L1}
\item \textbf{Post-exhaustion auto-recovery:} when retries and the chain are spent on a transient outage, the turn parks with a countdown and re-enters the retry loop (\texttt{auto\_\allowbreak{}recovery\_\allowbreak{}cycles} default 5). \texttt{agent/\allowbreak{}turn\_\allowbreak{}recovery\_\allowbreak{}autorecover.\allowbreak{}py:\allowbreak{}L1-\allowbreak{}-\allowbreak{}L7}
\item \textbf{Response validation + normalization:} \texttt{validate\_\allowbreak{}response\_\allowbreak{}shape} via the transport; \texttt{normalize\_\allowbreak{}model\_\allowbreak{}response} yields an assistant message with string content. \texttt{agent/\allowbreak{}turn\_\allowbreak{}recovery.\allowbreak{}py:\allowbreak{}L1524-\allowbreak{}-\allowbreak{}L1530}, \texttt{agent/\allowbreak{}turn\_\allowbreak{}response\_\allowbreak{}intake.\allowbreak{}py:\allowbreak{}L120-\allowbreak{}-\allowbreak{}L128}
\item \textbf{Degenerate-response ladders:} empty/thinking-only $\rightarrow$ partial-stream recovery $\rightarrow$ nudge $\rightarrow$ prefill continuation $\times$2 $\rightarrow$ budgeted retries $\rightarrow$ fallback provider; \texttt{finish\_\allowbreak{}reason=length} $\rightarrow$ continuation nudges (up to 4), truncated tool-call retries with max\_tokens boosts, repetition abort. \texttt{agent/\allowbreak{}turn\_\allowbreak{}empty\_\allowbreak{}response.\allowbreak{}py:\allowbreak{}L1-\allowbreak{}-\allowbreak{}L5}, \texttt{agent/\allowbreak{}turn\_\allowbreak{}truncation.\allowbreak{}py:\allowbreak{}L1-\allowbreak{}-\allowbreak{}L6}, \texttt{agent/\allowbreak{}repetition\_\allowbreak{}guard.\allowbreak{}py:\allowbreak{}L141}
\item \textbf{Usage receipt:} \texttt{record\_\allowbreak{}response\_\allowbreak{}usage} folds \texttt{response.\allowbreak{}usage} into the compressor, session counters and state.db. \texttt{agent/\allowbreak{}turn\_\allowbreak{}usage.\allowbreak{}py:\allowbreak{}L73-\allowbreak{}-\allowbreak{}L79}
\end{itemize}

\subsubsection*{Output}
\begin{itemize}
\item \textbf{Tool set:} core tools (terminal, process\_manage, read/write/patch/search files, web, vision, image, skills, browser\_\emph{, todo, memory, session\_search, clarify, execute\_code, delegate\_task, cronjob, kanban\_}, computer\_use). \texttt{toolsets.\allowbreak{}py:\allowbreak{}L11-\allowbreak{}-\allowbreak{}L39}
\item \textbf{Dispatch precedence:} inline agent-level tools (\texttt{INLINE\_\allowbreak{}TOOL\_\allowbreak{}EXECUTORS} table) $\rightarrow$ \texttt{delegate\_\allowbreak{}task} $\rightarrow$ context-engine tools $\rightarrow$ memory-provider tools $\rightarrow$ registry via \texttt{handle\_\allowbreak{}function\_\allowbreak{}call}. \texttt{agent/\allowbreak{}tool\_\allowbreak{}executor.\allowbreak{}py:\allowbreak{}L1669-\allowbreak{}-\allowbreak{}L1733}, \texttt{agent/\allowbreak{}inline\_\allowbreak{}tool\_\allowbreak{}executors.\allowbreak{}py:\allowbreak{}L253-\allowbreak{}-\allowbreak{}L313}
\item \textbf{Batch planning:} >1 calls are split into parallel/sequential segments; the parallel-safe set is read-only tools plus path-non-overlapping file ops. \texttt{run\_\allowbreak{}agent.\allowbreak{}py:\allowbreak{}L1299-\allowbreak{}-\allowbreak{}L1325}, \texttt{agent/\allowbreak{}tool\_\allowbreak{}dispatch\_\allowbreak{}helpers.\allowbreak{}py:\allowbreak{}L212-\allowbreak{}-\allowbreak{}L257}, \texttt{agent/\allowbreak{}tool\_\allowbreak{}dispatch\_\allowbreak{}helpers.\allowbreak{}py:\allowbreak{}L33-\allowbreak{}-\allowbreak{}L44}
\item \textbf{Concurrent execution:} \texttt{\_\allowbreak{}ConcurrentBatch} runs on \texttt{DaemonThreadPoolExecutor(\allowbreak{}max\_\allowbreak{}workers)} with a batch deadline; results appended in call order. \texttt{agent/\allowbreak{}tool\_\allowbreak{}executor.\allowbreak{}py:\allowbreak{}L1573-\allowbreak{}-\allowbreak{}L1620}, \texttt{agent/\allowbreak{}tool\_\allowbreak{}executor.\allowbreak{}py:\allowbreak{}L1500-\allowbreak{}-\allowbreak{}L1501}
\item \textbf{Registry call path:} \texttt{handle\_\allowbreak{}function\_\allowbreak{}call} $\rightarrow$ aliases $\rightarrow$ bridge unwrap $\rightarrow$ request middleware $\rightarrow$ \texttt{\_\allowbreak{}pre\_\allowbreak{}dispatch\_\allowbreak{}guards} (\texttt{pre\_\allowbreak{}tool\_\allowbreak{}call} hook, fail-closed on timeout) $\rightarrow$ \texttt{registry.\allowbreak{}dispatch} $\rightarrow$ \texttt{post\_\allowbreak{}tool\_\allowbreak{}call} $\rightarrow$ \texttt{transform\_\allowbreak{}tool\_\allowbreak{}result}. \texttt{model\_\allowbreak{}tools.\allowbreak{}py:\allowbreak{}L873-\allowbreak{}-\allowbreak{}L967}, \texttt{model\_\allowbreak{}tools.\allowbreak{}py:\allowbreak{}L823-\allowbreak{}-\allowbreak{}L846}, \texttt{hermes\_\allowbreak{}cli/\allowbreak{}plugins\_\allowbreak{}dispatch.\allowbreak{}py:\allowbreak{}L51}
\item \textbf{Terminal tool:} plan $\rightarrow$ acquire env $\rightarrow$ bounded pre-exec guard $\rightarrow$ approval guards $\rightarrow$ foreground or background spawn; description tells the model output is auto-truncated to a file. \texttt{tools/\allowbreak{}terminal\_\allowbreak{}tool.\allowbreak{}py:\allowbreak{}L1360-\allowbreak{}-\allowbreak{}L1475}, \texttt{tools/\allowbreak{}terminal\_\allowbreak{}tool.\allowbreak{}py:\allowbreak{}L160-\allowbreak{}-\allowbreak{}L165}
\item \textbf{Execution backends (sandboxes):} local, docker, singularity/apptainer, modal, daytona, vercel\_sandbox, ssh, plus plugin envs; built by \texttt{\_\allowbreak{}ENV\_\allowbreak{}BUILDERS}. \texttt{tools/\allowbreak{}terminal\_\allowbreak{}tool\_\allowbreak{}backends.\allowbreak{}py:\allowbreak{}L54}, \texttt{tools/\allowbreak{}terminal\_\allowbreak{}tool\_\allowbreak{}backends.\allowbreak{}py:\allowbreak{}L239-\allowbreak{}-\allowbreak{}L241}
\item \textbf{Programmatic tool calling:} \texttt{execute\_\allowbreak{}code} runs a model-written Python script in a persistent per-conversation kernel that calls Hermes tools over RPC; only stdout returns. \texttt{tools/\allowbreak{}code\_\allowbreak{}execution\_\allowbreak{}tool.\allowbreak{}py:\allowbreak{}L1-\allowbreak{}-\allowbreak{}L10}, \texttt{tools/\allowbreak{}code\_\allowbreak{}kernel.\allowbreak{}py:\allowbreak{}L1-\allowbreak{}-\allowbreak{}L5}
\item \textbf{Browser:} agent-browser CLI over local headless Chromium, cloud backends, a user CDP endpoint or Camofox; accessibility-tree snapshots with \texttt{@eN} refs. \texttt{tools/\allowbreak{}browser\_\allowbreak{}tool.\allowbreak{}py:\allowbreak{}L1-\allowbreak{}-\allowbreak{}L8}
\item \textbf{Persist-before-execute:} the assistant tool-call row is flushed to SQLite before any tool runs; a failed append ends the turn. \texttt{agent/\allowbreak{}turn\_\allowbreak{}tool\_\allowbreak{}round.\allowbreak{}py:\allowbreak{}L131-\allowbreak{}-\allowbreak{}L153}
\item \textbf{Rendering and delivery:} single-writer stream ownership with delta/reasoning fan-out; open stream boxes flushed before tools; gateway \texttt{delivery.\allowbreak{}py} routes responses to targets and \texttt{response\_\allowbreak{}filters.\allowbreak{}py} decides whether a turn is delivered to the chat. \texttt{agent/\allowbreak{}stream\_\allowbreak{}delivery.\allowbreak{}py:\allowbreak{}L1-\allowbreak{}-\allowbreak{}L3}, \texttt{agent/\allowbreak{}turn\_\allowbreak{}tool\_\allowbreak{}round.\allowbreak{}py:\allowbreak{}L160-\allowbreak{}-\allowbreak{}L164}, \texttt{gateway/\allowbreak{}delivery.\allowbreak{}py:\allowbreak{}L1-\allowbreak{}-\allowbreak{}L2}, \texttt{gateway/\allowbreak{}response\_\allowbreak{}filters.\allowbreak{}py:\allowbreak{}L1-\allowbreak{}-\allowbreak{}L4}
\item \textbf{Sub-agent spawn tool:} \texttt{delegate\_\allowbreak{}task} (output) whose effect is the control-flow delegation above. \texttt{tools/\allowbreak{}delegate\_\allowbreak{}tool.\allowbreak{}py:\allowbreak{}L488}
\end{itemize}

\subsubsection*{Production concerns}
\begin{longtable}{@{}p{0.13\linewidth}p{0.17\linewidth}p{0.36\linewidth}p{0.26\linewidth}@{}}
\toprule \textbf{Concern} & \textbf{Where it lives (primitive)} & \textbf{Mechanism} & \textbf{Evidence} \\ \midrule \endhead
Persistence & Output (persist-before-execute), Context (prompt/history restore), Input (session resume) & SQLite \texttt{state.\allowbreak{}db} in WAL mode: \texttt{sessions}, \texttt{messages} (role/content/tool\_calls/reasoning/\ldots{}), \texttt{system\_\allowbreak{}prompts}, \texttt{session\_\allowbreak{}model\_\allowbreak{}usage}, FTS5 search; \texttt{\_\allowbreak{}flush\_\allowbreak{}messages\_\allowbreak{}to\_\allowbreak{}session\_\allowbreak{}db} dedups by marker; memory in \texttt{MEMORY.\allowbreak{}md}/\texttt{USER.\allowbreak{}md}; spillover files; trajectories & \texttt{hermes\_\allowbreak{}state.\allowbreak{}py:\allowbreak{}L177}, \texttt{hermes\_\allowbreak{}state.\allowbreak{}py:\allowbreak{}L460-\allowbreak{}-\allowbreak{}L466}, \texttt{hermes\_\allowbreak{}state\_\allowbreak{}common.\allowbreak{}py:\allowbreak{}L451-\allowbreak{}-\allowbreak{}L470}, \texttt{agent/\allowbreak{}session\_\allowbreak{}persistence.\allowbreak{}py:\allowbreak{}L467-\allowbreak{}-\allowbreak{}L511}, \texttt{agent/\allowbreak{}turn\_\allowbreak{}finalizer.\allowbreak{}py:\allowbreak{}L608}, \texttt{agent/\allowbreak{}turn\_\allowbreak{}finalizer.\allowbreak{}py:\allowbreak{}L635} \\
Safety \& security & Output (sandbox, approvals, path denies), Context (injection scans, untrusted wrapping), Control flow (gate decision) & OS is the only boundary; hardline + dangerous command regexes with human gate; write-denied credential paths (\texttt{.\allowbreak{}ssh}, \texttt{.\allowbreak{}env}, \ldots{}); SSRF-safe URL policy; secret redaction of logs/tool output; threat-pattern scan of context files, skills and tool results; \texttt{UnscopedSecretError} secret scoping & \texttt{SECURITY.\allowbreak{}md:\allowbreak{}L58-\allowbreak{}-\allowbreak{}L65}, \texttt{tools/\allowbreak{}approval\_\allowbreak{}detection.\allowbreak{}py:\allowbreak{}L107}, \texttt{tools/\allowbreak{}approval\_\allowbreak{}detection.\allowbreak{}py:\allowbreak{}L228}, \texttt{agent/\allowbreak{}file\_\allowbreak{}safety.\allowbreak{}py:\allowbreak{}L176-\allowbreak{}-\allowbreak{}L196}, \texttt{agent/\allowbreak{}file\_\allowbreak{}safety.\allowbreak{}py:\allowbreak{}L210}, \texttt{tools/\allowbreak{}url\_\allowbreak{}safety.\allowbreak{}py:\allowbreak{}L1-\allowbreak{}-\allowbreak{}L9}, \texttt{agent/\allowbreak{}redact.\allowbreak{}py:\allowbreak{}L1-\allowbreak{}-\allowbreak{}L3}, \texttt{agent/\allowbreak{}prompt\_\allowbreak{}builder.\allowbreak{}py:\allowbreak{}L83-\allowbreak{}-\allowbreak{}L110}, \texttt{tools/\allowbreak{}skills\_\allowbreak{}guard.\allowbreak{}py:\allowbreak{}L240-\allowbreak{}-\allowbreak{}L253}, \texttt{agent/\allowbreak{}tool\_\allowbreak{}dispatch\_\allowbreak{}helpers.\allowbreak{}py:\allowbreak{}L544-\allowbreak{}-\allowbreak{}L562} \\
Observability & across (recorded at Model interface and Output) & per-call usage folded into state.db; \texttt{pre/\allowbreak{}post\_\allowbreak{}api\_\allowbreak{}request}, \texttt{pre/\allowbreak{}post\_\allowbreak{}tool\_\allowbreak{}call}, \texttt{pre/\allowbreak{}post\_\allowbreak{}auxiliary\_\allowbreak{}call} plugin hooks; gateway monitoring events exported over OTLP (opt-in); \texttt{HERMES\_\allowbreak{}DUMP\_\allowbreak{}REQUESTS}; middleware traces on requests & \texttt{agent/\allowbreak{}turn\_\allowbreak{}usage.\allowbreak{}py:\allowbreak{}L73-\allowbreak{}-\allowbreak{}L79}, \texttt{hermes\_\allowbreak{}state\_\allowbreak{}common.\allowbreak{}py:\allowbreak{}L484}, \texttt{hermes\_\allowbreak{}cli/\allowbreak{}plugins.\allowbreak{}py:\allowbreak{}L106-\allowbreak{}-\allowbreak{}L135}, \texttt{agent/\allowbreak{}monitoring/\allowbreak{}events.\allowbreak{}py:\allowbreak{}L16-\allowbreak{}-\allowbreak{}L66}, \texttt{agent/\allowbreak{}monitoring/\allowbreak{}otlp\_\allowbreak{}exporter.\allowbreak{}py:\allowbreak{}L1-\allowbreak{}-\allowbreak{}L5}, \texttt{agent/\allowbreak{}turn\_\allowbreak{}api\_\allowbreak{}request.\allowbreak{}py:\allowbreak{}L152-\allowbreak{}-\allowbreak{}L153} \\
Resilience & Model interface (retries, fallback, auto-recovery), Context (compaction on overflow), Output (bounded guards) & 3 retries with Retry-After/jittered backoff; fallback chain and credential pool; post-exhaustion ladder; 413/context-length $\rightarrow$ compress+restart; stale-stream breaker; turn liveness watchdog; cron inactivity watchdog; resume after restart from persisted tool-call rows & \texttt{agent/\allowbreak{}turn\_\allowbreak{}recovery.\allowbreak{}py:\allowbreak{}L1401-\allowbreak{}-\allowbreak{}L1430}, \texttt{agent/\allowbreak{}chat\_\allowbreak{}completion\_\allowbreak{}helpers.\allowbreak{}py:\allowbreak{}L2086-\allowbreak{}-\allowbreak{}L2121}, \texttt{agent/\allowbreak{}turn\_\allowbreak{}recovery\_\allowbreak{}autorecover.\allowbreak{}py:\allowbreak{}L1-\allowbreak{}-\allowbreak{}L7}, \texttt{agent/\allowbreak{}turn\_\allowbreak{}overflow.\allowbreak{}py:\allowbreak{}L1-\allowbreak{}-\allowbreak{}L6}, \texttt{agent/\allowbreak{}chat\_\allowbreak{}completion\_\allowbreak{}helpers.\allowbreak{}py:\allowbreak{}L4005}, \texttt{agent/\allowbreak{}turn\_\allowbreak{}liveness.\allowbreak{}py:\allowbreak{}L90}, \texttt{cron/\allowbreak{}scheduler.\allowbreak{}py:\allowbreak{}L1961-\allowbreak{}-\allowbreak{}L1963}, \texttt{agent/\allowbreak{}turn\_\allowbreak{}tool\_\allowbreak{}round.\allowbreak{}py:\allowbreak{}L131-\allowbreak{}-\allowbreak{}L153} \\
Performance & Context (cache tiers, truncation, spill), Output (parallel tools), Model interface (streaming, cache\_control) & system prompt byte-stable in stable$\rightarrow$context$\rightarrow$volatile tiers; compression is the only sanctioned cache break; head/tail truncation and spill files; parallel-safe tool segments; streaming by default; Anthropic cache markers; tool-definition cache & \texttt{agent/\allowbreak{}system\_\allowbreak{}prompt.\allowbreak{}py:\allowbreak{}L730-\allowbreak{}-\allowbreak{}L793}, \texttt{agent/\allowbreak{}system\_\allowbreak{}prompt.\allowbreak{}py:\allowbreak{}L808-\allowbreak{}-\allowbreak{}L826}, \texttt{tools/\allowbreak{}tool\_\allowbreak{}output\_\allowbreak{}truncate.\allowbreak{}py:\allowbreak{}L1-\allowbreak{}-\allowbreak{}L12}, \texttt{agent/\allowbreak{}tool\_\allowbreak{}dispatch\_\allowbreak{}helpers.\allowbreak{}py:\allowbreak{}L212-\allowbreak{}-\allowbreak{}L257}, \texttt{agent/\allowbreak{}turn\_\allowbreak{}api\_\allowbreak{}call.\allowbreak{}py:\allowbreak{}L46-\allowbreak{}-\allowbreak{}L64}, \texttt{agent/\allowbreak{}prompt\_\allowbreak{}caching.\allowbreak{}py:\allowbreak{}L298-\allowbreak{}-\allowbreak{}L308}, \texttt{model\_\allowbreak{}tools.\allowbreak{}py:\allowbreak{}L215} \\
Evaluation & across & 6 392 pytest files run per-file-isolated by \texttt{scripts/\allowbreak{}run\_\allowbreak{}tests.\allowbreak{}sh}; markers for \texttt{integration}/\texttt{live} provider canaries; \texttt{evals/\allowbreak{}} harnesses (compaction recall eval, provider\_fallback, tool\_search, token\_accounting, \ldots{}); \texttt{batch\_\allowbreak{}runner.\allowbreak{}py} runs the agent over JSONL datasets; \texttt{mini\_\allowbreak{}swe\_\allowbreak{}runner.\allowbreak{}py} runs SWE tasks in the execution environments & \texttt{scripts/\allowbreak{}run\_\allowbreak{}tests.\allowbreak{}sh:\allowbreak{}L1-\allowbreak{}-\allowbreak{}L15}, \texttt{pyproject.\allowbreak{}toml:\allowbreak{}L691-\allowbreak{}-\allowbreak{}L700}, \texttt{evals/\allowbreak{}compaction/\allowbreak{}README.\allowbreak{}md:\allowbreak{}L1-\allowbreak{}-\allowbreak{}L12}, \texttt{batch\_\allowbreak{}runner.\allowbreak{}py:\allowbreak{}L1-\allowbreak{}-\allowbreak{}L6}, \texttt{mini\_\allowbreak{}swe\_\allowbreak{}runner.\allowbreak{}py:\allowbreak{}L1-\allowbreak{}-\allowbreak{}L6}, \texttt{tests/\allowbreak{}evals/\allowbreak{}} (listing) \\
\bottomrule \end{longtable}

\subsubsection*{Design decisions (one line each: the choice this harness made within a primitive)}
\begin{itemize}
\item Input: one normalized \texttt{MessageEvent}/\texttt{run\_\allowbreak{}conversation(\allowbreak{}user\_\allowbreak{}message, conversation\_\allowbreak{}history, task\_\allowbreak{}id)} seam shared by CLI, \textasciitilde{}25 gateway adapters, HTTP, ACP and cron; mid-turn messages become steer/redirect rather than new turns.
\item Context: a three-tier byte-stable system prompt (identity $\rightarrow$ project files $\rightarrow$ volatile memory/skills/time) rebuilt only after compression; compaction at 50 \% of the window by prune-then-summarize with a separate auxiliary model; oversized tool output is spilled to files, not dropped.
\item Model interface: OpenAI chat-completions wire by default with pluggable transports (Anthropic, Responses, Bedrock) and provider plugins; streaming always on; 3 retries with backoff, then a configurable fallback chain, then a bounded auto-recovery wait.
\item Output: a fixed core tool set sent on every call plus a Tool Search bridge for the long tail; tool batches split into parallel-safe segments; terminal execution behind pluggable sandboxes (local/docker/modal/\ldots{}) and a regex-plus-human approval gate; tool-call rows persisted before execution.
\item Control flow: synchronous phase-verdict loop with an unlimited default iteration budget, verify-on-stop nudges that re-enter the loop, flat delegation by default (\texttt{max\_\allowbreak{}spawn\_\allowbreak{}depth=1}) with threaded children, and task-keyed auxiliary model routing.
\end{itemize}

\subsubsection*{Hard cases}
\begin{itemize}
\item \textbf{Tool Search bridge (\texttt{tool\_\allowbreak{}search}/\texttt{tool\_\allowbreak{}describe}/\texttt{tool\_\allowbreak{}call}):} a tool call the model makes (output) whose only job is to decide which tool schemas the model sees (context); resolved by "assign by function": the catalog read is context, the unwrapped \texttt{tool\_\allowbreak{}call} re-dispatch is output. \texttt{tools/\allowbreak{}tool\_\allowbreak{}search.\allowbreak{}py:\allowbreak{}L1-\allowbreak{}-\allowbreak{}L7}, \texttt{model\_\allowbreak{}tools.\allowbreak{}py:\allowbreak{}L701-\allowbreak{}-\allowbreak{}L709}
\item \textbf{\texttt{execute\_\allowbreak{}code} programmatic tool calling:} a program the model writes at run time that sequences only tools $\rightarrow$ output (rule 27); but the loop refunds the iteration it used, which is control flow. \texttt{tools/\allowbreak{}code\_\allowbreak{}execution\_\allowbreak{}tool.\allowbreak{}py:\allowbreak{}L1-\allowbreak{}-\allowbreak{}L10}, \texttt{agent/\allowbreak{}turn\_\allowbreak{}tool\_\allowbreak{}round.\allowbreak{}py:\allowbreak{}L199-\allowbreak{}-\allowbreak{}L200}
\item \textbf{Subdirectory hints and untrusted wrapping inside the tool result path:} code in the output executor that changes what the model is shown; assigned to context by rule 20. \texttt{agent/\allowbreak{}subdirectory\_\allowbreak{}hints.\allowbreak{}py:\allowbreak{}L1-\allowbreak{}-\allowbreak{}L5}, \texttt{agent/\allowbreak{}tool\_\allowbreak{}dispatch\_\allowbreak{}helpers.\allowbreak{}py:\allowbreak{}L572-\allowbreak{}-\allowbreak{}L595}
\item \textbf{Stop-gate nudges:} the harness appends a synthetic user row and goes again; the decision is control flow, the row's text is context, and it is not input (no person wrote it). \texttt{agent/\allowbreak{}turn\_\allowbreak{}stop\_\allowbreak{}gates.\allowbreak{}py:\allowbreak{}L105-\allowbreak{}-\allowbreak{}L140}
\item \textbf{Approval inside the terminal tool:} the gate runs inside the tool executor and blocks the worker thread (CLI) or returns \texttt{pending\_\allowbreak{}approval} and ends the call (gateway); the person's answer is consumed by the harness $\rightarrow$ control flow by rule 26, even though the code lives in \texttt{tools/\allowbreak{}}. \texttt{tools/\allowbreak{}terminal\_\allowbreak{}tool.\allowbreak{}py:\allowbreak{}L1013-\allowbreak{}-\allowbreak{}L1047}
\item \textbf{Reasoning-only stall $\rightarrow$ model fallback:} a model swap triggered by unusable output rather than an HTTP failure; placed in model interface because the request "could not be served" with content, not keyed to the work. \texttt{agent/\allowbreak{}turn\_\allowbreak{}truncation.\allowbreak{}py:\allowbreak{}L630-\allowbreak{}-\allowbreak{}L635}
\item \textbf{Prompt-cache decoration:} \texttt{cache\_\allowbreak{}control} markers are wire-format (model interface) but the tiered ordering of the system prompt exists only to keep them effective (context design); split across both. \texttt{agent/\allowbreak{}prompt\_\allowbreak{}caching.\allowbreak{}py:\allowbreak{}L298-\allowbreak{}-\allowbreak{}L308}, \texttt{agent/\allowbreak{}system\_\allowbreak{}prompt.\allowbreak{}py:\allowbreak{}L730-\allowbreak{}-\allowbreak{}L793}
\item \textbf{Gateway session hygiene (0.85 compaction):} compaction performed by the gateway outside the agent turn; still context (fits the window) but hosted in a different process layer. \texttt{gateway/\allowbreak{}run\_\allowbreak{}turn\_\allowbreak{}hygiene.\allowbreak{}py:\allowbreak{}L94-\allowbreak{}-\allowbreak{}L102}
\item \textbf{Codex app-server mode --- DOES NOT FIT cleanly:} with \texttt{api\_\allowbreak{}mode == "codex\_\allowbreak{}app\_\allowbreak{}server"} the whole turn (loop, tool execution, model calls) is handed to an external OpenAI Codex subprocess and Hermes only projects its events into rows and arms a fallback on failure; the five primitives collapse into one opaque external harness that is not a sub-agent (no delegation, no outer loop). \texttt{agent/\allowbreak{}conversation\_\allowbreak{}loop.\allowbreak{}py:\allowbreak{}L1650-\allowbreak{}-\allowbreak{}L1651}, \texttt{agent/\allowbreak{}conversation\_\allowbreak{}loop.\allowbreak{}py:\allowbreak{}L1537-\allowbreak{}-\allowbreak{}L1552}
\item \textbf{ACP / external-process "providers":} \texttt{acp:\allowbreak{}/\allowbreak{}/\allowbreak{}} base URLs make the "model" another agent process; streaming is disabled and the model interface boundary is another harness. \texttt{agent/\allowbreak{}turn\_\allowbreak{}api\_\allowbreak{}call.\allowbreak{}py:\allowbreak{}L46-\allowbreak{}-\allowbreak{}L56}
\item \textbf{Mixture-of-Agents:} several model calls per iteration issued by the harness (control flow) whose outputs are appended to the user message (context). \texttt{agent/\allowbreak{}moa\_\allowbreak{}loop.\allowbreak{}py:\allowbreak{}L1-\allowbreak{}-\allowbreak{}L5}
\end{itemize}

\subsubsection*{Summary row}
\begin{longtable}{@{}p{0.18\linewidth}p{0.18\linewidth}p{0.18\linewidth}p{0.18\linewidth}p{0.18\linewidth}@{}}
\toprule \textbf{Input} & \textbf{Context} & \textbf{Model interface} & \textbf{Output} & \textbf{Control flow} \\ \midrule \endhead
One MessageEvent seam: CLI, 25 gateways, HTTP, ACP, cron, steer & Tiered byte-stable prompt, memory files, 50 \% prune-then-summarize compaction & Pluggable transports, streaming, 3 retries, fallback chain, auto-recovery & Core tools + bridge, parallel segments, sandboxed terminal, approval gate & Phase-verdict loop, unlimited budget, stop-gate nudges, flat threaded delegation \\
\bottomrule \end{longtable}


\subsection{LangGraph (prebuilt ReAct agent on the Pregel runtime)}

\textbf{Source:} https://github.com/langchain-ai/langgraph $\cdot$ \textbf{Version:} commit 5eb60aef6f3837a1fb632a8abc6cf8e8da81785b (2026-10-14; packages \texttt{langgraph} 1.2.15, \texttt{langgraph-\allowbreak{}prebuilt} 1.1.1, \texttt{langgraph-\allowbreak{}checkpoint}; the local clone holds this single squashed commit, so no file history is available) $\cdot$ \textbf{Language:} Python
\textbf{Path conventions:} \texttt{LG/\allowbreak{}} = \texttt{libs/\allowbreak{}langgraph/\allowbreak{}langgraph/\allowbreak{}}; \texttt{PB/\allowbreak{}} = \texttt{libs/\allowbreak{}prebuilt/\allowbreak{}langgraph/\allowbreak{}prebuilt/\allowbreak{}}; \texttt{CP/\allowbreak{}} = \texttt{libs/\allowbreak{}checkpoint/\allowbreak{}langgraph/\allowbreak{}}; \texttt{T/\allowbreak{}} = \texttt{libs/\allowbreak{}*/\allowbreak{}tests/\allowbreak{}}. "User-written" marks a decision the framework leaves to the person building the graph.
\textbf{What it is:} LangGraph (LangChain Inc.) is a framework for writing agent control flow as a graph of nodes over a shared, checkpointed state; the Pregel runtime executes the graph in bulk-synchronous supersteps. The canonical harness deconstructed here is \texttt{create\_\allowbreak{}react\_\allowbreak{}agent} (\texttt{PB/\allowbreak{}chat\_\allowbreak{}agent\_\allowbreak{}executor.\allowbreak{}py}): an \texttt{agent} node that calls a LangChain chat model and a \texttt{tools} node (\texttt{ToolNode}) that runs the model's tool calls, cycling until the model stops calling tools. The model is reached through a \texttt{langchain\_\allowbreak{}core} \texttt{BaseChatModel} (\texttt{.\allowbreak{}invoke}/\texttt{.\allowbreak{}ainvoke}), which lives outside this repository; LangGraph itself contains no provider client.

\subsubsection*{Control flow}
\begin{itemize}
\item \textbf{ReAct loop shape (user-written graph, built by the prebuilt):} \texttt{agent} and \texttt{tools} nodes are added and the \texttt{tools} node's edge returns to the entrypoint, so the loop is an edge, not a \texttt{while}. \texttt{PB/\allowbreak{}chat\_\allowbreak{}agent\_\allowbreak{}executor.\allowbreak{}py:\allowbreak{}L867-\allowbreak{}-\allowbreak{}L872}, \texttt{L985-\allowbreak{}-\allowbreak{}L990}
\item \textbf{Go again / stop decision:} \texttt{should\_\allowbreak{}continue} reads the last message; no \texttt{tool\_\allowbreak{}calls} $\rightarrow$ \texttt{END} (or \texttt{post\_\allowbreak{}model\_\allowbreak{}hook} / \texttt{generate\_\allowbreak{}structured\_\allowbreak{}response}), otherwise route to \texttt{tools}. \texttt{PB/\allowbreak{}chat\_\allowbreak{}agent\_\allowbreak{}executor.\allowbreak{}py:\allowbreak{}L831-\allowbreak{}-\allowbreak{}L859}; attached as a conditional edge \texttt{L964-\allowbreak{}-\allowbreak{}L968}
\item \textbf{v2 fan-out, one task per tool call:} \texttt{should\_\allowbreak{}continue} returns one \texttt{Send(\allowbreak{}"tools", ToolCallWithContext(\allowbreak{}.\allowbreak{}.\allowbreak{}.\allowbreak{}))} per tool call, so each call becomes its own Pregel task. \texttt{PB/\allowbreak{}chat\_\allowbreak{}agent\_\allowbreak{}executor.\allowbreak{}py:\allowbreak{}L849-\allowbreak{}-\allowbreak{}L859}; \texttt{Send} type \texttt{LG/\allowbreak{}types.\allowbreak{}py:\allowbreak{}L738-\allowbreak{}-\allowbreak{}L826}
\item \textbf{\texttt{return\_\allowbreak{}direct} hand-back:} tools flagged \texttt{return\_\allowbreak{}direct} end the loop after they run; \texttt{route\_\allowbreak{}tool\_\allowbreak{}responses} returns \texttt{END} when the newest ToolMessages came from such a tool. \texttt{PB/\allowbreak{}chat\_\allowbreak{}agent\_\allowbreak{}executor.\allowbreak{}py:\allowbreak{}L597}, \texttt{L970-\allowbreak{}-\allowbreak{}L990}
\item \textbf{post\_model\_hook router (user-written hook, framework routing):} after the hook, pending tool calls without ToolMessages are re-sent to \texttt{tools}; a trailing ToolMessage routes back to the entrypoint; else END/structured. \texttt{PB/\allowbreak{}chat\_\allowbreak{}agent\_\allowbreak{}executor.\allowbreak{}py:\allowbreak{}L919-\allowbreak{}-\allowbreak{}L962}
\item \textbf{Step budget inside the agent (\texttt{remaining\_\allowbreak{}steps}):} a managed value computed as \texttt{scratchpad.\allowbreak{}stop -\allowbreak{} scratchpad.\allowbreak{}step}. \texttt{LG/\allowbreak{}managed/\allowbreak{}is\_\allowbreak{}last\_\allowbreak{}step.\allowbreak{}py:\allowbreak{}L18-\allowbreak{}-\allowbreak{}L24}; \texttt{stop = step + recursion\_\allowbreak{}limit + 1} \texttt{LG/\allowbreak{}pregel/\allowbreak{}\_\allowbreak{}loop.\allowbreak{}py:\allowbreak{}L1858-\allowbreak{}-\allowbreak{}L1859}; \texttt{\_\allowbreak{}are\_\allowbreak{}more\_\allowbreak{}steps\_\allowbreak{}needed} returns True when \texttt{< 2} steps remain and the response has tool calls \texttt{PB/\allowbreak{}chat\_\allowbreak{}agent\_\allowbreak{}executor.\allowbreak{}py:\allowbreak{}L620-\allowbreak{}-\allowbreak{}L634}; the model's response is then replaced by an AIMessage "Sorry, need more steps to process this request." \texttt{L684-\allowbreak{}-\allowbreak{}L694}
\item \textbf{Recursion limit (runtime step budget):} \texttt{tick(\allowbreak{})} sets \texttt{status = "out\_\allowbreak{}of\_\allowbreak{}steps"} when \texttt{step > stop}. \texttt{LG/\allowbreak{}pregel/\allowbreak{}\_\allowbreak{}loop.\allowbreak{}py:\allowbreak{}L640-\allowbreak{}-\allowbreak{}L643}; \texttt{stream(\allowbreak{})} then raises \texttt{GraphRecursionError} \texttt{LG/\allowbreak{}pregel/\allowbreak{}main.\allowbreak{}py:\allowbreak{}L3093-\allowbreak{}-\allowbreak{}L3102}; error class \texttt{LG/\allowbreak{}errors.\allowbreak{}py:\allowbreak{}L68-\allowbreak{}-\allowbreak{}L88}; default \texttt{DEFAULT\_\allowbreak{}RECURSION\_\allowbreak{}LIMIT = int(\allowbreak{}getenv(\allowbreak{}"LANGGRAPH\_\allowbreak{}DEFAULT\_\allowbreak{}RECURSION\_\allowbreak{}LIMIT", "10007"))} \texttt{LG/\allowbreak{}\_\allowbreak{}internal/\allowbreak{}\_\allowbreak{}config.\allowbreak{}py:\allowbreak{}L32}
\item \textbf{Pregel superstep loop:} \texttt{while loop.\allowbreak{}tick(\allowbreak{}):\allowbreak{} .\allowbreak{}.\allowbreak{}.\allowbreak{} runner.\allowbreak{}tick(\allowbreak{}.\allowbreak{}.\allowbreak{}.\allowbreak{}) .\allowbreak{}.\allowbreak{}.\allowbreak{} loop.\allowbreak{}after\_\allowbreak{}tick(\allowbreak{})}; channel updates from step N are visible only in step N+1. \texttt{LG/\allowbreak{}pregel/\allowbreak{}main.\allowbreak{}py:\allowbreak{}L3050-\allowbreak{}-\allowbreak{}L3079}
\item \textbf{Natural stop:} the loop ends when \texttt{prepare\_\allowbreak{}next\_\allowbreak{}tasks} yields no tasks (\texttt{status = "done"}). \texttt{LG/\allowbreak{}pregel/\allowbreak{}\_\allowbreak{}loop.\allowbreak{}py:\allowbreak{}L686-\allowbreak{}-\allowbreak{}L689}
\item \textbf{Which node runs next (the scheduler):} \texttt{prepare\_\allowbreak{}next\_\allowbreak{}tasks} first consumes pending \texttt{Send} packets from the reserved \texttt{TASKS} topic (PUSH tasks), then PULL tasks for nodes whose trigger channels were updated. \texttt{LG/\allowbreak{}pregel/\allowbreak{}\_\allowbreak{}algo.\allowbreak{}py:\allowbreak{}L391-\allowbreak{}-\allowbreak{}L512} (PUSH \texttt{L441-\allowbreak{}-\allowbreak{}L465}, PULL \texttt{L489-\allowbreak{}-\allowbreak{}L511}); a node triggers when a trigger channel's version exceeds the version it last saw \texttt{L1259-\allowbreak{}-\allowbreak{}L1276}
\item \textbf{Apply writes between steps:} \texttt{apply\_\allowbreak{}writes} sorts tasks by path for determinism, records \texttt{versions\_\allowbreak{}seen}, updates channels and bumps versions. \texttt{LG/\allowbreak{}pregel/\allowbreak{}\_\allowbreak{}algo.\allowbreak{}py:\allowbreak{}L254-\allowbreak{}-\allowbreak{}L258}, \texttt{L293-\allowbreak{}-\allowbreak{}L344}
\item \textbf{Conditional edges as routing-by-channel:} \texttt{add\_\allowbreak{}conditional\_\allowbreak{}edges} stores a \texttt{BranchSpec} \texttt{LG/\allowbreak{}graph/\allowbreak{}state.\allowbreak{}py:\allowbreak{}L982-\allowbreak{}-\allowbreak{}L1030}; at run time \texttt{\_\allowbreak{}route} reads fresh state, calls the user's path function, maps results through \texttt{path\_\allowbreak{}map}, and writes to \texttt{branch:\allowbreak{}to:\allowbreak{}<node>} channels \texttt{LG/\allowbreak{}graph/\allowbreak{}\_\allowbreak{}branch.\allowbreak{}py:\allowbreak{}L146-\allowbreak{}-\allowbreak{}L167}, \texttt{L192-\allowbreak{}-\allowbreak{}L225}; \texttt{attach\_\allowbreak{}branch} installs the branch as a writer on the source node \texttt{LG/\allowbreak{}graph/\allowbreak{}state.\allowbreak{}py:\allowbreak{}L1577-\allowbreak{}-\allowbreak{}L1624}; plain edges write the same way \texttt{L1551-\allowbreak{}-\allowbreak{}L1575}; each node is triggered by its own \texttt{branch:\allowbreak{}to:\allowbreak{}} ephemeral channel \texttt{L1525-\allowbreak{}-\allowbreak{}L1532}
\item \textbf{\texttt{Command(\allowbreak{}goto=.\allowbreak{}.\allowbreak{}.\allowbreak{})} from a node or tool:} \texttt{\_\allowbreak{}control\_\allowbreak{}branch} turns \texttt{goto} into TASKS (\texttt{Send}) or \texttt{branch:\allowbreak{}to:\allowbreak{}} writes; \texttt{Command.\allowbreak{}PARENT} raises \texttt{ParentCommand}. \texttt{LG/\allowbreak{}graph/\allowbreak{}state.\allowbreak{}py:\allowbreak{}L1749-\allowbreak{}-\allowbreak{}L1775}; \texttt{ParentCommand} is bubbled to the parent graph by \texttt{run\_\allowbreak{}with\_\allowbreak{}retry} \texttt{LG/\allowbreak{}pregel/\allowbreak{}\_\allowbreak{}retry.\allowbreak{}py:\allowbreak{}L618-\allowbreak{}-\allowbreak{}L631}; the exception class \texttt{LG/\allowbreak{}errors.\allowbreak{}py:\allowbreak{}L130-\allowbreak{}-\allowbreak{}L134}
\item \textbf{Send packets are tasks of the next (or same) superstep:} \texttt{prepare\_\allowbreak{}push\_\allowbreak{}task\_\allowbreak{}send} turns \texttt{TASKS[idx]} into a task for the named node \texttt{LG/\allowbreak{}pregel/\allowbreak{}\_\allowbreak{}algo.\allowbreak{}py:\allowbreak{}L937-\allowbreak{}-\allowbreak{}L996}; \texttt{accept\_\allowbreak{}push} schedules pushes emitted mid-step \texttt{LG/\allowbreak{}pregel/\allowbreak{}\_\allowbreak{}loop.\allowbreak{}py:\allowbreak{}L584-\allowbreak{}-\allowbreak{}L621}; \texttt{TASKS} is a non-accumulating \texttt{Topic(\allowbreak{}Send)} \texttt{LG/\allowbreak{}pregel/\allowbreak{}main.\allowbreak{}py:\allowbreak{}L814-\allowbreak{}-\allowbreak{}L819}
\item \textbf{Static breakpoints (\texttt{interrupt\_\allowbreak{}before}/\texttt{interrupt\_\allowbreak{}after}):} \texttt{should\_\allowbreak{}interrupt} returns tasks whose node is listed, if any channel changed since the last interrupt \texttt{LG/\allowbreak{}pregel/\allowbreak{}\_\allowbreak{}algo.\allowbreak{}py:\allowbreak{}L156-\allowbreak{}-\allowbreak{}L186}; checked before running tasks \texttt{LG/\allowbreak{}pregel/\allowbreak{}\_\allowbreak{}loop.\allowbreak{}py:\allowbreak{}L696-\allowbreak{}-\allowbreak{}L701} and after the step is checkpointed \texttt{L787-\allowbreak{}-\allowbreak{}L792}; the prebuilt passes the user's lists straight to \texttt{compile} \texttt{PB/\allowbreak{}chat\_\allowbreak{}agent\_\allowbreak{}executor.\allowbreak{}py:\allowbreak{}L995-\allowbreak{}-\allowbreak{}L1002}
\item \textbf{Dynamic pause (\texttt{interrupt(\allowbreak{})}):} first call in a task raises \texttt{GraphInterrupt} carrying an \texttt{Interrupt} whose id hashes the task namespace \texttt{LG/\allowbreak{}types.\allowbreak{}py:\allowbreak{}L1038-\allowbreak{}-\allowbreak{}L1046}; on re-run it returns the stored resume value instead \texttt{L1019-\allowbreak{}-\allowbreak{}L1036}; \texttt{GraphInterrupt} \texttt{LG/\allowbreak{}errors.\allowbreak{}py:\allowbreak{}L103-\allowbreak{}-\allowbreak{}L108}
\item \textbf{Pause bookkeeping:} the runner saves the interrupt as an \texttt{INTERRUPT} pending write \texttt{LG/\allowbreak{}pregel/\allowbreak{}\_\allowbreak{}runner.\allowbreak{}py:\allowbreak{}L586-\allowbreak{}-\allowbreak{}L592}; at the root graph the exception is swallowed, state is emitted and the run returns \texttt{LG/\allowbreak{}pregel/\allowbreak{}\_\allowbreak{}loop.\allowbreak{}py:\allowbreak{}L1491-\allowbreak{}-\allowbreak{}L1530}
\item \textbf{Resume (harness-consumed answer):} \texttt{Command(\allowbreak{}resume=.\allowbreak{}.\allowbreak{}.\allowbreak{})} maps to a \texttt{RESUME} write \texttt{LG/\allowbreak{}pregel/\allowbreak{}\_\allowbreak{}io.\allowbreak{}py:\allowbreak{}L56-\allowbreak{}-\allowbreak{}L78}; \texttt{\_\allowbreak{}first} requires a checkpointer and demands interrupt ids when several interrupts are pending \texttt{LG/\allowbreak{}pregel/\allowbreak{}\_\allowbreak{}loop.\allowbreak{}py:\allowbreak{}L992-\allowbreak{}-\allowbreak{}L1009}; the task scratchpad looks up its resume value \texttt{LG/\allowbreak{}pregel/\allowbreak{}\_\allowbreak{}algo.\allowbreak{}py:\allowbreak{}L1279-\allowbreak{}-\allowbreak{}L1344}; resume is inferred from \texttt{None} input or a \texttt{Command} \texttt{LG/\allowbreak{}pregel/\allowbreak{}\_\allowbreak{}loop.\allowbreak{}py:\allowbreak{}L915-\allowbreak{}-\allowbreak{}L928}
\item \textbf{Approval gate is user-written:} the framework ships only the \texttt{HumanInterrupt}/\texttt{HumanResponse} shapes (\texttt{accept}/\texttt{ignore}/\texttt{response}/\texttt{edit}); no code in the harness consumes them. \texttt{PB/\allowbreak{}interrupt.\allowbreak{}py:\allowbreak{}L51-\allowbreak{}-\allowbreak{}L84}, \texttt{L87-\allowbreak{}-\allowbreak{}L105}
\item \textbf{Model routing keyed to the work (user-written):} \texttt{model} may be a callable \texttt{(\allowbreak{}state, runtime) -\allowbreak{}> BaseChatModel} resolved on every \texttt{agent} step. \texttt{PB/\allowbreak{}chat\_\allowbreak{}agent\_\allowbreak{}executor.\allowbreak{}py:\allowbreak{}L563-\allowbreak{}-\allowbreak{}L564}, \texttt{L599-\allowbreak{}-\allowbreak{}L618}, \texttt{L674-\allowbreak{}-\allowbreak{}L677}
\item \textbf{Second model call for structured output:} after the loop ends, \texttt{generate\_\allowbreak{}structured\_\allowbreak{}response} runs as its own node. \texttt{PB/\allowbreak{}chat\_\allowbreak{}agent\_\allowbreak{}executor.\allowbreak{}py:\allowbreak{}L744-\allowbreak{}-\allowbreak{}L767}, wiring \texttt{L899-\allowbreak{}-\allowbreak{}L910}
\item \textbf{Node concurrency within a superstep:} all ready tasks are submitted to the executor and awaited with \texttt{FIRST\_\allowbreak{}COMPLETED}; the first non-interrupt failure cancels the rest. \texttt{LG/\allowbreak{}pregel/\allowbreak{}\_\allowbreak{}runner.\allowbreak{}py:\allowbreak{}L259-\allowbreak{}-\allowbreak{}L290}, \texttt{L618-\allowbreak{}-\allowbreak{}L636}, \texttt{L652-\allowbreak{}-\allowbreak{}L699}
\item \textbf{Cooperative drain (operator stop):} \texttt{RunControl.\allowbreak{}request\_\allowbreak{}drain(\allowbreak{})} is checked at the superstep boundary; the loop sets \texttt{status="draining"} and \texttt{stream(\allowbreak{})} raises \texttt{GraphDrained} after checkpointing. \texttt{LG/\allowbreak{}runtime.\allowbreak{}py:\allowbreak{}L79-\allowbreak{}-\allowbreak{}L104}; \texttt{LG/\allowbreak{}pregel/\allowbreak{}\_\allowbreak{}loop.\allowbreak{}py:\allowbreak{}L706-\allowbreak{}-\allowbreak{}L708}; \texttt{LG/\allowbreak{}pregel/\allowbreak{}main.\allowbreak{}py:\allowbreak{}L3103-\allowbreak{}-\allowbreak{}L3106}; \texttt{LG/\allowbreak{}errors.\allowbreak{}py:\allowbreak{}L55-\allowbreak{}-\allowbreak{}L65}
\item \textbf{Per-node error-handler nodes (user-written handler, framework routing):} \texttt{compile} assigns default/explicit handler nodes \texttt{LG/\allowbreak{}graph/\allowbreak{}state.\allowbreak{}py:\allowbreak{}L1290-\allowbreak{}-\allowbreak{}L1344}; the runner reroutes a failed task to its handler instead of panicking \texttt{LG/\allowbreak{}pregel/\allowbreak{}\_\allowbreak{}runner.\allowbreak{}py:\allowbreak{}L298-\allowbreak{}-\allowbreak{}L324}
\item \textbf{Sub-agents / subgraphs:} a compiled agent is used as a node of another graph via \texttt{name} \texttt{PB/\allowbreak{}chat\_\allowbreak{}agent\_\allowbreak{}executor.\allowbreak{}py:\allowbreak{}L466-\allowbreak{}-\allowbreak{}L468}; nesting is detected by \texttt{CONFIG\_\allowbreak{}KEY\_\allowbreak{}TASK\_\allowbreak{}ID} in config \texttt{LG/\allowbreak{}pregel/\allowbreak{}\_\allowbreak{}loop.\allowbreak{}py:\allowbreak{}L351}; a child's \texttt{Command(\allowbreak{}graph=PARENT, goto=[Send.\allowbreak{}.\allowbreak{}.\allowbreak{}])} is bubbled and merged by \texttt{ToolNode} so a tool can hand work to the parent \texttt{PB/\allowbreak{}tool\_\allowbreak{}node.\allowbreak{}py:\allowbreak{}L903-\allowbreak{}-\allowbreak{}L927}. Results return as channel writes to the parent task; depth is unbounded (namespace \texttt{parent|child:\allowbreak{}task\_\allowbreak{}id}).
\item \textbf{\texttt{tools\_\allowbreak{}condition} helper:} standalone routing function (\texttt{"tools"} or \texttt{"\_\allowbreak{}\_\allowbreak{}end\_\allowbreak{}\_\allowbreak{}"}) for user-built loops. \texttt{PB/\allowbreak{}tool\_\allowbreak{}node.\allowbreak{}py:\allowbreak{}L1616-\allowbreak{}-\allowbreak{}L1619}
\item \textbf{Hook nodes (user-written):} \texttt{pre\_\allowbreak{}model\_\allowbreak{}hook} runs before every \texttt{agent} step, \texttt{post\_\allowbreak{}model\_\allowbreak{}hook} after it; both are ordinary nodes the user supplies. \texttt{PB/\allowbreak{}chat\_\allowbreak{}agent\_\allowbreak{}executor.\allowbreak{}py:\allowbreak{}L876-\allowbreak{}-\allowbreak{}L881}, \texttt{L891-\allowbreak{}-\allowbreak{}L896}
\end{itemize}

\subsubsection*{Input}
\begin{itemize}
\item \textbf{Person's message enters as graph input:} \texttt{map\_\allowbreak{}input} turns the input dict into writes on the input channels. \texttt{LG/\allowbreak{}pregel/\allowbreak{}\_\allowbreak{}io.\allowbreak{}py:\allowbreak{}L81-\allowbreak{}-\allowbreak{}L97}; applied as a null-task write set on the first tick \texttt{LG/\allowbreak{}pregel/\allowbreak{}\_\allowbreak{}loop.\allowbreak{}py:\allowbreak{}L1086-\allowbreak{}-\allowbreak{}L1121}; the hidden \texttt{START} node forwards it into state \texttt{LG/\allowbreak{}graph/\allowbreak{}state.\allowbreak{}py:\allowbreak{}L1508-\allowbreak{}-\allowbreak{}L1514}
\item \textbf{Coercion of raw messages:} \texttt{add\_\allowbreak{}messages} converts dicts/tuples to \texttt{BaseMessage}s and assigns UUIDs to id-less messages. \texttt{LG/\allowbreak{}graph/\allowbreak{}message.\allowbreak{}py:\allowbreak{}L190-\allowbreak{}-\allowbreak{}L211}; ids are fixed before writes are stored \texttt{LG/\allowbreak{}pregel/\allowbreak{}\_\allowbreak{}loop.\allowbreak{}py:\allowbreak{}L1087-\allowbreak{}-\allowbreak{}L1089}, \texttt{L461-\allowbreak{}-\allowbreak{}L463}
\item \textbf{Tool results captured as \texttt{ToolMessage}:} raw tool returns are normalized (\texttt{Command}, \texttt{ToolMessage}, or list) \texttt{PB/\allowbreak{}tool\_\allowbreak{}node.\allowbreak{}py:\allowbreak{}L1466-\allowbreak{}-\allowbreak{}L1487}; non-string outputs are JSON-dumped \texttt{L317-\allowbreak{}-\allowbreak{}L344}; failures become \texttt{ToolMessage(\allowbreak{}status="error")} \texttt{L1022-\allowbreak{}-\allowbreak{}L1029}
\item \textbf{Human answer the model consumes:} when a tool calls \texttt{interrupt(\allowbreak{})}, the resume value is returned to the tool code on re-run (\texttt{L1023-\allowbreak{}-\allowbreak{}L1036} of \texttt{LG/\allowbreak{}types.\allowbreak{}py}) and becomes input only if the tool puts it in its \texttt{ToolMessage} (user-written). Shape: \texttt{PB/\allowbreak{}interrupt.\allowbreak{}py:\allowbreak{}L87-\allowbreak{}-\allowbreak{}L105}
\item \textbf{State edits from outside (\texttt{Command(\allowbreak{}update=.\allowbreak{}.\allowbreak{}.\allowbreak{})}):} mapped to null-task channel writes. \texttt{LG/\allowbreak{}pregel/\allowbreak{}\_\allowbreak{}io.\allowbreak{}py:\allowbreak{}L76-\allowbreak{}-\allowbreak{}L78}
\item \textbf{Run-scoped \texttt{context} (operator-supplied, not model-visible unless user code shows it):} built into \texttt{Runtime} and passed to nodes and tools. \texttt{LG/\allowbreak{}pregel/\allowbreak{}main.\allowbreak{}py:\allowbreak{}L2963-\allowbreak{}-\allowbreak{}L2973}; \texttt{ToolRuntime.\allowbreak{}context} \texttt{PB/\allowbreak{}tool\_\allowbreak{}node.\allowbreak{}py:\allowbreak{}L814-\allowbreak{}-\allowbreak{}L825}
\end{itemize}

\subsubsection*{Context}
\begin{itemize}
\item \textbf{Prompt assembly (user chooses form):} \texttt{\_\allowbreak{}get\_\allowbreak{}prompt\_\allowbreak{}runnable} prepends a \texttt{SystemMessage} (from \texttt{str}/\texttt{SystemMessage}) to \texttt{state["messages"]}, or runs a user callable/Runnable over the whole state. \texttt{PB/\allowbreak{}chat\_\allowbreak{}agent\_\allowbreak{}executor.\allowbreak{}py:\allowbreak{}L137-\allowbreak{}-\allowbreak{}L170}
\item \textbf{What the model is shown each call:} \texttt{state["messages"]}, or \texttt{llm\_\allowbreak{}input\_\allowbreak{}messages} produced by \texttt{pre\_\allowbreak{}model\_\allowbreak{}hook} (which does not alter stored history). \texttt{PB/\allowbreak{}chat\_\allowbreak{}agent\_\allowbreak{}executor.\allowbreak{}py:\allowbreak{}L636-\allowbreak{}-\allowbreak{}L658}; contract documented \texttt{L396-\allowbreak{}-\allowbreak{}L424}
\item \textbf{Tool descriptions:} \texttt{bind\_\allowbreak{}tools(\allowbreak{}tool\_\allowbreak{}classes + llm\_\allowbreak{}builtin\_\allowbreak{}tools)} unless the user already bound them; \texttt{\_\allowbreak{}should\_\allowbreak{}bind\_\allowbreak{}tools} cross-checks names (OpenAI- and Anthropic-style schemas). \texttt{PB/\allowbreak{}chat\_\allowbreak{}agent\_\allowbreak{}executor.\allowbreak{}py:\allowbreak{}L582-\allowbreak{}-\allowbreak{}L589}, \texttt{L173-\allowbreak{}-\allowbreak{}L217}
\item \textbf{Conversation as reducer-merged state:} \texttt{AgentState.\allowbreak{}messages} uses \texttt{add\_\allowbreak{}messages}, which upserts by id and applies \texttt{RemoveMessage}. \texttt{PB/\allowbreak{}chat\_\allowbreak{}agent\_\allowbreak{}executor.\allowbreak{}py:\allowbreak{}L57-\allowbreak{}-\allowbreak{}L62}; \texttt{LG/\allowbreak{}graph/\allowbreak{}message.\allowbreak{}py:\allowbreak{}L218-\allowbreak{}-\allowbreak{}L237}
\item \textbf{Window fitting (left to the user):} the only primitives are \texttt{RemoveMessage(\allowbreak{}id=REMOVE\_\allowbreak{}ALL\_\allowbreak{}MESSAGES)} to replace history \texttt{LG/\allowbreak{}graph/\allowbreak{}message.\allowbreak{}py:\allowbreak{}L212-\allowbreak{}-\allowbreak{}L216} and the \texttt{pre\_\allowbreak{}model\_\allowbreak{}hook} slot \texttt{PB/\allowbreak{}chat\_\allowbreak{}agent\_\allowbreak{}executor.\allowbreak{}py:\allowbreak{}L396-\allowbreak{}-\allowbreak{}L424}; trimming/summarization code: not found in this repository (grep for \texttt{trim\_\allowbreak{}messages}/\texttt{summariz} hits only docstrings).
\item \textbf{History validation before the call:} every \texttt{AIMessage.\allowbreak{}tool\_\allowbreak{}calls} id must have a \texttt{ToolMessage}, else \texttt{INVALID\_\allowbreak{}CHAT\_\allowbreak{}HISTORY}. \texttt{PB/\allowbreak{}chat\_\allowbreak{}agent\_\allowbreak{}executor.\allowbreak{}py:\allowbreak{}L243-\allowbreak{}-\allowbreak{}L271}
\item \textbf{Node input is read from channels:} \texttt{\_\allowbreak{}proc\_\allowbreak{}input} collects the node's subscribed channels plus managed values, then applies the schema mapper. \texttt{LG/\allowbreak{}pregel/\allowbreak{}\_\allowbreak{}algo.\allowbreak{}py:\allowbreak{}L1347-\allowbreak{}-\allowbreak{}L1391}; conditional edges read a fresh copy including the node's own writes \texttt{L189-\allowbreak{}-\allowbreak{}L225}; node reads/writes go through \texttt{CONFIG\_\allowbreak{}KEY\_\allowbreak{}READ}/\texttt{CONFIG\_\allowbreak{}KEY\_\allowbreak{}SEND} \texttt{LG/\allowbreak{}pregel/\allowbreak{}\_\allowbreak{}read.\allowbreak{}py:\allowbreak{}L74-\allowbreak{}-\allowbreak{}L91}, \texttt{LG/\allowbreak{}pregel/\allowbreak{}\_\allowbreak{}write.\allowbreak{}py:\allowbreak{}L106-\allowbreak{}-\allowbreak{}L126}
\item \textbf{Long-term memory for the model (cross-thread store):} \texttt{BaseStore.\allowbreak{}get/\allowbreak{}put/\allowbreak{}search} \texttt{CP/\allowbreak{}store/\allowbreak{}base/\allowbreak{}\_\allowbreak{}\_\allowbreak{}init\_\allowbreak{}\_\allowbreak{}.\allowbreak{}py:\allowbreak{}L708}, \texttt{L756}, \texttt{L784}, \texttt{L866}; injected into tools via \texttt{InjectedStore} \texttt{PB/\allowbreak{}tool\_\allowbreak{}node.\allowbreak{}py:\allowbreak{}L1441-\allowbreak{}-\allowbreak{}L1449}; retrieval and what to show the model are user-written.
\item \textbf{Hidden state injected into tools (not shown to the model):} \texttt{InjectedState}/\texttt{ToolRuntime} args are filled from graph state at call time and any model-supplied values for those names are stripped. \texttt{PB/\allowbreak{}tool\_\allowbreak{}node.\allowbreak{}py:\allowbreak{}L1400-\allowbreak{}-\allowbreak{}L1439}, \texttt{L1455-\allowbreak{}-\allowbreak{}L1463}; \texttt{ToolRuntime} fields \texttt{L1756-\allowbreak{}-\allowbreak{}L1764}
\item \textbf{State snapshot travels with each tool task:} \texttt{ToolCallWithContext} carries \texttt{state} inside the \texttt{Send}; \texttt{\_\allowbreak{}extract\_\allowbreak{}state} hydrates it (or reads channels). \texttt{PB/\allowbreak{}tool\_\allowbreak{}node.\allowbreak{}py:\allowbreak{}L306-\allowbreak{}-\allowbreak{}L314}, \texttt{L1315-\allowbreak{}-\allowbreak{}L1347}
\item \textbf{Structured-output prompt:} an optional system prompt is prepended for the final structured call. \texttt{PB/\allowbreak{}chat\_\allowbreak{}agent\_\allowbreak{}executor.\allowbreak{}py:\allowbreak{}L756-\allowbreak{}-\allowbreak{}L758}
\item \textbf{Agent name stamped on responses:} \texttt{response.\allowbreak{}name = name} (used in multi-agent transcripts). \texttt{PB/\allowbreak{}chat\_\allowbreak{}agent\_\allowbreak{}executor.\allowbreak{}py:\allowbreak{}L681-\allowbreak{}-\allowbreak{}L682}
\item \textbf{Injection of \texttt{config}/\texttt{store}/\texttt{writer}/\texttt{runtime} into node functions by signature:} \texttt{LG/\allowbreak{}\_\allowbreak{}internal/\allowbreak{}\_\allowbreak{}runnable.\allowbreak{}py:\allowbreak{}L168-\allowbreak{}-\allowbreak{}L219}, \texttt{L392-\allowbreak{}-\allowbreak{}L419}
\end{itemize}

\subsubsection*{Model interface}
\begin{itemize}
\item \textbf{The call:} \texttt{static\_\allowbreak{}model.\allowbreak{}invoke(\allowbreak{}model\_\allowbreak{}input, config)} (async \texttt{ainvoke}); the model is a \texttt{langchain\_\allowbreak{}core} \texttt{BaseChatModel}, outside this repo. \texttt{PB/\allowbreak{}chat\_\allowbreak{}agent\_\allowbreak{}executor.\allowbreak{}py:\allowbreak{}L672-\allowbreak{}-\allowbreak{}L679}, \texttt{L699-\allowbreak{}-\allowbreak{}L707}; string ids resolved with \texttt{init\_\allowbreak{}chat\_\allowbreak{}model} \texttt{L569-\allowbreak{}-\allowbreak{}L580}
\item \textbf{Structured output call:} \texttt{model.\allowbreak{}with\_\allowbreak{}structured\_\allowbreak{}output(\allowbreak{}schema).\allowbreak{}invoke(\allowbreak{}messages, config)}. \texttt{PB/\allowbreak{}chat\_\allowbreak{}agent\_\allowbreak{}executor.\allowbreak{}py:\allowbreak{}L761-\allowbreak{}-\allowbreak{}L766}
\item \textbf{Receiving a token stream:} \texttt{StreamMessagesHandler} is a LangChain callback: \texttt{on\_\allowbreak{}chat\_\allowbreak{}model\_\allowbreak{}start} records namespace/metadata, \texttt{on\_\allowbreak{}llm\_\allowbreak{}new\_\allowbreak{}token} forwards each \texttt{ChatGenerationChunk}, \texttt{on\_\allowbreak{}llm\_\allowbreak{}end} emits the final message once. \texttt{LG/\allowbreak{}pregel/\allowbreak{}\_\allowbreak{}messages.\allowbreak{}py:\allowbreak{}L139-\allowbreak{}-\allowbreak{}L158}, \texttt{L160-\allowbreak{}-\allowbreak{}L173}, \texttt{L175-\allowbreak{}-\allowbreak{}L188}, \texttt{L106-\allowbreak{}-\allowbreak{}L113}; installed when \texttt{"messages"} is in \texttt{stream\_\allowbreak{}mode} \texttt{LG/\allowbreak{}pregel/\allowbreak{}main.\allowbreak{}py:\allowbreak{}L2902-\allowbreak{}-\allowbreak{}L2918}; modes \texttt{LG/\allowbreak{}types.\allowbreak{}py:\allowbreak{}L131-\allowbreak{}-\allowbreak{}L133}
\item \textbf{Retries of a failed request (per node, user-configured):} \texttt{RetryPolicy(\allowbreak{}initial\_\allowbreak{}interval=0.\allowbreak{}5, backoff\_\allowbreak{}factor=2.\allowbreak{}0, max\_\allowbreak{}interval=128.\allowbreak{}0, max\_\allowbreak{}attempts=3, jitter=True, retry\_\allowbreak{}on=default\_\allowbreak{}retry\_\allowbreak{}on)} \texttt{LG/\allowbreak{}types.\allowbreak{}py:\allowbreak{}L427-\allowbreak{}-\allowbreak{}L446}; default predicate retries \texttt{ConnectionError} and HTTP 5xx, not \texttt{ValueError}/\texttt{TypeError}/\texttt{OSError} etc. \texttt{LG/\allowbreak{}\_\allowbreak{}internal/\allowbreak{}\_\allowbreak{}retry.\allowbreak{}py:\allowbreak{}L1-\allowbreak{}-\allowbreak{}L29}; loop with exponential backoff and jitter \texttt{LG/\allowbreak{}pregel/\allowbreak{}\_\allowbreak{}retry.\allowbreak{}py:\allowbreak{}L641-\allowbreak{}-\allowbreak{}L682}; policy comes from the node or the graph \texttt{L579}; \texttt{create\_\allowbreak{}react\_\allowbreak{}agent} sets none and \texttt{Pregel} defaults to \texttt{(\allowbreak{})} \texttt{LG/\allowbreak{}pregel/\allowbreak{}main.\allowbreak{}py:\allowbreak{}L786}, \texttt{L832-\allowbreak{}-\allowbreak{}L834}
\item \textbf{Backup model when the request cannot be served:} not found. The dynamic-model callable is keyed to state/context, not to failure (\texttt{PB/\allowbreak{}chat\_\allowbreak{}agent\_\allowbreak{}executor.\allowbreak{}py:\allowbreak{}L599-\allowbreak{}-\allowbreak{}L606}).
\item \textbf{Timeouts per attempt:} \texttt{TimeoutPolicy(\allowbreak{}run\_\allowbreak{}timeout, idle\_\allowbreak{}timeout, refresh\_\allowbreak{}on)} \texttt{LG/\allowbreak{}types.\allowbreak{}py:\allowbreak{}L461-\allowbreak{}-\allowbreak{}L523}; async watchdog tasks race the node \texttt{LG/\allowbreak{}pregel/\allowbreak{}\_\allowbreak{}retry.\allowbreak{}py:\allowbreak{}L422-\allowbreak{}-\allowbreak{}L470}; \texttt{NodeTimeoutError} is deliberately retryable \texttt{LG/\allowbreak{}errors.\allowbreak{}py:\allowbreak{}L191-\allowbreak{}-\allowbreak{}L207}; sync nodes cannot use timeouts \texttt{LG/\allowbreak{}pregel/\allowbreak{}\_\allowbreak{}retry.\allowbreak{}py:\allowbreak{}L580-\allowbreak{}-\allowbreak{}L583}
\item \textbf{Sync/async mismatch guard:} an async model callable invoked synchronously raises. \texttt{PB/\allowbreak{}chat\_\allowbreak{}agent\_\allowbreak{}executor.\allowbreak{}py:\allowbreak{}L664-\allowbreak{}-\allowbreak{}L670}
\end{itemize}

\subsubsection*{Output}
\begin{itemize}
\item \textbf{Tool executor (\texttt{ToolNode}):} parses the last \texttt{AIMessage.\allowbreak{}tool\_\allowbreak{}calls} (or a single \texttt{ToolCallWithContext}) \texttt{PB/\allowbreak{}tool\_\allowbreak{}node.\allowbreak{}py:\allowbreak{}L1258-\allowbreak{}-\allowbreak{}L1300}; runs all calls concurrently with a thread pool \texttt{executor.\allowbreak{}map} \texttt{L829-\allowbreak{}-\allowbreak{}L832} or \texttt{asyncio.\allowbreak{}gather} \texttt{L863-\allowbreak{}-\allowbreak{}L866}; results combined into \texttt{\{"messages":\allowbreak{} [.\allowbreak{}.\allowbreak{}.\allowbreak{}]\}} \texttt{L870-\allowbreak{}-\allowbreak{}L895}
\item \textbf{One call's execution:} \texttt{tool.\allowbreak{}invoke(\allowbreak{}call\_\allowbreak{}args, config)} after injection \texttt{PB/\allowbreak{}tool\_\allowbreak{}node.\allowbreak{}py:\allowbreak{}L960-\allowbreak{}-\allowbreak{}L966}; unknown tool name $\rightarrow$ error \texttt{ToolMessage} listing valid tools \texttt{L1302-\allowbreak{}-\allowbreak{}L1313}, template \texttt{L116-\allowbreak{}-\allowbreak{}L118}
\item \textbf{Error-to-message policy (user-configurable):} default handler only converts argument-validation errors (\texttt{ToolInvocationError}) and re-raises everything else \texttt{PB/\allowbreak{}tool\_\allowbreak{}node.\allowbreak{}py:\allowbreak{}L391-\allowbreak{}-\allowbreak{}L399}; \texttt{handle\_\allowbreak{}tool\_\allowbreak{}errors} may be bool/str/callable/exception types \texttt{L1002-\allowbreak{}-\allowbreak{}L1029}; templates \texttt{L119-\allowbreak{}-\allowbreak{}L129}
\item \textbf{Interceptor around each call (user-written):} \texttt{wrap\_\allowbreak{}tool\_\allowbreak{}call(\allowbreak{}request, execute)} can retry, cache, rewrite or short-circuit a call. \texttt{PB/\allowbreak{}tool\_\allowbreak{}node.\allowbreak{}py:\allowbreak{}L1061-\allowbreak{}-\allowbreak{}L1088}
\item \textbf{Tool returning \texttt{Command}:} validated and passed through so the tool's EFFECT (state update / goto) reaches the scheduler. \texttt{PB/\allowbreak{}tool\_\allowbreak{}node.\allowbreak{}py:\allowbreak{}L1473-\allowbreak{}-\allowbreak{}L1474}, \texttt{L897-\allowbreak{}-\allowbreak{}L928}
\item \textbf{Interrupt raised inside a tool is never caught:} \texttt{GraphBubbleUp} passes through the error handling. \texttt{PB/\allowbreak{}tool\_\allowbreak{}node.\allowbreak{}py:\allowbreak{}L986-\allowbreak{}-\allowbreak{}L996}
\item \textbf{Rendering to the person (stream modes):} \texttt{values} after each step \texttt{LG/\allowbreak{}pregel/\allowbreak{}\_\allowbreak{}loop.\allowbreak{}py:\allowbreak{}L751-\allowbreak{}-\allowbreak{}L759}; \texttt{updates}/\texttt{tasks} per task \texttt{L1571-\allowbreak{}-\allowbreak{}L1621}; \texttt{\_\allowbreak{}emit} filters by subscribed modes \texttt{L1535-\allowbreak{}-\allowbreak{}L1569}; final value of \texttt{invoke} is the last \texttt{values} chunk plus collected interrupts \texttt{LG/\allowbreak{}pregel/\allowbreak{}main.\allowbreak{}py:\allowbreak{}L4000-\allowbreak{}-\allowbreak{}L4043}
\item \textbf{Custom and tool-progress streams:} \texttt{stream\_\allowbreak{}writer} for \texttt{"custom"} \texttt{LG/\allowbreak{}pregel/\allowbreak{}main.\allowbreak{}py:\allowbreak{}L2932-\allowbreak{}-\allowbreak{}L2951}; \texttt{"tools"} mode installs \texttt{StreamToolCallHandler} (\texttt{on\_\allowbreak{}tool\_\allowbreak{}start}/\texttt{on\_\allowbreak{}tool\_\allowbreak{}end}) \texttt{L2921-\allowbreak{}-\allowbreak{}L2929}, \texttt{LG/\allowbreak{}pregel/\allowbreak{}\_\allowbreak{}tools.\allowbreak{}py:\allowbreak{}L35}, \texttt{L228-\allowbreak{}-\allowbreak{}L258}; tools push partial output with \texttt{ToolRuntime.\allowbreak{}emit\_\allowbreak{}output\_\allowbreak{}delta} \texttt{PB/\allowbreak{}tool\_\allowbreak{}node.\allowbreak{}py:\allowbreak{}L1766-\allowbreak{}-\allowbreak{}L1784}
\item \textbf{Sandboxing of tool execution:} not found; tools run in-process in the caller's executor (\texttt{PB/\allowbreak{}tool\_\allowbreak{}node.\allowbreak{}py:\allowbreak{}L829-\allowbreak{}-\allowbreak{}L832}).
\item \textbf{Argument validation without execution:} \texttt{ValidationNode} checks tool calls against schemas and returns error \texttt{ToolMessage}s for the model to fix. \texttt{PB/\allowbreak{}tool\_\allowbreak{}validator.\allowbreak{}py:\allowbreak{}L47}
\end{itemize}

\subsubsection*{Production concerns}
\begin{longtable}{@{}p{0.13\linewidth}p{0.17\linewidth}p{0.36\linewidth}p{0.26\linewidth}@{}}
\toprule \textbf{Concern} & \textbf{Where it lives (primitive)} & \textbf{Mechanism} & \textbf{Evidence} \\ \midrule \endhead
Persistence & Control flow (loop state) and Context (conversation state) & \texttt{Checkpoint} = channel values + versions + versions\_seen per thread; \texttt{BaseCheckpointSaver.\allowbreak{}put} after every superstep and \texttt{put\_\allowbreak{}writes} per task; \texttt{durability} \texttt{sync}/\texttt{async}/\texttt{exit} (default \texttt{async}); \texttt{thread\_\allowbreak{}id} required; backends \texttt{InMemorySaver}, \texttt{SqliteSaver}, \texttt{PostgresSaver}; resume loads latest checkpoint or a given \texttt{checkpoint\_\allowbreak{}id} (time travel); cross-thread \texttt{BaseStore} & \texttt{CP/\allowbreak{}checkpoint/\allowbreak{}base/\allowbreak{}\_\allowbreak{}\_\allowbreak{}init\_\allowbreak{}\_\allowbreak{}.\allowbreak{}py:\allowbreak{}L93-\allowbreak{}-\allowbreak{}L124}, \texttt{L301-\allowbreak{}-\allowbreak{}L342}; \texttt{LG/\allowbreak{}pregel/\allowbreak{}\_\allowbreak{}loop.\allowbreak{}py:\allowbreak{}L785-\allowbreak{}-\allowbreak{}L786}, \texttt{L498-\allowbreak{}-\allowbreak{}L528}, \texttt{L1261-\allowbreak{}-\allowbreak{}L1263}, \texttt{L1479-\allowbreak{}-\allowbreak{}L1489}; \texttt{LG/\allowbreak{}types.\allowbreak{}py:\allowbreak{}L98-\allowbreak{}-\allowbreak{}L104}; \texttt{LG/\allowbreak{}pregel/\allowbreak{}main.\allowbreak{}py:\allowbreak{}L2682-\allowbreak{}-\allowbreak{}L2686}, \texttt{L2695-\allowbreak{}-\allowbreak{}L2696}, \texttt{L3078-\allowbreak{}-\allowbreak{}L3079}; \texttt{LG/\allowbreak{}pregel/\allowbreak{}\_\allowbreak{}loop.\allowbreak{}py:\allowbreak{}L1783-\allowbreak{}-\allowbreak{}L1813}; \texttt{CP/\allowbreak{}checkpoint/\allowbreak{}memory/\allowbreak{}\_\allowbreak{}\_\allowbreak{}init\_\allowbreak{}\_\allowbreak{}.\allowbreak{}py:\allowbreak{}L34}; \texttt{libs/\allowbreak{}checkpoint-\allowbreak{}sqlite/\allowbreak{}langgraph/\allowbreak{}checkpoint/\allowbreak{}sqlite/\allowbreak{}\_\allowbreak{}\_\allowbreak{}init\_\allowbreak{}\_\allowbreak{}.\allowbreak{}py:\allowbreak{}L49}; \texttt{libs/\allowbreak{}checkpoint-\allowbreak{}postgres/\allowbreak{}langgraph/\allowbreak{}checkpoint/\allowbreak{}postgres/\allowbreak{}\_\allowbreak{}\_\allowbreak{}init\_\allowbreak{}\_\allowbreak{}.\allowbreak{}py:\allowbreak{}L45} \\
Safety \& security & Output (tool args), Control flow (pauses), Persistence layer & No sandbox (not found). Approval = user-written \texttt{interrupt(\allowbreak{})} before/inside a tool. Model cannot forge \texttt{InjectedState}/\texttt{InjectedStore} args (stripped). Resume values failing an interrupt's \texttt{response\_\allowbreak{}schema} are never fed back to the model. Checkpoint deserializer warns it must not see untrusted data; \texttt{LANGGRAPH\_\allowbreak{}STRICT\_\allowbreak{}MSGPACK} / msgpack allowlist; \texttt{EncryptedSerializer} wraps payloads; \texttt{TAG\_\allowbreak{}HIDDEN} keeps nodes out of streams/breakpoints & \texttt{PB/\allowbreak{}tool\_\allowbreak{}node.\allowbreak{}py:\allowbreak{}L1455-\allowbreak{}-\allowbreak{}L1463}, \texttt{L997-\allowbreak{}-\allowbreak{}L1001}; \texttt{CP/\allowbreak{}checkpoint/\allowbreak{}serde/\allowbreak{}jsonplus.\allowbreak{}py:\allowbreak{}L104-\allowbreak{}-\allowbreak{}L135}; \texttt{CP/\allowbreak{}checkpoint/\allowbreak{}base/\allowbreak{}\_\allowbreak{}\_\allowbreak{}init\_\allowbreak{}\_\allowbreak{}.\allowbreak{}py:\allowbreak{}L742-\allowbreak{}-\allowbreak{}L771}; \texttt{CP/\allowbreak{}checkpoint/\allowbreak{}serde/\allowbreak{}encrypted.\allowbreak{}py:\allowbreak{}L8-\allowbreak{}-\allowbreak{}L35}; \texttt{LG/\allowbreak{}pregel/\allowbreak{}\_\allowbreak{}algo.\allowbreak{}py:\allowbreak{}L175-\allowbreak{}-\allowbreak{}L182} \\
Observability & Across (traces via LangChain callbacks) & \texttt{on\_\allowbreak{}chain\_\allowbreak{}start/\allowbreak{}end/\allowbreak{}error} per run with \texttt{ls\_\allowbreak{}integration:\allowbreak{} langgraph} metadata; per-task metadata \texttt{langgraph\_\allowbreak{}step/\allowbreak{}node/\allowbreak{}triggers/\allowbreak{}path/\allowbreak{}checkpoint\_\allowbreak{}ns}; \texttt{checkpoints}/\texttt{tasks}/\texttt{debug} stream modes; retry attempts logged; \texttt{get\_\allowbreak{}state}/\texttt{get\_\allowbreak{}state\_\allowbreak{}history}; \texttt{ExecutionInfo} (thread/run/task ids, attempt). Token accounting: not found in LangGraph (left to \texttt{langchain\_\allowbreak{}core}/LangSmith). & \texttt{LG/\allowbreak{}pregel/\allowbreak{}main.\allowbreak{}py:\allowbreak{}L2865-\allowbreak{}-\allowbreak{}L2890}, \texttt{L3108-\allowbreak{}-\allowbreak{}L3112}; \texttt{LG/\allowbreak{}pregel/\allowbreak{}\_\allowbreak{}algo.\allowbreak{}py:\allowbreak{}L653-\allowbreak{}-\allowbreak{}L659}; \texttt{LG/\allowbreak{}pregel/\allowbreak{}\_\allowbreak{}loop.\allowbreak{}py:\allowbreak{}L665-\allowbreak{}-\allowbreak{}L684}, \texttt{L711}; \texttt{LG/\allowbreak{}pregel/\allowbreak{}\_\allowbreak{}retry.\allowbreak{}py:\allowbreak{}L676-\allowbreak{}-\allowbreak{}L680}; \texttt{LG/\allowbreak{}pregel/\allowbreak{}main.\allowbreak{}py:\allowbreak{}L1464}, \texttt{L1514}; \texttt{LG/\allowbreak{}runtime.\allowbreak{}py:\allowbreak{}L27-\allowbreak{}-\allowbreak{}L52} \\
Resilience & Model interface (retry/timeout), Control flow (resume, error handlers), Output (tool errors) & \texttt{RetryPolicy} with backoff; \texttt{TimeoutPolicy} with idle/run watchdogs; \texttt{step\_\allowbreak{}timeout} on a whole superstep; on resume, finished tasks' pending writes are re-applied so only unfinished tasks re-run; failed tasks routed to error-handler nodes; first failure cancels sibling tasks; \texttt{GraphDrained} on SIGTERM-style drain; tool errors returned to the model as messages. Compaction on overflow: not in framework (user \texttt{pre\_\allowbreak{}model\_\allowbreak{}hook}). & \texttt{LG/\allowbreak{}pregel/\allowbreak{}\_\allowbreak{}retry.\allowbreak{}py:\allowbreak{}L641-\allowbreak{}-\allowbreak{}L682}, \texttt{L422-\allowbreak{}-\allowbreak{}L470}; \texttt{LG/\allowbreak{}pregel/\allowbreak{}main.\allowbreak{}py:\allowbreak{}L781}, \texttt{L3060}; \texttt{LG/\allowbreak{}pregel/\allowbreak{}\_\allowbreak{}loop.\allowbreak{}py:\allowbreak{}L805-\allowbreak{}-\allowbreak{}L821}, \texttt{L691-\allowbreak{}-\allowbreak{}L694}; \texttt{LG/\allowbreak{}pregel/\allowbreak{}\_\allowbreak{}runner.\allowbreak{}py:\allowbreak{}L298-\allowbreak{}-\allowbreak{}L324}, \texttt{L652-\allowbreak{}-\allowbreak{}L699}; \texttt{LG/\allowbreak{}pregel/\allowbreak{}\_\allowbreak{}loop.\allowbreak{}py:\allowbreak{}L706-\allowbreak{}-\allowbreak{}L708}; \texttt{PB/\allowbreak{}tool\_\allowbreak{}node.\allowbreak{}py:\allowbreak{}L1022-\allowbreak{}-\allowbreak{}L1029} \\
Performance & Output (parallel tools), Control flow (parallel nodes, caching), Model interface (streaming) & Tool calls run concurrently (thread pool / gather); nodes in a superstep run concurrently; node-level \texttt{CachePolicy(\allowbreak{}key\_\allowbreak{}func, ttl)} with \texttt{BaseCache} lookup before execution; per-step input cache; \texttt{trigger\_\allowbreak{}to\_\allowbreak{}nodes} avoids scanning all nodes; token streaming; \texttt{durability="async"/\allowbreak{}"exit"} defers checkpoint writes & \texttt{PB/\allowbreak{}tool\_\allowbreak{}node.\allowbreak{}py:\allowbreak{}L829-\allowbreak{}-\allowbreak{}L832}, \texttt{L863-\allowbreak{}-\allowbreak{}L866}; \texttt{LG/\allowbreak{}pregel/\allowbreak{}\_\allowbreak{}runner.\allowbreak{}py:\allowbreak{}L259-\allowbreak{}-\allowbreak{}L277}; \texttt{LG/\allowbreak{}types.\allowbreak{}py:\allowbreak{}L529-\allowbreak{}-\allowbreak{}L538}; \texttt{LG/\allowbreak{}pregel/\allowbreak{}\_\allowbreak{}loop.\allowbreak{}py:\allowbreak{}L1701-\allowbreak{}-\allowbreak{}L1714}; \texttt{LG/\allowbreak{}pregel/\allowbreak{}main.\allowbreak{}py:\allowbreak{}L3056-\allowbreak{}-\allowbreak{}L3057}; \texttt{LG/\allowbreak{}pregel/\allowbreak{}\_\allowbreak{}algo.\allowbreak{}py:\allowbreak{}L667-\allowbreak{}-\allowbreak{}L686}, \texttt{L1357-\allowbreak{}-\allowbreak{}L1359}, \texttt{L474-\allowbreak{}-\allowbreak{}L481}; \texttt{CP/\allowbreak{}cache/\allowbreak{}base/\allowbreak{}\_\allowbreak{}\_\allowbreak{}init\_\allowbreak{}\_\allowbreak{}.\allowbreak{}py:\allowbreak{}L15} \\
Evaluation & Across & Unit/integration test suites per package (prebuilt: \texttt{test\_\allowbreak{}react\_\allowbreak{}agent.\allowbreak{}py} 46 tests, \texttt{test\_\allowbreak{}tool\_\allowbreak{}node.\allowbreak{}py} 42; runtime: \texttt{test\_\allowbreak{}pregel.\allowbreak{}py} 172, \texttt{test\_\allowbreak{}retry.\allowbreak{}py} 105, \texttt{test\_\allowbreak{}time\_\allowbreak{}travel.\allowbreak{}py} 50, \texttt{test\_\allowbreak{}interruption.\allowbreak{}py} 6); a \texttt{checkpoint-\allowbreak{}conformance} package for savers. Agent-quality evals: not found in this repository. & \texttt{libs/\allowbreak{}prebuilt/\allowbreak{}tests/\allowbreak{}test\_\allowbreak{}react\_\allowbreak{}agent.\allowbreak{}py}; \texttt{libs/\allowbreak{}prebuilt/\allowbreak{}tests/\allowbreak{}test\_\allowbreak{}tool\_\allowbreak{}node.\allowbreak{}py}; \texttt{libs/\allowbreak{}langgraph/\allowbreak{}tests/\allowbreak{}test\_\allowbreak{}pregel.\allowbreak{}py}; \texttt{libs/\allowbreak{}langgraph/\allowbreak{}tests/\allowbreak{}test\_\allowbreak{}retry.\allowbreak{}py}; \texttt{libs/\allowbreak{}checkpoint-\allowbreak{}conformance/\allowbreak{}tests/\allowbreak{}test\_\allowbreak{}validate\_\allowbreak{}memory.\allowbreak{}py} \\
\bottomrule \end{longtable}

\subsubsection*{Design decisions (one line each: the choice this harness made within a primitive)}
\begin{itemize}
\item Input: everything (user message, tool result, resume value, external \texttt{Command(\allowbreak{}update)}) enters as a write to a named channel of a typed state, merged by a user-chosen reducer (\texttt{add\_\allowbreak{}messages} upserts by message id).
\item Context: the model sees exactly \texttt{state["messages"]} (optionally prefixed by one system message); fitting to the window is delegated entirely to a user-written \texttt{pre\_\allowbreak{}model\_\allowbreak{}hook}, with \texttt{REMOVE\_\allowbreak{}ALL\_\allowbreak{}MESSAGES} as the only built-in rewrite primitive.
\item Model interface: no client of its own; a LangChain \texttt{BaseChatModel} is invoked once per \texttt{agent} node, with retries/timeouts expressed as per-node policies that are off by default.
\item Output: \texttt{ToolNode} runs every tool call of a response concurrently in-process, converts results and (configurable) errors into \texttt{ToolMessage}s, and lets a tool return a \texttt{Command} to reach the scheduler.
\item Control flow: the loop is a user-written graph executed as Pregel supersteps; "go again" is a conditional edge, "stop" is an empty task set or \texttt{recursion\_\allowbreak{}limit}, pauses are exceptions persisted as checkpoint writes and resumed by re-running the node.
\end{itemize}

\subsubsection*{Hard cases}
\begin{itemize}
\item \textbf{Conditional edge \texttt{should\_\allowbreak{}continue}:} user code, but it decides go-again vs. stop and is executed by the scheduler as a writer on the node; placed in control flow by function (rule: assign by function). Note the whole graph is user-authored control flow; the prebuilt merely authors it for you.
\item \textbf{\texttt{remaining\_\allowbreak{}steps} managed value:} it is injected into the node's \emph{state} (so it looks like context), but no model ever sees it; \texttt{call\_\allowbreak{}model} uses it to replace the model's answer with a canned message. Placed in control flow (step budget). Its side effect --- a synthetic AIMessage written into the conversation --- is a context effect of a control-flow decision.
\item \textbf{\texttt{Send(\allowbreak{}"tools", ToolCallWithContext(\allowbreak{}state=.\allowbreak{}.\allowbreak{}.\allowbreak{}))}:} a scheduling packet (control flow) that also carries a state snapshot (context) to the tool task. Split: the \texttt{Send} is control flow; the snapshot and \texttt{\_\allowbreak{}extract\_\allowbreak{}state} are context.
\item \textbf{Tool returning \texttt{Command(\allowbreak{}goto=.\allowbreak{}.\allowbreak{}.\allowbreak{})} or \texttt{Command(\allowbreak{}graph=PARENT)}:} the tool call is output; its effect is routing, so \texttt{\_\allowbreak{}combine\_\allowbreak{}tool\_\allowbreak{}outputs} and \texttt{\_\allowbreak{}control\_\allowbreak{}branch} are control flow (rule: a tool's EFFECT is assigned to what it changes).
\item \textbf{\texttt{interrupt(\allowbreak{})} inside a tool and \texttt{HumanResponse}:} the pause/resume machinery is control flow, but the harness never consumes an approval; the framework hands the raw resume value to user code in the tool, which decides whether to execute. So the approval gate is user-written, and the value is input only if the tool shows it to the model.
\item \textbf{Dynamic model callable:} a model change keyed to state/context (planning vs execution) is control flow by the rules, but the selection logic is user-written; the framework only re-resolves it before each call. No failure-keyed fallback exists, so "backup model" is absent from the model-interface primitive.
\item \textbf{Checkpointer as conversation memory:} the same mechanism that makes the loop resumable (persistence of control-flow state) is what makes a \texttt{thread\_\allowbreak{}id} a chat history (persistence of context). Recorded once under persistence, spanning both primitives.
\item \textbf{Pregel channel/version scheduler:} it is a general dataflow scheduler (nodes fire when subscribed channels change), not a loop with explicit go-again/hand-back/stop. It fits control flow, but the "shape of the loop" is emergent from edges and versions rather than coded; the only explicit stop conditions are "no tasks" and the step budget.
\item \textbf{\texttt{DEFAULT\_\allowbreak{}RECURSION\_\allowbreak{}LIMIT = "10007"}:} the value found at this commit (\texttt{LG/\allowbreak{}\_\allowbreak{}internal/\allowbreak{}\_\allowbreak{}config.\allowbreak{}py:\allowbreak{}L32}) differs from the widely documented default of 25; the clone has a single squashed commit, so the change cannot be traced. Recorded as found, not interpreted.
\item \textbf{Run \texttt{context} / \texttt{Runtime}:} operator-supplied static data that is neither a person's message nor a tool result; it is closest to configuration (out of scope) but is delivered on the request path to tools, so it is noted under Input without a strong claim.
\item \textbf{Error-handler nodes and \texttt{wrap\_\allowbreak{}tool\_\allowbreak{}call}:} both are user-written code that the framework schedules; they are resilience mechanisms whose primitive depends on what the user's code does (a retry inside \texttt{wrap\_\allowbreak{}tool\_\allowbreak{}call} is output-side; an error handler that routes is control flow).
\end{itemize}

\subsubsection*{Summary row}
\begin{longtable}{@{}p{0.18\linewidth}p{0.18\linewidth}p{0.18\linewidth}p{0.18\linewidth}p{0.18\linewidth}@{}}
\toprule \textbf{Input} & \textbf{Context} & \textbf{Model interface} & \textbf{Output} & \textbf{Control flow} \\ \midrule \endhead
Messages and tool results become reducer-merged channel writes & System prompt plus state messages; fitting left to user hook & One LangChain chat-model invoke per step; per-node retries optional & ToolNode runs tool calls concurrently, errors become ToolMessages & User-written graph run as Pregel supersteps; interrupts checkpointed \\
\bottomrule \end{longtable}


\normalsize

\endgroup
\end{document}
